% This LaTeX document needs to be compiled with XeLaTeX.
\documentclass[10pt]{article}
\usepackage[utf8]{inputenc}
\usepackage{amsmath}
\usepackage{amsfonts}
\usepackage{amssymb}
\usepackage[version=4]{mhchem}
\usepackage{stmaryrd}
\usepackage{hyperref}
\hypersetup{colorlinks=true, linkcolor=blue, filecolor=magenta, urlcolor=cyan,}
\urlstyle{same}
\usepackage{graphicx}
\usepackage[export]{adjustbox}
\graphicspath{ {./images/} }
\usepackage{caption}
\usepackage{bbold}
\usepackage{multirow}
\usepackage{fvextra, csquotes}
% \usepackage[fallback]{xeCJK}
\usepackage{polyglossia}
\usepackage{fontspec}
\usepackage{newunicodechar}
\IfFontExistsTF{Noto Serif CJK TC}
{\setCJKmainfont{Noto Serif CJK TC}}
{\IfFontExistsTF{STSong}
  {\setCJKmainfont{STSong}}
  {\IfFontExistsTF{Droid Sans Fallback}
    {\setCJKmainfont{Droid Sans Fallback}}
    {\setCJKmainfont{SimSun}}
}}

\setmainlanguage{english}
\IfFontExistsTF{CMU Serif}
{\setmainfont{CMU Serif}}
{\IfFontExistsTF{DejaVu Sans}
  {\setmainfont{DejaVu Sans}}
  {\setmainfont{Georgia}}
}

 Finite-Sample Properties of OLS }
%New command to display footnote whose markers will always be hidden
\let\svthefootnote\thefootnote
\newcommand\blfootnotetext[1]{%
  \let\thefootnote\relax\footnote{#1}%
  \addtocounter{footnote}{-1}%
  \let\thefootnote\svthefootnote%
}
%Overriding the \footnotetext command to hide the marker if its value is `0`
\let\svfootnotetext\footnotetext
\renewcommand\footnotetext[2][?]{%
  \if\relax#1\relax%
    \ifnum\value{footnote}=0\blfootnotetext{#2}\else\svfootnotetext{#2}\fi%
  \else%
    \if?#1\ifnum\value{footnote}=0\blfootnotetext{#2}\else\svfootnotetext{#2}\fi%
    \else\svfootnotetext[#1]{#2}\fi%
  \fi
}
\newunicodechar{×}{\ifmmode\times\else{$\times$}\fi}
\newunicodechar{⋮}{\ifmmode\vdots\else{$\vdots$}\fi}
\newunicodechar{⇔}{\ifmmode\Leftrightarrow\else{$\Leftrightarrow$}\fi}
\newunicodechar{¥}{\ifmmode\text{\yen}\else\yen\fi}
\title{Chapter 3: Single-Equation GMM}

\author{}
\date{}

\begin{document}
\maketitle

\section*{CHAPTER 3}
\section*{Single-Equation GMM}
\begin{abstract}
The most important assumption made for the OLS is the orthogonality between the error term and regressors. Without it, the OLS estimator is not even consistent. Since in many important applications the orthogonality condition is not satisfied, it is imperative to be able to deal with endogenous regressors. The estimation method called the Generalized Method of Moments (GMM), which includes OLS as a special case, provides a solution. This chapter presents GMM for single-equation models, while the next chapter develops its multiple-equation extension. These chapters deal only with linear equations. GMM estimation of nonlinear equations will be covered in Chapter 7. The major alternative to GMM, the maximum likelihood (ML) estimation, will be covered in Chapter 8.\\
Reflecting the prevalence of endogenous regressors in economics, this chapter starts out with a number of examples. This is followed by a general formulation of endogenous regressors in Section 3.3. Section 3.4 introduces the GMM procedure for the model of Section 3.3. Sections 3.5-3.7 are devoted to developing the large sample properties of the GMM estimator and associated test statistics. Under conditional homoskedasticity, the formulas derived in these sections can be simplified. Section 3.8 collects those simplified formulas. In particular, the two-stage least squares (2SLS), the techniques originally designed for the estimation of simultaneous equations models, is a special case of GMM. The ML counterpart of 2SLS is called limited-information maximum likelihood (LIML), which will be covered in Section 8.6.\\
The empirical exercise of the chapter takes up the most widely estimated equation in economics, the wage equation. The equation relates the wage rate to the individual's various characteristics such as education and ability. Because education is a choice made by individuals, and also because ability is imperfectly measured in data, the regressors are likely to be endogenous. We apply the estimation techniques introduced in this chapter to the wage equation and verify that the parameter estimates depend on whether we correct for endogeneity.
\end{abstract}

\subsection*{3.1 Endogeneity Bias: Working's Example}
\section*{A Simultaneous Equations Model of Market Equilibrium}
The classic illustration of biases created by endogenous regressors is Working (1927). Consider the following simple model of demand and supply:

\[
\begin{array}{cl}
q_{i}^{d}=\alpha_{0}+\alpha_{1} p_{i}+u_{i}, & \text { (demand equation) } \\
q_{i}^{s}=\beta_{0}+\beta_{1} p_{i}+v_{i}, & \text { (supply equation) } \\
q_{i}^{d}=q_{i}^{s}, & \text { (market equilibrium) } \tag{3.1.1c}
\end{array}
\]

where $q_{i}^{d}$ is the quantity demanded for the commodity in question (coffee, say) in period $i, q_{i}^{s}$ is the quantity supplied, and $p_{i}$ is the price. The error term $u_{i}$ in the demand equation represents factors that influence coffee demand other than price, such as the public's mood for coffee. Depending on the value of $u_{i}$, the demand curve in the price-quantity plane shifts up or down. Similarly, $v_{i}$ represents supply factors other than price. We assume that $\mathrm{E}\left(u_{i}\right)=0$ and $\mathrm{E}\left(v_{i}\right)=0$ (if not, include the nonzero means in the intercepts $\alpha_{0}$ and $\beta_{0}$ ). To avoid inessential complications, we also assume $\operatorname{Cov}\left(u_{i}, v_{i}\right)=0$. If we define $q_{i}=q_{i}^{d}=q_{i}^{s}$, the three-equation system (3.1.1) can be reduced to a two-equation system:

\[
\begin{array}{ll}
q_{i}=\alpha_{0}+\alpha_{1} p_{i}+u_{i}, & \text { (demand equation) } \\
q_{i}=\beta_{0}+\beta_{1} p_{i}+v_{i} . & \text { (supply equation) } \tag{3.1.2b}
\end{array}
\]

We say that a regressor is endogenous if it is not predetermined (i.e., not orthogonal to the error term), that is, if it does not satisfy the orthogonality condition. When the equation includes the intercept, the orthogonality condition is violated and hence the regressor is endogenous, if and only if the regressor is correlated with the error term. In the present example, the regressor $p_{i}$ is necessarily endogenous in both equations. To see why, treat (3.1.2) as a system of two simultaneous equations and solve for ( $p_{i}, q_{i}$ ) as


\begin{gather*}
p_{i}=\frac{\beta_{0}-\alpha_{0}}{\alpha_{1}-\beta_{1}}+\frac{v_{i}-u_{i}}{\alpha_{1}-\beta_{1}}  \tag{3.1.3a}\\
q_{i}=\frac{\alpha_{1} \beta_{0}-\alpha_{0} \beta_{1}}{\alpha_{1}-\beta_{1}}+\frac{\alpha_{1} v_{i}-\beta_{1} u_{i}}{\alpha_{1}-\beta_{1}} . \tag{3.1.3b}
\end{gather*}


So price is a function of the two error terms. From (3.1.3a), we can calculate the covariance of the regressor $p_{i}$ with the demand shifter $u_{i}$ and the supply shifter $v_{i}$ to be


\begin{equation*}
\operatorname{Cov}\left(p_{i}, u_{i}\right)=-\frac{\operatorname{Var}\left(u_{i}\right)}{\alpha_{1}-\beta_{1}}, \quad \operatorname{Cov}\left(p_{i}, v_{i}\right)=\frac{\operatorname{Var}\left(v_{i}\right)}{\alpha_{1}-\beta_{1}}, \tag{3.1.4}
\end{equation*}


which are not zero (unless $\operatorname{Var}\left(u_{i}\right)=0$ and $\operatorname{Var}\left(v_{i}\right)=0$ ). Therefore, price is correlated positively with the demand shifter and negatively with the supply shifter, if the demand curve is downward-sloping $\left(\alpha_{1}<0\right)$ and the supply curve upwardsloping $\left(\beta_{1}>0\right)$. In this example, endogeneity is a result of market equilibrium.

\section*{Endogeneity Bias}
When quantity is regressed on a constant and price, does it estimate the demand curve or the supply curve? The answer is neither, because price is endogenous in both the demand and supply equations. Recall from Section 2.9 that the OLS estimator is consistent for the least squares projection coefficients. In the least squares projection of $q_{i}$ on a constant and $p_{i}$, the coefficient of $p_{i}$ is $\operatorname{Cov}\left(p_{i}, q_{i}\right) / \operatorname{Var}\left(p_{i}\right),{ }^{1}$ so,


\begin{equation*}
\text { plim of the OLS estimate of the price coefficient }=\frac{\operatorname{Cov}\left(p_{i}, q_{i}\right)}{\operatorname{Var}\left(p_{i}\right)} . \tag{3.1.5}
\end{equation*}


To rewrite this ratio in relation to the price effect in the demand curve ( $\alpha_{1}$ ), use the demand equation (3.1.2a) to calculate $\operatorname{Cov}\left(p_{i}, q_{i}\right)$ as


\begin{equation*}
\operatorname{Cov}\left(p_{i}, q_{i}\right)=\alpha_{1} \operatorname{Var}\left(p_{i}\right)+\operatorname{Cov}\left(p_{i}, u_{i}\right) \tag{3.1.6}
\end{equation*}


Substituting (3.1.6) into (3.1.5), we obtain the expression for the asymptotic bias for $\alpha_{1}$


\begin{equation*}
\text { plim of the OLS estimate of the price coefficient }-\alpha_{1}=\frac{\operatorname{Cov}\left(p_{i}, u_{i}\right)}{\operatorname{Var}\left(p_{i}\right)} . \tag{3.1.7}
\end{equation*}


Similarly, the asymptotic bias for $\beta_{1}$, the price effect in the supply curve, is


\begin{equation*}
\text { plim of the OLS estimate of the price coefficient }-\beta_{1}=\frac{\operatorname{Cov}\left(p_{i}, v_{i}\right)}{\operatorname{Var}\left(p_{i}\right)} \text {. } \tag{3.1.8}
\end{equation*}


But, as seen from (3.1.4), $\operatorname{Cov}\left(p_{i}, u_{i}\right) \neq 0$ and $\operatorname{Cov}\left(p_{i}, v_{i}\right) \neq 0$, so the OLS estimate is consistent for neither $\alpha_{1}$ nor $\beta_{1}$. This phenomenon is known as the endogeneity bias. It is also known as the simultaneous equations bias or simultaneity bias, because the regressor and the error term are often related to each other through a system of simultaneous equations, as in the present example.

\footnotetext{${ }^{1}$ Fact: Let $\boldsymbol{\gamma}$ be the least squares coefficients in $\widehat{\mathrm{E}}^{*}(\boldsymbol{y} \mid \mathbf{I}, \mathbf{x})=\alpha+\mathbf{x}^{\prime} \boldsymbol{\gamma}$. Then $\boldsymbol{\gamma}=\operatorname{Var}(\mathbf{x})^{-1} \operatorname{Cov}(\mathbf{x}, \boldsymbol{y})$. Proving this was an analytical exercise for Chapter 2.
}In the extreme case of no demand shifters (so that $u_{i}=0$ for all $i$ ), we have $\operatorname{Cov}\left(p_{i}, u_{i}\right)=0$, and the formula (3.1.7) indicates that the OLS estimate is consistent for the demand parameter $\alpha_{1}$. In this case, the demand curve does not shift, and, as illustrated in Figure 3.1(a), all the observed combinations of price and quantity fall on the demand curve as the supply curve shifts. In the other extreme case of no supply shifters, the observed price-quantity pairs trace out the supply curve as the demand curve shifts (see Figure 3.1(b)). In the general case of both curves shifting, the OLS estimate should be consistent for a weighted average of $\alpha_{1}$ and $\beta_{1}$. This can be seen analytically by deriving yet another expression for the plim of the OLS estimate:


\begin{equation*}
\text { plim of the OLS estimate of the price coefficient }=\frac{\alpha_{1} \operatorname{Var}\left(v_{i}\right)+\beta_{1} \operatorname{Var}\left(u_{i}\right)}{\operatorname{Var}\left(v_{i}\right)+\operatorname{Var}\left(u_{i}\right)} . \tag{3.1.9}
\end{equation*}


Proving this is a review question.

\section*{Observable Supply Shifters}
The reason neither the demand curve nor the supply curve is consistently estimated in the general case is that we cannot infer from data whether the change in price and quantity is due to a demand shift or a supply shift. This suggests that it might be possible to estimate the demand curve (resp. the supply curve) if some of the factors that shift the supply curve (resp. the demand curve) are observable. So suppose the supply shifter, $v_{i}$, can be divided into an observable factor $x_{i}$ and an unobservable factor $\zeta_{i}$ uncorrelated with $x_{i} .^{2}$


\begin{equation*}
q_{i}=\beta_{0}+\beta_{1} p_{i}+\beta_{2} x_{i}+\zeta_{i} \quad \text { with } \beta_{2} \neq 0 . \quad \text { (supply) } \tag{3.1.2b́}
\end{equation*}


Now imagine that this observed supply shifter $x_{i}$ is predetermined (i.e., uncorrelated with the error term) in the demand equation; think of $x_{i}$ as the temperature in coffee-growing regions. If the temperature ( $x_{i}$ ) is uncorrelated with the unobserved taste for coffee ( $u_{i}$ ), it should be possible to extract from price movements a component that is related to the temperature (the observed supply shifter) but uncorrelated with the demand shifter. We can then estimate the demand curve by examining the relationship between coffee consumption and that component of price. Let us formalize this argument.

\footnotetext{${ }^{2}$ This decomposition is always possible. If the least squares projection of $v_{i}$ on a constant and $x_{i}$ is $\gamma_{0}+\beta_{2} x_{i}$. define $\zeta_{i} \equiv v_{i}-\gamma_{0}-\beta_{2} x_{i}$, so $v_{i}=\zeta_{i}+\gamma_{0}+\beta_{2} x_{i}$. By definition, $\zeta_{i}$ is uncorrelated with $x_{i}$. Substitute this equation into the original supply equation (3.1.2b) and submerge $\gamma_{0}$ in the intercept. This produces (3.1.2b').
}\begin{figure}[h]
\begin{center}
  \includegraphics[alt={},max width=\textwidth]{14fb0795-6d81-4e76-9788-792298f1ddae-207_1277_924_219_426}
\captionsetup{labelformat=empty}
\caption{Figure 3.1: Supply and Demand Curves}
\end{center}
\end{figure}

For the equation in question, a predetermined variable that is correlated with the endogenous regressor is called an instrumental variable or an instrument. We sometimes call it a valid instrument to emphasize that its correlation with the endogenous regressor is not zero. In this example, the observable supply shifter $x_{i}$ can serve as an instrument for the demand equation. This can be seen easily. Solve the system of two simultaneous equations (3.1.2a) and (3.1.2b') for ( $p_{i}, q_{i}$ ):


\begin{gather*}
p_{i}=\frac{\beta_{0}-\alpha_{0}}{\alpha_{1}-\beta_{1}}+\frac{\beta_{2}}{\alpha_{1}-\beta_{1}} x_{i}+\frac{\zeta_{i}-u_{i}}{\alpha_{1}-\beta_{1}},  \tag{3.1.10a}\\
q_{i}=\frac{\alpha_{1} \beta_{0}-\alpha_{0} \beta_{1}}{\alpha_{1}-\beta_{1}}+\frac{\alpha_{1} \beta_{2}}{\alpha_{1}-\beta_{1}} x_{i}+\frac{\alpha_{1} \zeta_{i}-\beta_{1} u_{i}}{\alpha_{1}-\beta_{1}} . \tag{3.1.10b}
\end{gather*}


Since $\operatorname{Cov}\left(x_{i}, \zeta_{i}\right)=0$ by construction and $\operatorname{Cov}\left(x_{i}, u_{i}\right)=0$ by assumption, it follows from (3.1.10a) that

$$
\operatorname{Cov}\left(x_{i}, p_{i}\right)=\frac{\beta_{2}}{\alpha_{1}-\beta_{1}} \operatorname{Var}\left(x_{i}\right) \neq 0
$$

So $x_{i}$ is indeed a valid instrument.\\
With a valid instrument at hand, we can estimate the price coefficient $\alpha_{1}$ of the demand curve consistently. Use the demand equation (3.1.2a) to calculate, not $\operatorname{Cov}\left(p_{i}, q_{i}\right)$ as in (3.1.6), but $\operatorname{Cov}\left(x_{i}, q_{i}\right)$ :

$$
\begin{aligned}
\operatorname{Cov}\left(x_{i}, q_{i}\right) & =\alpha_{1} \operatorname{Cov}\left(x_{i}, p_{i}\right)+\operatorname{Cov}\left(x_{i}, u_{i}\right) \\
& =\alpha_{1} \operatorname{Cov}\left(x_{i}, p_{i}\right) \quad\left(\operatorname{Cov}\left(x_{i}, u_{i}\right)=0 \text { by assumption }\right) .
\end{aligned}
$$

As we just verified, $\operatorname{Cov}\left(x_{i}, p_{i}\right) \neq 0$. So we can divide both sides by $\operatorname{Cov}\left(x_{i}, p_{i}\right)$ to obtain


\begin{equation*}
\alpha_{1}=\frac{\operatorname{Cov}\left(x_{i}, q_{i}\right)}{\operatorname{Cov}\left(x_{i}, p_{i}\right)} \tag{3.1.11}
\end{equation*}


A natural estimator that suggests itself is


\begin{equation*}
\hat{\alpha}_{1, \mathrm{IV}}=\frac{\text { sample covariance between } x_{i} \text { and } q_{i}}{\text { sample covariance between } x_{i} \text { and } p_{i}} . \tag{3.1.12}
\end{equation*}


This estimator is called the instrumental variables (IV) estimator with $x_{i}$ as the instrument. We sometimes say "the endogenous regressor $p_{i}$ is instrumented by $x_{i}$."

Another popular procedure for consistently estimating $\alpha_{1}$ is the two-stage least squares (2SLS). It is so called because the procedure consists of running two\\
regressions. In the first stage, the endogenous regressor $p_{i}$ is regressed on a constant and the predetermined variable $x_{i}$, to obtain the fitted value, $\hat{p}_{i}$, of $p_{i}$. The second stage is to regress the dependent variable $q_{i}$ on a constant and $\hat{p}_{i}$. The use of $\hat{p}_{i}$ rather than $p_{i}$ distinguishes 2 SLS from the simple-minded application of OLS to the demand equation. The 2SLS estimator of $\alpha_{1}$ is the OLS estimate of the $\hat{p}_{i}$ coefficient in the second-stage regression. So it equals ${ }^{3}$


\begin{equation*}
\hat{\alpha}_{1,2 \mathrm{SLS}}=\frac{\text { sample covariance between } \hat{p}_{i} \text { and } q_{i}}{\text { sample variance of } \hat{p}_{i}} . \tag{3.1.13}
\end{equation*}


To relate the regression in the second stage to the demand equation, rewrite (3.1.2a) as


\begin{equation*}
q_{i}=\alpha_{0}+\alpha_{1} \hat{p}_{i}+\left[u_{i}+\alpha_{1}\left(p_{i}-\hat{p}_{i}\right)\right] \tag{3.1.14}
\end{equation*}


The second stage regression estimates this equation, treating the bracketed term as the error term. The OLS estimate of $\alpha_{1}$ is consistent for the following reason. If the fitted value $\hat{p}_{i}$ were exactly equal to the least squares projection $\widehat{\mathrm{E}}^{*}\left(p_{i} \mid 1, x_{i}\right)$, then neither $u_{i}$ nor ( $p_{i}-\hat{p}_{i}$ ) would be correlated with $\hat{p}_{i}: u_{i}$ is uncorrelated because it is uncorrelated with $x_{i}$ and $\hat{p}_{i}$ is a linear function of $x_{i}$, and ( $p_{i}-\hat{p}_{i}$ ) is uncorrelated because it is a least squares projection error. The fitted value $\hat{p}_{i}$ is not exactly equal to $\widehat{\mathrm{E}}^{*}\left(p_{i} \mid 1, x_{i}\right)$, but the difference between the two vanishes as the sample gets larger. Therefore, asymptotically, $\hat{p}_{i}$ is uncorrelated with the error term in (3.1.14), making the 2SLS estimator consistent. Furthermore, since the projection is the best linear predictor, the fitted value incorporates all of the effect of the supply shifter on price. This suggests that minimizing the asymptotic variance is minimized for the 2SLS estimator.

In the present example, the IV and 2SLS estimators are numerically the same (we will prove this in a more general context). More generally, the 2SLS estimator can be written as an IV estimator with an appropriate choice of instruments, and the IV estimator, in turn, is a particular GMM estimator.

\section*{QUESTIONS FOR REVIEW}
\begin{enumerate}
  \item In the simultaneous equations model (3.1.2), suppose $\operatorname{Cov}\left(u_{i}, v_{i}\right)$ is not necessarily zero. Is it necessarily true that price and the demand shifter $u_{i}$ are positively correlated when $\alpha_{1}<0$ and $\beta_{1}>0$ ? [Answer: No.] Why?
\end{enumerate}

\footnotetext{${ }^{3}$ Fact: In the regression of $y_{i}$ on a constant and $x_{i}$, the OLS estimator of the $x_{i}$ coefficient is the ratio of the sample covariance between $x_{i}$ and $y_{i}$ to the sample variance of $x_{i}$. Proving this was a review question of Section 1.2.
}
2. (3.1.7) shows that the OLS estimate of the price coefficient in the regression of quantity on a constant and price is biased for the price coefficient $\alpha_{1}$. Is the OLS estimate of the intercept biased for $\alpha_{0}$, the intercept in the demand equation? Hint: The plim of the OLS estimate of the intercept in the regression is
$$
\mathrm{E}\left(q_{i}\right)-\frac{\operatorname{Cov}\left(p_{i}, q_{i}\right)}{\operatorname{Var}\left(p_{i}\right)} \mathrm{E}\left(p_{i}\right)
$$

But from (3.1.2a), $\mathrm{E}\left(q_{i}\right)=\alpha_{0}+\alpha_{1} \mathrm{E}\left(p_{i}\right)$.\\
3. Derive (3.1.9). Hint: Show:

$$
\operatorname{Var}\left(p_{i}\right)=\frac{\operatorname{Var}\left(v_{i}\right)+\operatorname{Var}\left(u_{i}\right)}{\left(\alpha_{1}-\beta_{1}\right)^{2}}, \quad \operatorname{Cov}\left(p_{i}, q_{i}\right)=\frac{\alpha_{1} \operatorname{Var}\left(v_{i}\right)+\beta_{1} \operatorname{Var}\left(u_{i}\right)}{\left(\alpha_{1}-\beta_{1}\right)^{2}} .
$$

\begin{enumerate}
  \setcounter{enumi}{3}
  \item For the market equilibrium model $(3.1 .2 \mathrm{a}),\left(3.1 .2 \mathrm{~b}^{\prime}\right)$ (on page 189) with $\operatorname{Cov}\left(u_{i}, \zeta_{i}\right)=0, \operatorname{Cov}\left(x_{i}, u_{i}\right)=0$, and $\operatorname{Cov}\left(x_{i}, \zeta_{i}\right)=0$, verify that price is endogenous in the demand equation. Is it in the supply equation? Hint: Look at (3.1.10a). Do we need the assumption that the demand and supply shifters are uncorrelated (i.e., $\operatorname{Cov}\left(u_{i}, \zeta_{i}\right)=0$ ) for $\hat{\alpha}_{1, \text { IV }}$ and $\hat{\alpha}_{1,2 \text { SLS }}$ to be consistent? Hint: Is $x_{i}$ a valid instrument without the assumption?
\end{enumerate}

\subsection*{3.2 More Examples}
Endogeneity bias arises in a wide variety of situations. We examine a few more examples.

\section*{A Simple Macroeconometric Model}
Haavelmo's (1943) illustrative model is an extremely simple macroeconometric model:

$$
\begin{aligned}
C_{i} & =\alpha_{0}+\alpha_{1} Y_{i}+u_{i}, \quad 0<\alpha_{1}<1 \quad \text { (consumption function) } \\
Y_{i} & =C_{i}+I_{i} \quad(\mathrm{GNP} \text { identity }),
\end{aligned}
$$

where $C_{i}$ is aggregate consumption in year $i, Y_{i}$ is GNP, $I_{i}$ is investment, and $\alpha_{1}$ is the Marginal Propensity to Consume out of income (the MPC). As is familiar from\\
introductory macroeconomics, GNP in equilibrium is


\begin{equation*}
Y_{i}=\frac{\alpha_{0}}{1-\alpha_{1}}+\frac{I_{i}}{1-\alpha_{1}}+\frac{u_{i}}{1-\alpha_{1}} \tag{3.2.1}
\end{equation*}


If investment is predetermined in that $\operatorname{Cov}\left(I_{i}, u_{i}\right)=0$, it follows from (3.2.1) that

$$
\begin{aligned}
& \operatorname{Cov}\left(Y_{i}, u_{i}\right)=\frac{\operatorname{Var}\left(u_{i}\right)}{1-\alpha_{1}}>0 \\
& \operatorname{Cov}\left(I_{i}, Y_{i}\right)=\frac{\operatorname{Var}\left(I_{i}\right)}{1-\alpha_{1}}>0
\end{aligned}
$$

So income is endogenous in the consumption function, but investment is a valid instrument for the endogenous regressor. A straightforward calculation similar to the one we used to derive (3.1.7) shows that the OLS estimator of the MPC obtained from regressing consumption on a constant and income is asymptotically biased:


\begin{equation*}
\operatorname{plim} \hat{\alpha}_{1, \mathrm{OLS}}-\alpha_{1}=\frac{\operatorname{Cov}\left(Y_{i}, u_{i}\right)}{\operatorname{Var}\left(Y_{i}\right)}=\frac{1-\alpha_{1}}{1+\frac{\operatorname{Var}\left(I_{i}\right)}{\operatorname{Var}\left(u_{i}\right)}}>0 \tag{3.2.2}
\end{equation*}


This is perhaps the clearest example of simultaneity bias. The difference from Working's example of market equilibrium is that here the second equation is an identity, which only makes it easier to see the endogeneity of the regressor. The asymptotic bias can be corrected for by the use of investment as the instrument for income in the consumption function. The role played by the observable supply shifter in Working's example is played here by investment.

\section*{Errors-in-Variables}
The term errors-in-variables refers to the phenomenon that an otherwise predetermined regressor necessarily becomes endogenous when measured with error. This problem is ubiquitous, particularly in micro data on households. For example, in the Panel Study of Income Dynamics (PSID), information on variables such as food consumption and income is collected over the telephone. It is perhaps too much to hope that the respondent can recall on the spot how much was spent on food over a specified period of time.

The cross-section version of M. Friedman's (1957) Permanent Income Hypothesis can be formulated neatly as an errors-in-variables problem. The hypothesis states that "permanent consumption" $C_{i}^{*}$ for household $i$ is proportional to "permanent income" $Y_{i}^{*}$ :


\begin{equation*}
C_{i}^{*}=k Y_{i}^{*} \quad \text { with } 0<k<1 . \tag{3.2.3}
\end{equation*}


It is assumed that measured consumption $C_{i}$ and measured income $Y_{i}$ are errorridden measures of permanent consumption and income:


\begin{equation*}
C_{i}=C_{i}^{*}+c_{i} \quad \text { and } \quad Y_{i}=Y_{i}^{*}+y_{i} . \tag{3.2.4}
\end{equation*}


Measurement errors $c_{i}$ and $y_{i}$ are assumed to be zero mean and uncorrelated with permanent consumption and income:


\begin{gather*}
\mathrm{E}\left(c_{i}\right)=0, \mathrm{E}\left(y_{i}\right)=0, \mathrm{E}\left(c_{i} y_{i}\right)=0,  \tag{3.2.5}\\
\mathrm{E}\left(C_{i}^{*} c_{i}\right)=0, \mathrm{E}\left(Y_{i}^{*} y_{i}\right)=0, \mathrm{E}\left(C_{i}^{*} y_{i}\right)=0, \mathrm{E}\left(Y_{i}^{*} c_{i}\right)=0 . \tag{3.2.6}
\end{gather*}


Substituting (3.2.4) into (3.2.3), the relationship can be expressed in terms of measured consumption and income:


\begin{equation*}
C_{i}=k Y_{i}+u_{i} \quad \text { with } u_{i} \equiv c_{i}-k y_{i} . \tag{3.2.7}
\end{equation*}


This example differs from the previous ones in that the equation does not include a constant. So we should examine the cross moment $\mathrm{E}\left(Y_{i} u_{i}\right)$ rather than the covariance $\operatorname{Cov}\left(Y_{i}, u_{i}\right)$ to determine whether income is predetermined. It is straightforward to derive from (3.2.4)-(3.2.7) that

$$
\mathrm{E}\left(Y_{i} u_{i}\right)=-k \mathrm{E}\left(y_{i}^{2}\right)<0
$$

So measured income is endogenous in the consumption function (3.2.7). Unlike in the previous examples, the endogeneity is due to measurement errors. Using the fact that the OLS estimator of the $Y_{i}$ coefficient in (3.2.7) is consistent for the corresponding least squares projection coefficient $\mathrm{E}\left(C_{i} Y_{i}\right) / \mathrm{E}\left(Y_{i}^{2}\right)$, we can also derive from (3.2.4)-(3.2.7) that


\begin{equation*}
\operatorname{plim} \widehat{k}_{\mathrm{OLS}}=\frac{k \mathrm{E}\left[\left(Y_{i}^{*}\right)^{2}\right]}{\mathrm{E}\left[\left(Y_{i}^{*}\right)^{2}\right]+\mathrm{E}\left(y_{i}^{2}\right)}<k . \tag{3.2.8}
\end{equation*}


So the regression of measured consumption on measured income (without the intercept) underestimates $k$. Friedman used this result to explain why the MPC from the cross-section regression of consumption on income is lower than the MPC estimated from aggregate time-series data.

Let us suppose for a moment that there exists a valid instrument $x_{i}$. So $\mathrm{E}\left(x_{i} u_{i}\right)=0$ and $\mathrm{E}\left(x_{i} Y_{i}\right) \neq 0$. A similar argument we used for deriving (3.1.11) establishes that


\begin{equation*}
k=\frac{\mathrm{E}\left(x_{i} C_{i}\right)}{\mathrm{E}\left(x_{i} Y_{i}\right)} \tag{3.2.9}
\end{equation*}


So a consistent IV estimator is


\begin{equation*}
\widehat{k}_{\mathrm{IV}}=\frac{\text { sample mean of } x_{i} C_{i}}{\text { sample mean of } x_{i} Y_{i}} \tag{3.2.10}
\end{equation*}


But is there a valid instrument? Yes, and it is a constant. Substituting $x_{i}=1$ into (3.2.10), we obtain a consistent estimator of $k$ which is the ratio of the sample mean of measured consumption to the sample mean of measured income. This is how Friedman estimated $k$.

\section*{Production Function}
In many contexts, the error term includes factors that are observable to the economic agent under study but unobservable to the econometrician. Endogeneity arises when regressors are decisions made by the agent on the basis of such factors. Consider a cross-sectional sample of firms choosing labor input to maximize profits. The production function of the $i$-th firm is


\begin{equation*}
Q_{i}=A_{i} \cdot\left(L_{i}\right)^{\phi_{1}} \cdot \exp \left(v_{i}\right), \quad 0<\phi_{1}<1, \tag{3.2.11}
\end{equation*}


where $Q_{i}$ is the firm's output, $L_{i}$ is labor input, $A_{i}$ is the firm's efficiency level known to the firm, and $v_{i}$ is the technology shock. In contrast to $A_{i}, v_{i}$ is not observable to the firm when it chooses $L_{i}$. Neither $A_{i}$ nor $v_{i}$ is observable to the econometrician.

Assume that, for each firm $i, v_{i}$ is serially independent over time. Therefore, $B \equiv \mathrm{E}\left[\exp \left(v_{i}\right)\right]$ is the same for all $i,{ }^{4}$ and the level of output the firm expects when it chooses $L_{i}$ is

$$
A_{i} \cdot\left(L_{i}\right)^{\phi_{1}} \cdot B .
$$

Let $p$ be the output price and $w$ the wage rate. For simplicity, assume that all the firms are from the same competitive industry so that $p$ and $w$ are constant across firms. Firm $i$ 's objective is to choose $L_{i}$ to maximize expected profits

$$
p \cdot A_{i} \cdot\left(L_{i}\right)^{\phi_{1}} \cdot B-w \cdot L_{i}
$$

Take the derivative of this with respect to $L_{i}$, set it equal to zero, and solve it for

\footnotetext{${ }^{4}$ If $v_{i}$ for firm $i$ were correlated over time, the firm would use past values of $v_{i}$ to forecast the current value of $v_{i}$ when choosing labor input, and so $B$ would differ across firms. Also, since the expected value of a nonlinear function of a random variable whose mean is zero is not necessarily zero, $B$ is not necessarily zero even if $E\left(v_{i}\right)=0$.
}
$L_{i}$, to obtain the profit-maximizing labor input level:

\begin{equation*}
L_{i}=\left(\frac{w}{p}\right)^{\frac{1}{\phi_{1}-1}}\left(A_{i} B \phi_{1}\right)^{\frac{1}{1-\phi_{1}}} \tag{3.2.12}
\end{equation*}


Let $u_{i}$ be firm $i$ 's deviation from the mean $\log$ efficiency: $u_{i} \equiv \log \left(A_{i}\right)- \mathrm{E}\left[\log \left(A_{i}\right)\right]$, and let $\phi_{0} \equiv \mathrm{E}\left[\log \left(A_{i}\right)\right]$ (so $\mathrm{E}\left(u_{i}\right)=0$ and $A_{i}=\exp \left(\phi_{0}+u_{i}\right)$ ). Then, the production function (3.2.11) and the labor demand function (3.2.12) in logs can be written as


\begin{gather*}
\log \left(Q_{i}\right)=\phi_{0}+\phi_{1} \log \left(L_{i}\right)+\left(v_{i}+u_{i}\right)  \tag{3.2.13}\\
\log \left(L_{i}\right)=\beta_{0}+\frac{1}{1-\phi_{1}} u_{i} \tag{3.2.14}
\end{gather*}


where

$$
\beta_{0}=\frac{1}{\phi_{1}-1}\left[\log (w / p)-\phi_{0}-\log \left(\phi_{1} B\right)\right]
$$

Because all firms face the same prices, $\beta_{0}$ is constant across firms. (3.2.14) shows that, in the $\log$-linear production function (3.2.13), $\log \left(L_{i}\right)$ is an endogenous regressor positively correlated with the error term $\left(v_{i}+u_{i}\right)$ through $u_{i}$. Thus, the OLS estimator of $\phi_{1}$ in the estimation of the log-linear production function confounds the contribution to output of $u_{i}$ with labor's contribution. This example illustrates yet another source of endogeneity: a variable chosen by the agent taking into account some error component unobservable to the econometrician.

\section*{QUESTIONS FOR REVIEW}
\begin{enumerate}
  \item In Friedman's Permanent Income Hypothesis, consider the regression of $C_{i}$ on a constant and $Y_{i}$. Derive the plim of the OLS estimator of the $Y_{i}$ coefficient in this regression with a constant. Hint: The plim equals the ratio of $\operatorname{Cov}\left(C_{i}, Y_{i}\right)$ to $\operatorname{Var}\left(Y_{i}\right)$. Show that it equals
\end{enumerate}

$$
\frac{k \operatorname{Var}\left(Y_{i}^{*}\right)}{\operatorname{Var}\left(Y_{i}^{*}\right)+\operatorname{Var}\left(y_{i}\right)}
$$

\begin{enumerate}
  \setcounter{enumi}{1}
  \item In the production function example, show that $\operatorname{plim}_{n \rightarrow \infty} \hat{\phi}_{1, \mathrm{OLS}}=1$, where $\hat{\phi}_{1, \mathrm{OLS}}$ is the OLS estimate of $\phi_{1}$ from (3.2.13). Hint: Eliminate $u_{i}$ from the log output equation (3.2.13) by using the labor demand equation (3.2.14).
  \item In the production function example, suppose the firm can observe $v_{i}$ as well as $u_{i}$ before deciding on labor input. How does the demand equation for labor\\
(3.2.12) change? Show that $\log \left(Q_{i}\right)$ and $\log \left(L_{i}\right)$ are perfectly correlated. Hint: $\log \left(Q_{i}\right)$ will be an exact linear function of $\log \left(L_{i}\right)$ without errors.
\end{enumerate}

\subsection*{3.3 The General Formulation}
We now provide a general formulation. It is a model described by the following set of assumptions and is a generalization of the model of Chapter 2.

\section*{Regressors and Instruments}
Assumption 3.1 (linearity): The equation to be estimated is linear:

$$
y_{i}=\mathbf{z}_{i}^{\prime} \boldsymbol{\delta}+\varepsilon_{i} \quad(i=1,2, \ldots, n),
$$

where $\mathbf{z}_{i}$ is an $L$-dimensional vector of regressors, $\delta$ is an $L$-dimensional coefficient vector, and $\varepsilon_{i}$ is an unobservable error term.

Assumption 3.2 (ergodic stationarity): Let $\mathbf{x}_{i}$ be a $K$-dimensional vector to be referred to as the vector of instruments, and let $\mathbf{w}_{i}$ be the unique and nonconstant elements of $\left(y_{i}, \mathbf{z}_{i}, \mathbf{x}_{i}\right) .^{5}\left\{\mathbf{w}_{i}\right\}$ is jointly stationary and ergodic.

Assumption 3.3 (orthogonality conditions): All the $K$ variables in $\mathbf{x}_{i}$ are predetermined in the sense that they are all orthogonal to the current error term: $\mathrm{E}\left(x_{i k} \varepsilon_{i}\right)=0$ for all $i$ and $k(k=1,2, \ldots, K) .^{6}$ This can be written as

$$
\mathrm{E}\left[\mathbf{x}_{i} \cdot\left(y_{i}-\mathbf{z}_{i}^{\prime} \boldsymbol{\delta}\right)\right]=\mathbf{0} \quad \text { or } \quad \mathrm{E}\left(\mathbf{g}_{i}\right)=\mathbf{0},
$$

where $\mathbf{g}_{i} \equiv \mathbf{x}_{i} \cdot \varepsilon_{i}$.

We remarked in Section 2.3 that, when the regressors include a constant, the orthogonality conditions about the regressors are equivalent to the condition that $\mathrm{E}\left(\varepsilon_{i}\right)=$ 0 and that the regressors be uncorrelated with the error term. Here, the orthogonality conditions are about the instrumental variables. So if one of the instruments is

\footnotetext{${ }^{5}$ See examples below for illustrations of $\mathbf{w}_{i}$.\\
${ }^{6}$ As was noted in Section 2.3, our use of the term predetermined may not be standard in some quarters of the profession. All we require for the term is that the current error term be orthogonal to the current regressors. We do not require that the current error term be orthogonal to the past regressors.
}
a constant, Assumption 3.3 is equivalent to the condition that $\mathrm{E}\left(\varepsilon_{i}\right)=0$ and that nonconstant instruments be uncorrelated with the error term.

The examples examined in the previous two sections can be written in this general format. For example,

Example 3.1 (Working's example with an observable supply shifter):\\
Consider the market equilibrium model of (3.1.2a) and (3.1.2b') on page 189. Suppose the equation to be estimated is the demand equation. The example can be cast in the general format by setting

$$
y_{i}=q_{i}, L=2, \delta=\left[\begin{array}{l}
\alpha_{0} \\
\alpha_{1}
\end{array}\right], \mathbf{z}_{i}=\left[\begin{array}{c}
1 \\
p_{i}
\end{array}\right], \varepsilon_{i}=u_{i}, K=2, \mathbf{x}_{i}=\left[\begin{array}{c}
1 \\
x_{i}
\end{array}\right],
$$

and $\mathbf{w}_{i}=\left(q_{i}, p_{i}, x_{i}\right)^{\prime}$. Since the mean of the error term is zero, a constant is orthogonal to the error term and so can be included in $\mathbf{x}_{i}$. It follows that $x_{i}$, which is assumed to be uncorrelated with the error term, satisfies the orthogonality condition $\mathrm{E}\left(x_{i} \varepsilon_{i}\right)=0$. So it too can be included in $\mathbf{x}_{i}$.

The other examples of the previous sections can be similarly written as special cases of the general model.

Having two separate symbols, $\mathbf{x}_{i}$ and $\mathbf{z}_{i}$, may give you the impression that the regressors and the instruments do not share the same variables, but that is not always the case. Indeed in the above example, $\mathbf{x}_{i}$ and $\mathbf{z}_{i}$ share the same variable (a constant). The instruments that are also regressors are called predetermined regressors, and the rest of the regressors, those that are not included in $\mathbf{x}_{i}$, are called endogenous regressors. A good example to make the point is

Example 3.2 (wage equation): A simplified version of the wage equation to be estimated later in this chapter is

$$
L W_{i}=\delta_{1}+\delta_{2} S_{i}+\delta_{3} E X P R_{i}+\delta_{4} I Q_{i}+\varepsilon_{i}
$$

where $L W_{i}$ is the log wage of individual $i, S_{i}$ is completed years of schooling, $E X P R_{i}$ is experience in years, $I Q_{i}$ is IQ, and $\varepsilon_{i}$ is unobservable individual characteristics relevant for the wage rate. We assume $\mathrm{E}\left(\varepsilon_{i}\right)=0$ (if not, include the mean of $\varepsilon_{i}$ in $\delta_{1}$ ). In one of the specifications we estimate later, we assume that $S_{i}$ is predetermined but that $I Q_{i}$, an error-ridden measure of the individual's ability, is endogenous due to the errors-in-variables problem. We\\
also assume that $E X P R_{i}, A G E_{i}$ (age of the individual), and $M E D_{i}$ (mother's education in years) are predetermined. $A G E_{i}$ is excluded from the wage equation, reflecting the underlying assumption that, once experience is controlled for, age has no effect on the wage rate. In terms of the general model,

$$
y_{i}=L W_{i}, \quad L=4, \quad \mathbf{z}_{i}=\left[\begin{array}{c}
1 \\
S_{i} \\
E X P R_{i} \\
I Q_{i}
\end{array}\right], \quad K=5, \quad \mathbf{x}_{i}=\left[\begin{array}{c}
1 \\
S_{i} \\
E X P R_{i} \\
A G E_{i} \\
M E D_{i}
\end{array}\right],
$$

and $\mathbf{w}_{i}=\left(L W_{i}, S_{i}, E X P R_{i}, I Q_{i}, A G E_{i}, M E D_{i}\right)^{\prime}$. As in Example 3.1, we can include a constant in $\mathbf{x}_{i}$ because $\mathrm{E}\left(\varepsilon_{i}\right)=0$. In this example, $\mathbf{x}_{i}$ and $\mathbf{z}_{i}$ share three variables ( $1, S_{i}, E X P R_{i}$ ). If, for example, $S_{i}$ were not included in $\mathbf{x}_{i}$, that would amount to treating $S_{i}$ as endogenous.

If, as in this example, some of the regressors $\mathbf{z}_{i}$ are predetermined, those predetermined regressors should be included in the list of instruments $\mathbf{x}_{i}$. The GMM estimation of the parameter vector is about how to exploit the information afforded by the orthogonality conditions. Not including predetermined regressors as elements of $\mathbf{x}_{i}$ amounts to throwing away the orthogonality conditions that could have been exploited.

\section*{Identification}
As seen from the examples of the previous two sections, an instrument must be not only predetermined (i.e., orthogonal to the error term) but also correlated with the regressors. Otherwise, the instrumental variables estimator cannot be defined (see, e.g., (3.1.11)). The generalization to the case with more than one regressor and more than one predetermined variable is

Assumption 3.4 (rank condition for identification): The $K \times L$ matrix $\mathrm{E}\left(\mathbf{x}_{i} \mathbf{z}_{i}^{\prime}\right)$ is of full column rank (i.e., its rank equals $L$, the number of its columns). We denote this matrix by $\mathbf{\Sigma}_{\mathbf{x z}} \cdot{ }^{7}$

To see that this is indeed a generalization, consider Example 3.1. With $\mathbf{z}_{i}= \left(1, p_{i}\right)^{\prime}, \mathbf{x}_{i}=\left(1, x_{i}\right)^{\prime}$, the $\boldsymbol{\Sigma}_{\mathbf{x z}}$ matrix is

\footnotetext{${ }^{7}$ So the cross moment $\mathrm{E}\left(\mathbf{x}_{i} \mathbf{z}_{i}^{\prime}\right)$ is assumed to exist and is finite. If a moment is indicated, as here, then by implication the moment is assumed to exist and is finite.
}
$$
\boldsymbol{\Sigma}_{\mathbf{x z}}=\left[\begin{array}{cc}
1 & \mathrm{E}\left(p_{i}\right) \\
\mathrm{E}\left(x_{i}\right) & \mathrm{E}\left(x_{i} p_{i}\right)
\end{array}\right] .
$$

The determinant of this matrix is not zero (and hence is of full column rank) if and only if $\operatorname{Cov}\left(x_{i}, p_{i}\right)=\mathrm{E}\left(x_{i} p_{i}\right)-\mathrm{E}\left(x_{i}\right) \mathrm{E}\left(p_{i}\right) \neq 0$.

Assumption 3.4 is called the rank condition for identification for the following reason. To emphasize the dependence of $\mathbf{g}_{i}\left(\equiv \mathbf{x}_{i} \cdot \varepsilon_{i}\right)$ on the data and the parameter vector, rewrite $\mathbf{g}_{i}$ as


\begin{equation*}
\mathbf{g}_{i}=\mathbf{g}\left(\mathbf{w}_{i} ; \delta\right) \equiv \mathbf{x}_{i} \cdot\left(y_{i}-\mathbf{z}_{i}^{\prime} \delta\right) \tag{3.3.1}
\end{equation*}


So the orthogonality conditions can be rewritten as


\begin{equation*}
\mathrm{E}\left[\mathbf{g}\left(\mathbf{w}_{i} ; \delta\right)\right]=\underset{(K \times 1)}{\mathbf{0}} \tag{3.3.2}
\end{equation*}


Now let $\tilde{\delta}(L \times 1)$ be a hypothetical value of $\delta$ and consider the system of $K$ simultaneous equations in $L$ unknowns (the $L$ elements of $\tilde{\delta}$ ):


\begin{equation*}
\mathrm{E}\left[\mathbf{g}\left(\mathbf{w}_{i} ; \tilde{\boldsymbol{\delta}}\right)\right]=\underset{(K \times 1)}{\mathbf{0}} \tag{3.3.3}
\end{equation*}


The orthogonality conditions (3.3.2) mean that the true value of the coefficient vector, $\delta$, is a solution to this system of $K$ simultaneous equations (3.3.3). So the assumptions we have made so far guarantee that there exists a solution to the simultaneous equations system. We say that the coefficient vector (or the equation) is identified if $\tilde{\delta}=\delta$ is the only solution.

Because the estimation equation is linear in our model, the function $\mathbf{g}\left(\mathbf{w}_{i} ; \tilde{\delta}\right)$ is linear in $\tilde{\boldsymbol{\delta}}$, as it can be written as $\mathbf{x}_{i} \cdot y_{i}-\mathbf{x}_{i} \mathbf{z}_{i}^{\prime} \tilde{\boldsymbol{\delta}}$. So (3.3.3) is a system of $K$ linear equations:


\begin{equation*}
\mathrm{E}\left(\mathbf{x}_{i} \cdot y_{i}\right)-\mathrm{E}\left(\mathbf{x}_{i} \mathbf{z}_{i}^{\prime}\right) \tilde{\boldsymbol{\delta}}=\mathbf{0} \quad \text { or } \quad \boldsymbol{\Sigma}_{(K \mathbf{x} \mathbf{z})(L \times 1)}^{\tilde{\delta}}=\underset{(K \times 1)}{\sigma_{\mathbf{x y}}} \tag{3.3.4}
\end{equation*}


where

$$
\sigma_{\mathbf{x y}} \equiv \mathrm{E}\left(\mathbf{x}_{i} \cdot y_{i}\right), \boldsymbol{\Sigma}_{\mathbf{x z}} \equiv \mathrm{E}\left(\mathbf{x}_{i} \mathbf{z}_{i}^{\prime}\right)
$$

A necessary and sufficient condition that $\tilde{\delta}=\delta$ is the only solution can be derived from the following result from matrix algebra. ${ }^{8}$

Suppose there exists a solution to a system of linear simultaneous equations

\footnotetext{${ }^{8}$ See, e.g.. Searle (1982, pp. 233-234).
}
in $\mathbf{x}: \mathbf{A x}=\mathbf{b}$. A necessary and sufficient condition that it is the only solution is that $\mathbf{A}$ is of full column rank.

Therefore, $\tilde{\delta}=\delta$ is the only solution to (3.3.4) if and only if $\boldsymbol{\Sigma}_{\mathbf{x z}}$ is of full column rank, which is Assumption 3.4.

\section*{Order Condition for Identification}
Since $\operatorname{rank}\left(\boldsymbol{\Sigma}_{\mathbf{x z}}\right)<L$ if $K<L$, a necessary condition for identification is that


\begin{equation*}
K(=\# \text { predetermined variables }) \geq L(=\# \text { regressors }) . \tag{3.3.5a}
\end{equation*}


This is called the order condition for identification. It can be stated in different ways. Since $K$ is also the number of orthogonality conditions and $L$ is the number of parameters, the order condition \href{http://can.be}{can.be} stated equivalently as


\begin{equation*}
\text { \#orthogonality conditions } \geq \text { \#parameters. } \tag{3.3.5b}
\end{equation*}


By subtracting the number of predetermined regressors from both sides of the inequality, we obtain another equivalent statement:

\begin{displayquote}
\#predetermined variables excluded from the equation $\geq$ \#endogenous regressors.
\end{displayquote}

Depending on whether the order condition is satisfied, we say that the equation is

\begin{itemize}
  \item overidentified if the rank condition is satisfied and $K>L$,
  \item exactly identified or just identified if the rank condition is satisfied and $K=L$,
  \item underidentified (or not identified) if the order condition is not satisfied (i.e., if $K<L)$.
\end{itemize}

Since the order condition is a necessary condition for identification, a failure of the order condition means that the equation is not identified.

\section*{The Assumption for Asymptotic Normality}
As in Chapter 2, we need to strengthen the orthogonality conditions for the estimator to be asymptotically normal.

Assumption 3.5 ( $\mathbf{g}_{i}$ is a martingale difference sequence with finite second moments): Let $\mathbf{g}_{i} \equiv \mathbf{x}_{i} \cdot \varepsilon_{i}$. $\left\{\mathbf{g}_{i}\right\}$ is a martingale difference sequence (so a for-\\
tiori $\mathrm{E}\left(\mathbf{g}_{i}\right)=\mathbf{0}$ ). The $K \times K$ matrix of cross moments, $\mathrm{E}\left(\mathbf{g}_{i} \mathbf{g}_{i}^{\prime}\right)$, is nonsingular. We use $\mathbf{S}$ for $\operatorname{Avar}(\overline{\mathbf{g}})$ (i.e., the variance of the limiting distribution of $\sqrt{n} \overline{\mathbf{g}}$, where $\overline{\mathbf{g}} \equiv \frac{1}{n} \sum_{i=1}^{n} \mathbf{g}_{i}$ ). By Assumption 3.2 and the ergodic stationary Martingale Differences CLT, $\mathbf{S}=\mathrm{E}\left(\mathbf{g}_{i} \mathbf{g}_{i}^{\prime}\right)$.

This is the same as Assumption 2.5, and the same comments apply:

\begin{itemize}
  \item If the instruments include a constant, then this assumption implies that the error is a martingale difference sequence (and a fortiori serially uncorrelated).
  \item A sufficient and perhaps easier to understand condition for Assumption 3.5 is that
\end{itemize}


\begin{equation*}
\mathrm{E}\left(\varepsilon_{i} \mid \varepsilon_{i-1}, \ldots, \varepsilon_{1}, \mathbf{x}_{i}, \mathbf{x}_{i-1}, \ldots, \mathbf{x}_{1}\right)=0, \tag{3.3.6}
\end{equation*}


which means that, besides being a martingale difference sequence, the error term is orthogonal not only to the current but also to the past instruments.

\begin{itemize}
  \item Since $\mathbf{g}_{i} \mathbf{g}_{i}^{\prime}=\varepsilon_{i}^{2} \mathbf{x}_{i} \mathbf{x}_{i}^{\prime}$, $\mathbf{S}$ is a matrix of fourth moments. Consistent estimation of $\mathbf{S}$ will require a fourth-moment assumption to be specified in Assumption 3.6 below.
  \item If $\left\{\mathbf{g}_{i}\right\}$ is serially correlated, then $\mathbf{S}$ (which is defined to be $\operatorname{Avar}(\overline{\mathbf{g}})$ )) does not equal $\mathrm{E}\left(\mathbf{g}_{i} \mathbf{g}_{i}^{\prime}\right)$ and will take a more complicated form, as we will see in Chapter 6.
\end{itemize}

\section*{QUESTIONS FOR REVIEW}
\begin{enumerate}
  \item In the Working example of Example 3.1, is the demand equation identified? Overidentified? Is the supply equation identified?
  \item Suppose the rank condition is satisfied for the wage equation of Example 3.2. Is the equation overidentified?
  \item In the production function example, no instrument other than a constant is specified. So $K=1$ and $L=2$. The order condition informs us that the log output equation (3.2.13) cannot be identified. Write down the orthogonality condition and verify that there are infinitely many combinations of ( $\phi_{0}, \phi_{1}$ ) that satisfy the orthogonality condition.
  \item Verify that the examples of Section 3.2 are special cases of the model of this section by specifying ( $y_{i}, \mathbf{x}_{i}, \mathbf{z}_{i}$ ) and writing down the rank condition for each example.
  \item Show that the model implies
\end{enumerate}

$$
\operatorname{rank}\left(\underset{(K \times L)}{\boldsymbol{\Sigma}_{\mathbf{x z}}}\right)=\operatorname{rank}\left(\underset{(K \times L)}{\boldsymbol{\Sigma}_{\mathbf{x z}}} \vdots \underset{(K \times 1)}{\sigma_{\mathbf{x y}}}\right) .
$$

Hint: $\tilde{\delta}=\delta$ is a solution to (3.3.4). (Some textbooks add this equation to the rank condition for identification, but the equation is implied by the other assumptions of the model.)\\
6. (Irrelevant instruments) Consider adding another variable, call it $\xi_{i}$, to $\mathbf{x}_{i}$. Although it is predetermined, the variable is unrelated to the regressors in that $\mathrm{E}\left(\xi_{i} z_{i \ell}\right)=0$ for all $\ell(=1,2, \ldots, L)$. Is the rank condition still satisfied? Hint: If a $K \times L$ matrix is of full column rank, then adding any $L$-dimensional row vector to the rows of the matrix does not alter the rank.\\
7. (Linearly dependent instruments) In Example 3.2, suppose $A G E_{i}=E X P R_{i}+ S_{i}$ for all the individuals in the population. Does that necessarily mean a failure of the rank condition? [Answer: No.] Is the full-rank condition in Assumption 3.5 (that $\mathrm{E}\left(\mathbf{g}_{i} \mathbf{g}_{i}^{\prime}\right)$ be nonsingular) satisfied? Hint: There exists a $K$-dimensional vector $\boldsymbol{\alpha} \neq \mathbf{0}$ such that $\boldsymbol{\alpha}^{\prime} \mathbf{x}_{i}=0$. Show that $\boldsymbol{\alpha}^{\prime} \mathrm{E}\left(\mathbf{g}_{i} \mathbf{g}_{i}^{\prime}\right)=\mathbf{0}$.\\
8. (Linear combinations of instruments) Let $\mathbf{A}$ be a $q \times K$ matrix of full row rank (so $q \leq K$ ) such that $\mathbf{A} \boldsymbol{\Sigma}_{\mathbf{x z}}$ is of full column rank. Let $\hat{\mathbf{x}}_{i} \equiv \mathbf{A} \mathbf{x}_{i}$ (so $\tilde{\mathbf{x}}_{i}$ is a vector of $q$ transformed instruments). Verify: Assumptions $3.3-3.5$ hold for $\hat{\mathbf{x}}_{i}$ if they do for $\mathbf{x}_{i}$.\\
9. Verify that the model, consisting of Assumptions 3.1-3.5, reduces to the regression model of Chapter 2 if $\mathbf{z}_{i}=\mathbf{x}_{i}$.

\subsection*{3.4 Generalized Method of Moments Defined}
The orthogonality conditions state that a set of population moments are all equal to zero. The basic principle of the method of moments is to choose the parameter estimate so that the corresponding sample moments are also equal to zero. The population moments in the orthogonality conditions are $\mathrm{E}\left[\mathbf{g}\left(\mathbf{w}_{i} ; \delta\right)\right]$. Its sample analogue is the sample mean of $\mathbf{g}\left(\mathbf{w}_{i} ; \delta\right)$ (where $\mathbf{g}\left(\mathbf{w}_{i} ; \delta\right)$ is defined in (3.3.1))\\
evaluated at some hypothetical value $\tilde{\delta}$ of $\delta$ :


\begin{equation*}
\underset{(K \times 1)}{\mathbf{g}_{n}(\tilde{\boldsymbol{\delta}})} \equiv \frac{1}{n} \sum_{i=1}^{n} \mathbf{g}\left(\mathbf{w}_{i} ; \tilde{\boldsymbol{\delta}}\right) \tag{3.4.1}
\end{equation*}


Applying the method of moments principle to our model amounts to choosing a $\tilde{\boldsymbol{\delta}}$ that solves the system of $K$ simultaneous equations in $L$ unknowns: $\mathbf{g}_{n}(\tilde{\boldsymbol{\delta}})=\mathbf{0}$, which is the sample analog of (3.3.3). Because the estimation equation is linear, $\mathbf{g}_{n}(\tilde{\boldsymbol{\delta}})$ can be written as


\begin{align*}
\mathbf{g}_{n}(\tilde{\boldsymbol{\delta}}) & =\frac{1}{n} \sum_{i=1}^{n} \mathbf{x}_{i} \cdot\left(y_{i}-\mathbf{z}_{i}^{\prime} \tilde{\boldsymbol{\delta}}\right) \quad\left(\text { by definition (3.3.1) of } \mathbf{g}\left(\mathbf{w}_{i} ; \boldsymbol{\delta}\right)\right) \\
& =\frac{1}{n} \sum_{i=1}^{n} \mathbf{x}_{i} \cdot y_{i}-\left(\frac{1}{n} \sum_{i=1}^{n} \mathbf{x}_{i} \mathbf{z}_{i}^{\prime}\right) \tilde{\boldsymbol{\delta}} \equiv \mathbf{s}_{\mathbf{x y}}-\mathbf{S}_{\mathbf{x z}} \tilde{\boldsymbol{\delta}}, \tag{3.4.2}
\end{align*}


where $\mathbf{S}_{\mathbf{x y}}$ and $\mathbf{S}_{\mathbf{x z}}$ are the corresponding sample moments of $\boldsymbol{\sigma}_{\mathbf{x y}}$ and $\boldsymbol{\Sigma}_{\mathbf{x z}}$ :

$$
\underset{(K \times 1)}{\mathbf{S}_{\mathbf{x y}}} \equiv \frac{1}{n} \sum_{i=1}^{n} \mathbf{x}_{i} \cdot y_{i} \quad \text { and } \quad \underset{(K \times L)}{\mathbf{S}_{\mathbf{x z}}} \equiv \frac{1}{n} \sum_{i=1}^{n} \mathbf{x}_{i} \mathbf{z}_{i}^{\prime}
$$

So the sample analog $\mathbf{g}_{n}(\tilde{\boldsymbol{\delta}})=\mathbf{0}$ is a system of $K$ linear equations in $L$ unknowns:


\begin{equation*}
\mathbf{S}_{\mathbf{x z}} \tilde{\delta}=\mathbf{S}_{\mathbf{x y}} \tag{3.4.3}
\end{equation*}


This is the sample analogue of (3.3.4). If there are more orthogonality conditions than parameters, that is, if $K>L$, then the system may not have a solution. The extension of the method of moments to cover this case as well is the generalized method of moments (GMM).

\section*{Method of Moments}
If the equation is exactly identified, then $K=L$ and $\mathbf{\Sigma}_{\mathbf{x z}}$ is square and invertible. Since under Assumption $3.2 \mathbf{S}_{\mathbf{x z}}$ converges to $\boldsymbol{\Sigma}_{\mathbf{x z}}$ almost surely, $\mathbf{S}_{\mathbf{x z}}$ is invertible for sufficiently large sample size $n$ with probability one. Thus, when the sample size is large, the system of simultaneous equations (3.4.3) has a unique solution given by


\begin{equation*}
\hat{\delta}_{\mathrm{IV}}=\mathbf{S}_{\mathbf{x z}}^{-1} \mathbf{s}_{\mathbf{x y}}=\left(\frac{1}{n} \sum_{i=1}^{n} \mathbf{x}_{i} \mathbf{z}_{i}^{\prime}\right)^{-1} \frac{1}{n} \sum_{i=1}^{n} \mathbf{x}_{i} \cdot y_{i} \tag{3.4.4}
\end{equation*}


This estimator is also called the instrumental variables estimator with $\mathbf{x}_{i}$ serving as instruments. Because it is defined for the exactly identified case, the formula assumes that there are as many instruments as regressors. If, moreover, $\mathbf{z}_{i}=\mathbf{x}_{i}$ (all the regressors are predetermined or orthogonal to the error term), then $\hat{\delta}_{\text {IV }}$ reduces to the OLS estimator. Thus, the OLS estimator is a method of moments estimator.

\section*{Generalized Method of Moments}
If the equation is overidentified so that $K>L$, we cannot in general choose an $L$-dimensional $\tilde{\delta}$ to satisfy the $K$ equations in (3.4.3). If we cannot set $\mathbf{g}_{n}(\tilde{\delta})$ exactly equal to $\mathbf{0}$, we can at least choose $\tilde{\delta}$ so that $\mathbf{g}_{n}(\tilde{\delta})$ is as close to $\mathbf{0}$ as possible. To make precise what we mean by "close," we define the distance between any two $K$-dimensional vectors $\boldsymbol{\xi}$ and $\boldsymbol{\eta}$ by the quadratic form $(\boldsymbol{\xi}-\boldsymbol{\eta})^{\prime} \widehat{\mathbf{W}}(\boldsymbol{\xi}-\boldsymbol{\eta})$, where $\widehat{\mathbf{W}}$, sometimes called the weighting matrix, is a symmetric and positive definite matrix defining the distance. ${ }^{9}$

Definition 3.1 (GMM estimator): Let $\widehat{\mathbf{W}}$ be a $K \times K$ symmetric positive definite matrix, possibly dependent on the sample, such that $\widehat{\mathbf{W}} \rightarrow_{\mathrm{p}} \mathbf{W}$ as the sample size $n$ goes to infinity with $\mathbf{W}$ symmetric and positive definite. The GMM estimator of $\boldsymbol{\delta}$, denoted $\hat{\boldsymbol{\delta}}(\widehat{\mathbf{W}})$, is


\begin{equation*}
\hat{\boldsymbol{\delta}}(\widehat{\mathbf{W}}) \equiv \underset{\tilde{\delta}}{\operatorname{argmin}} J(\tilde{\boldsymbol{\delta}}, \widehat{\mathbf{W}}), \tag{3.4.5}
\end{equation*}


where

$$
J(\tilde{\boldsymbol{\delta}}, \widehat{\mathbf{W}}) \equiv n \cdot \mathbf{g}_{n}(\tilde{\boldsymbol{\delta}})^{\prime} \widehat{\mathbf{W}} \mathbf{g}_{n}(\tilde{\boldsymbol{\delta}})
$$

(The reason the distance $\mathbf{g}_{n}(\tilde{\boldsymbol{\delta}})^{\prime} \widehat{\mathbf{W}} \mathbf{g}_{n}(\tilde{\boldsymbol{\delta}})$ is multiplied by the sample size ( $n$ ) becomes clear in Section 3.6.) The weighting matrix is allowed to be random and depend on the sample size, to cover the possibility that the matrix is estimated from the sample. The definition makes clear that the GMM is a special case of minimum distance estimation. In minimum distance estimation, $\operatorname{plim} \mathbf{g}_{n}(\boldsymbol{\delta})=\mathbf{0}$, as here, but the $\mathbf{g}_{n}(\cdot)$ function is not necessarily a sample mean.

Since $\mathbf{g}_{n}(\tilde{\boldsymbol{\delta}})$ is linear in $\tilde{\boldsymbol{\delta}}$ (see (3.4.2)), the objective function is quadratic in $\tilde{\boldsymbol{\delta}}$ when the equation is linear:


\begin{equation*}
J(\tilde{\boldsymbol{\delta}}, \widehat{\mathbf{W}})=n \cdot\left(\mathbf{s}_{\mathbf{x y}}-\mathbf{S}_{\mathbf{x z}} \tilde{\boldsymbol{\delta}}\right)^{\prime} \widehat{\mathbf{W}}\left(\mathbf{s}_{\mathbf{x y}}-\mathbf{S}_{\mathbf{x z}} \tilde{\boldsymbol{\delta}}\right) . \tag{3.4.6}
\end{equation*}


\footnotetext{${ }^{9}$ Do not let the word "weight" confuse you between GMM and weighted least squares (WLS). In GMM, the weighting is applied to the sample mean $\overline{\mathbf{g}}$, while in WLS it applies to each observation.
}It is left to you to show that the first-order condition for the minimization of this with respect to $\tilde{\delta}$ is


\begin{equation*}
\underset{(L \times K)}{\mathbf{S}_{\mathbf{x z}}^{\prime}} \underset{(K \times K)}{\widehat{\mathbf{W}}} \underset{(K \times 1)}{\mathbf{S}_{\mathbf{x y}}}=\underset{(L \times K)}{\mathbf{S}_{\mathbf{x z}}^{\prime}} \underset{(K \times K)}{\widehat{\mathbf{W}}} \underset{(K \times L)}{\mathbf{S}_{\mathbf{x z}}} \underset{(L \times 1)}{\tilde{\boldsymbol{\delta}}} . \tag{3.4.7}
\end{equation*}


If Assumptions 3.2 and 3.4 hold, then $\mathbf{S}_{\mathbf{x z}}$ is of full column rank for sufficiently large $n$ with probability one. Then, since $\widehat{\mathbf{W}}$ is positive definite, the $L \times L$ matrix $\mathbf{S}_{\mathbf{x z}}^{\prime} \widehat{\mathbf{W}} \mathbf{S}_{\mathbf{x z}}$ is nonsingular. So the unique solution can be obtained by multiplying both sides by the inverse of $\mathbf{S}_{\mathbf{x z}}^{\prime} \widehat{\mathbf{W}} \mathbf{S}_{\mathbf{x z}}$. That unique solution is the GMM estimator:


\begin{equation*}
\text { GMM estimator: } \hat{\delta}(\widehat{\mathbf{W}})=\left(\mathbf{S}_{\mathbf{x z}}^{\prime} \widehat{\mathbf{W}} \mathbf{S}_{\mathbf{x z}}\right)^{-1} \mathbf{S}_{\mathbf{x z}}^{\prime} \widehat{\mathbf{W}} \mathbf{S}_{\mathbf{x y}} . \tag{3.4.8}
\end{equation*}


If $K=L$, then $\mathbf{S}_{\mathbf{x z}}$ is a square matrix and (3.4.8) reduces to the IV estimator (3.4.4). The GMM is indeed a generalization of the method of moments.

\section*{Sampling Error}
For later use, we derive the expression for the sampling error. Multiplying both sides of the estimation equation $y_{i}=\mathbf{z}_{i}^{\prime} \delta+\varepsilon_{i}$ from left by $\mathbf{x}_{i}$ and taking the averages yields


\begin{equation*}
\mathrm{s}_{\mathrm{xy}}=\mathrm{S}_{\mathrm{xz}} \delta+\overline{\mathrm{g}}, \tag{3.4.9}
\end{equation*}


where


\begin{equation*}
\overline{\mathbf{g}} \equiv \frac{1}{n} \sum_{i=1}^{n} \mathbf{x}_{i} \cdot \varepsilon_{i}=\frac{1}{n} \sum_{i=1}^{n} \mathbf{g}\left(\mathbf{w}_{i} ; \boldsymbol{\delta}\right)=\mathbf{g}_{n}(\boldsymbol{\delta}) . \tag{3.4.10}
\end{equation*}


Substituting (3.4.9) into (3.4.8), we obtain


\begin{equation*}
\hat{\delta}(\widehat{\mathbf{W}})-\delta=\left(\mathbf{S}_{\mathbf{x z}}^{\prime} \widehat{\mathbf{W}} \mathbf{S}_{\mathbf{x z}}\right)^{-1} \mathbf{S}_{\mathbf{x z}}^{\prime} \widehat{\mathbf{W}} \overline{\mathbf{g}} . \tag{3.4.11}
\end{equation*}


\section*{QUESTIONS FOR REVIEW}
\begin{enumerate}
  \item Verify that (3.1.12) is the IV estimator of $\alpha_{1}$ when the method of moments is applied to the demand equation (3.1.2a) with ( $1, x_{i}$ ) as instruments.
  \item If the equation is just identified, what is the minimized value of $J(\tilde{\boldsymbol{\delta}}, \widehat{\mathbf{W}})$ ?
  \item Even if the equation is overidentified, the population version, (3.3.4), of the system of $K$ equations (3.4.3) has a solution. Then, why does not the sample version, (3.4.3), have a solution? Hint: The population version has a solution\\
because the $K \times(L+1)$ matrix [ $\Sigma_{\mathbf{x z}} \vdots \sigma_{\mathbf{x y}}$ ] is of rank $L$. The rank condition is a set of equality conditions on the elements of $\boldsymbol{\Sigma}_{\mathbf{x z}}$ and $\boldsymbol{\sigma}_{\mathbf{x y}}$. Even if $\mathbf{S}_{\mathbf{x z}}$ and $\mathbf{s}_{\mathbf{x y}}$ converge to $\boldsymbol{\Sigma}_{\mathbf{x z}}$ and $\sigma_{\mathbf{x y}}$, respectively, it does not necessarily follow that [ $\mathbf{S}_{\mathbf{x z}} \vdots \mathbf{s}_{\mathbf{x y}}$ ] is of rank $L$ for sufficiently large $n$. In contrast, $\mathbf{S}_{\mathbf{x z}}$ is of full column rank for sufficiently large $n$ when its limit, $\boldsymbol{\Sigma}_{\mathbf{x z}}$ is of full column rank. This is because the full column rank condition for a matrix is a set of inequality conditions on the elements of the matrix.
  \item What is wrong with the following argument?
\end{enumerate}

Even if the equation is overidentified, there is no problem finding a solution to (3.4.3). Just premultiply both sides by $\mathbf{S}_{\mathbf{x z}}^{\prime}$ to obtain


\begin{equation*}
\mathbf{S}_{\mathbf{x z}}^{\prime} \mathbf{S}_{\mathbf{x z}} \tilde{\delta}=\mathbf{S}_{\mathbf{x z}}^{\prime} \mathbf{S}_{\mathbf{x y}} \tag{3.4.12}
\end{equation*}


Since $\mathbf{S}_{\mathbf{x z}}$ is of full column rank, $\mathbf{S}_{\mathbf{x z}}^{\prime} \mathbf{S}_{\mathbf{x z}}$ is invertible. So the solution is

$$
\tilde{\boldsymbol{\delta}}=\left(\mathbf{S}_{\mathbf{x z}}^{\prime} \mathbf{S}_{\mathbf{x z}}\right)^{-1} \mathbf{S}_{\mathbf{x z}}^{\prime} \mathbf{S}_{\mathbf{x y}} .
$$

Hint: This $\tilde{\delta}$ certainly solves (3.4.12), but does it solve (3.4.3)?\\
5. (Singular W) Verify that the GMM estimator (3.4.8) remains well-defined for sufficiently large $n$ even if $\mathbf{W}(\equiv \operatorname{plim} \widehat{\mathbf{W}})$ is singular, as long as $\boldsymbol{\Sigma}_{\mathbf{x z}}^{\prime} \mathbf{W} \boldsymbol{\Sigma}_{\mathbf{x z}}$ is nonsingular.

\subsection*{3.5 Large-Sample Properties of GMM}
The GMM formula (3.4.8) defines GMM estimators, which are a set of estimators, each indexed by the weighting matrix $\widehat{\mathbf{W}}$. You will be delighted to know that every estimator to be introduced in the next few chapters is a GMM estimator for some choice of $\widehat{\mathbf{W}}$. The task of this section is to develop the large-sample theory for the GMM estimator for any given choice of $\widehat{\mathbf{W}}$, which can be carried out quite easily with the techniques of the previous chapter. The first half of this section extends Propositions 2.1-2.4 of Chapter 2 to GMM estimators. One issue that did not arise in those propositions is which GMM estimators should be preferred to other GMM estimators. This is a question of choosing $\widehat{\mathbf{W}}$ optimally and will be taken up in the latter part of this section.

\section*{Asymptotic Distribution of the GMM Estimator}
The large-sample theory for $\hat{\delta}(\widehat{\mathbf{W}})$ valid for any given choice of $\widehat{\mathbf{W}}$ is

\section*{Proposition 3.1 (asymptotic distribution of the GMM estimator):}
(a) (Consistency) Under Assumptions 3.1-3.4, $\operatorname{plim}_{n \rightarrow \infty} \hat{\delta}(\widehat{\mathbf{W}})=\boldsymbol{\delta}$.\\
(b) (Asymptotic Normality) If Assumption 3.3 is strengthened as Assumption 3.5, then

$$
\sqrt{n}(\hat{\boldsymbol{\delta}}(\widehat{\mathbf{W}})-\boldsymbol{\delta}) \underset{\mathrm{d}}{\rightarrow} N(\mathbf{0} . \operatorname{Avar}(\hat{\boldsymbol{\delta}}(\widehat{\mathbf{W}}))) \quad \text { as } n \rightarrow \infty
$$

where

\[
\begin{array}{r}
\operatorname{Avar}(\hat{\delta}(\widehat{\mathbf{W}}))=\left(\boldsymbol{\Sigma}_{\mathbf{x z}}^{\prime} \mathbf{W} \boldsymbol{\Sigma}_{\mathbf{x z}}\right)^{-1} \boldsymbol{\Sigma}_{\mathbf{x z}}^{\prime} \mathbf{W S W} \boldsymbol{\Sigma}_{\mathbf{x z}}\left(\boldsymbol{\Sigma}_{\mathbf{x z}}^{\prime} \mathbf{W} \boldsymbol{\Sigma}_{\mathbf{x z}}\right)^{-1}  \tag{3.5.1}\\
\left(\text { Recall }: \boldsymbol{\Sigma}_{\mathbf{x z}} \equiv \mathrm{E}\left(\mathbf{x}_{i} \mathbf{z}_{i}^{\prime}\right), \mathbf{S}=\mathrm{E}\left(\mathbf{g}_{i} \mathbf{g}_{i}^{\prime}\right)=\mathrm{E}\left(\varepsilon_{i}^{2} \mathbf{x}_{i} \mathbf{x}_{i}^{\prime}\right), \mathbf{W} \equiv \operatorname{plim} \widehat{\mathbf{W}} .\right)
\end{array}
\]

(c) (Consistent Estimate of $\operatorname{Avar}(\hat{\delta}(\widehat{\mathbf{W}})))$ Suppose there is available a consistent estimator, $\widehat{\mathbf{S}}$, of $\mathbf{S}(K \times K)$. Then, under Assumption 3.2, $\operatorname{Avar}(\hat{\boldsymbol{\delta}}(\widehat{\mathbf{W}}))$ is consistently estimated by


\begin{equation*}
\widehat{\operatorname{Avar}(\hat{\delta}(\widehat{\mathbf{W}}))} \equiv\left(\mathbf{S}_{\mathbf{x z}}^{\prime} \widehat{\mathbf{W}} \mathbf{S}_{\mathbf{x z}}\right)^{-1} \mathbf{S}_{\mathbf{x z}}^{\prime} \widehat{\mathbf{W}} \widehat{\mathbf{S}} \widehat{\mathbf{W}} \mathbf{S}_{\mathbf{x z}}\left(\mathbf{S}_{\mathbf{x z}}^{\prime} \widehat{\mathbf{W}} \mathbf{S}_{\mathbf{x z}}\right)^{-1} \tag{3.5.2}
\end{equation*}


where $\mathbf{S}_{\mathbf{x z}}$ is the sample mean of $\mathbf{x}_{i} \mathbf{z}_{i}^{\prime}$ :

$$
\underset{(K \times L)}{\mathbf{S}_{\mathbf{x z}}} \equiv \frac{1}{n} \sum_{i=1}^{n} \mathbf{x}_{i} \mathbf{z}_{i}^{\prime}
$$

The ugly looking expression for the asymptotic variance will become much prettier when we choose the weighting matrix optimally. If you have gone through the proof of Proposition 2.1, you should find it a breeze to prove Proposition 3.1. The key observations are:\\
(1) $\mathbf{S}_{\mathbf{x z}} \rightarrow_{\mathrm{p}} \boldsymbol{\Sigma}_{\mathbf{x z}}$ (by ergodic stationarity)\\
(2) $\overline{\mathbf{g}}\left(\equiv \frac{1}{n} \sum_{i=1}^{n} \mathbf{g}_{i}\right) \rightarrow_{\mathrm{p}} \mathbf{0}$ (by ergodic stationarity and the orthogonality conditions)\\
(3) $\sqrt{n} \overline{\mathbf{g}} \rightarrow_{\mathrm{d}} N(\mathbf{0}, \mathbf{S})$ (by Assumption 3.5).

Consistency immediately follows if you use (1), (2), and Lemma 2.3(a) on the expression for the sampling error (3.4.11). To prove asymptotic normality, multiply\\
both sides of (3.4.11) by $\sqrt{n}$ to obtain


\begin{equation*}
\sqrt{n}(\hat{\boldsymbol{\delta}}(\widehat{\mathbf{W}})-\boldsymbol{\delta})=\left(\mathbf{S}_{\mathbf{x z}}^{\prime} \widehat{\mathbf{W}} \mathbf{S}_{\mathbf{x z}}\right)^{-1} \mathbf{S}_{\mathbf{x z}}^{\prime} \widehat{\mathbf{W}} \sqrt{n} \overline{\mathbf{g}} \tag{3.5.3}
\end{equation*}


and use (3), Lemma 2.3(a), and Lemma 2.4(c). Part (c) of Proposition 3.1 follows immediately from Lemma 2.3(a).

\section*{Estimation of Error Variance}
In Section 2.3, we proved the consistency of the OLS estimator, $s^{2}$, of the error variance. As was noted there, the result holds as long as the residual is from some consistent estimator of the coefficient vector. The same is true here.

Proposition 3.2 (consistent estimation of error variance): For any consistent estimator, $\hat{\boldsymbol{\delta}}$, of $\boldsymbol{\delta}$, define $\hat{\varepsilon}_{i} \equiv y_{i}-\mathbf{z}_{i}^{\prime} \hat{\boldsymbol{\delta}}$. Under Assumptions 3.1, 3.2, plus the assumption that $\mathrm{E}\left(\mathbf{z}_{i} \mathbf{z}_{i}^{\prime}\right)$ (second moments of the regressors) exists and is finite,

$$
\frac{1}{n} \sum_{i=1}^{n} \hat{\varepsilon}_{i}^{2} \underset{\mathrm{p}}{\rightarrow} \mathrm{E}\left(\varepsilon_{i}^{2}\right)
$$

provided $\mathrm{E}\left(\varepsilon_{i}^{2}\right)$ exists and is finite.\\
The proof is very similar to the proof of Proposition 2.2. The relationship between $\widehat{\varepsilon_{i}}$ and $\varepsilon_{i}$ is given by


\begin{equation*}
\hat{\varepsilon}_{i} \equiv y_{i}-\mathbf{z}_{i}^{\prime} \hat{\boldsymbol{\delta}}=\varepsilon_{i}-\mathbf{z}_{i}^{\prime}(\hat{\boldsymbol{\delta}}-\boldsymbol{\delta}) \tag{3.5.4}
\end{equation*}


so that


\begin{equation*}
\hat{\varepsilon}_{i}^{2}=\varepsilon_{i}^{2}-2(\hat{\boldsymbol{\delta}}-\boldsymbol{\delta})^{\prime} \mathbf{z}_{i} \cdot \varepsilon_{i}+(\hat{\boldsymbol{\delta}}-\boldsymbol{\delta})^{\prime} \mathbf{z}_{i} \mathbf{z}_{i}^{\prime}(\hat{\boldsymbol{\delta}}-\boldsymbol{\delta}) . \tag{3.5.5}
\end{equation*}


Summing over $i$, we obtain


\begin{equation*}
\frac{1}{n} \sum_{i=1}^{n} \hat{\varepsilon}_{i}^{2}=\frac{1}{n} \sum_{i=1}^{n} \varepsilon_{i}^{2}-2(\hat{\boldsymbol{\delta}}-\boldsymbol{\delta})^{\prime} \frac{1}{n} \sum_{i=1}^{n} \mathbf{z}_{i} \cdot \varepsilon_{i}+(\hat{\boldsymbol{\delta}}-\boldsymbol{\delta})^{\prime}\left(\frac{1}{n} \sum_{i=1}^{n} \mathbf{z}_{i} \mathbf{z}_{i}^{\prime}\right)(\hat{\boldsymbol{\delta}}-\boldsymbol{\delta}) . \tag{3.5.6}
\end{equation*}


As usual, $\frac{1}{n} \Sigma_{i} \varepsilon_{i}^{2} \rightarrow{ }_{\mathrm{p}} \mathrm{E}\left(\varepsilon_{i}^{2}\right)$. Since $\hat{\delta}$ is consistent and $\frac{1}{n} \Sigma_{i} \mathbf{Z}_{i} \mathbf{z}_{i}^{\prime}$ converges in probability to some finite matrix by assumption, the last term vanishes. It is easy to show that $\mathrm{E}\left(\mathbf{z}_{i} \cdot \varepsilon_{i}\right)$ exists and is finite. ${ }^{10}$ Then, by ergodic stationarity,

\footnotetext{${ }^{10}$ By the Cauchy-Schwartz inequality, $\mathrm{E}\left(\left|z_{i \ell} \varepsilon_{i}\right|\right) \leq \sqrt{\mathrm{E}\left(z_{i \ell}^{2}\right) \cdot \mathrm{E}\left(\varepsilon_{i}^{2}\right)}$, where $z_{i \ell}$ is the $\ell$-th element of $\mathbf{z}_{i}$. Both $E\left(z_{i \ell}^{2}\right)$ and $E\left(\varepsilon_{i}^{2}\right)$ are finite by assumption.
}
$$
\frac{1}{n} \Sigma_{i} \mathbf{z}_{i} \cdot \varepsilon_{i} \rightarrow \text { some finite vector. }
$$

So the second term on the RHS of (3.5.6), too, vanishes.

\section*{Hypothesis Testing}
It should now be routine to derive from Proposition 3.1(b) and 3.1(c) the following.

Proposition 3.3 (robust $t$-ratio and Wald statistics): Suppose Assumptions 3.13.5 hold, and suppose there is available a consistent estimate $\widehat{\mathbf{S}}$ of $\mathbf{S}(\equiv \operatorname{Avar}(\overline{\mathbf{g}})= \left.\mathrm{E}\left(\mathbf{g}_{i} \mathbf{g}_{i}^{\prime}\right)\right)$. Let

Then:\\
(a) Under the null hypothesis $\mathrm{H}_{0}: \delta_{\ell}=\bar{\delta}_{\ell}$,

$$
t_{\ell} \equiv \frac{\sqrt{n}\left(\hat{\delta}_{\ell}(\widehat{\mathbf{W}})-\bar{\delta}_{\ell}\right)}{\sqrt{(\operatorname{Avar}(\hat{\delta}(\widehat{\mathbf{W}})))_{\ell \ell}}}=\frac{\hat{\delta}_{\ell}(\widehat{\mathbf{W}})-\bar{\delta}_{\ell}}{S E_{\ell}^{*}} \underset{\mathrm{~d}}{\rightarrow} N(0,1)
$$

where $(\widehat{\operatorname{Avar}(\hat{\delta}(\widehat{\mathbf{W}}))})_{\ell \ell}$ is the $(\ell, \ell)$ element of $\widehat{\operatorname{Avar}(\hat{\delta}(\widehat{\mathbf{W}}))}$ [which is $\left.\widehat{\operatorname{Avar}\left(\hat{\delta}_{\ell}(\widehat{\mathbf{W}})\right)}\right)$ ] and


\begin{equation*}
S E_{\ell}^{*}(\text { robust standard error }) \equiv \sqrt{\frac{1}{n} \cdot(\operatorname{Avar}(\hat{\delta}(\widehat{\mathbf{W}})))_{\ell \ell}} \tag{3.5.7}
\end{equation*}


(b) Under the null hypothesis $\mathrm{H}_{0}: \mathbf{R} \boldsymbol{\delta}=\mathbf{r}$ where \# $\mathbf{r}$ is the number of restrictions (the dimension of $\mathbf{r}$ ) and $\mathbf{R}$ (\# $\mathbf{r} \times L$ ) is of full row rank,


\begin{equation*}
\left.W \equiv n \cdot(\mathbf{R} \hat{\delta}(\widehat{\mathbf{W}})-\mathbf{r})^{\prime}\{\mathbf{R}[\widehat{\operatorname{Avar}(\hat{\delta}(\widehat{\mathbf{W}}))})] \mathbf{R}^{\prime}\right\}^{-1}(\mathbf{R} \hat{\delta}(\widehat{\mathbf{W}})-\mathbf{r}) \underset{\mathrm{d}}{\rightarrow} \chi^{2}(\# \mathbf{r}) \tag{3.5.8}
\end{equation*}


(c) Under the null hypothesis $\mathrm{H}_{0}: \mathbf{a}(\boldsymbol{\delta})=\mathbf{0}$ such that $\mathbf{A}(\boldsymbol{\delta})$, the $\# \mathbf{a} \times L$ matrix of first derivatives of $\mathbf{a}(\boldsymbol{\delta})$ (where \#a is the dimension of $\mathbf{a}$ ), is continuous and of full row rank,


\begin{equation*}
\left.W \equiv n \cdot \mathbf{a}(\hat{\delta}(\widehat{\mathbf{W}}))^{\prime}\{\mathbf{A}(\hat{\delta}(\widehat{\mathbf{W}}))[\widehat{\operatorname{Avar}(\hat{\delta}(\widehat{\mathbf{W}}))})] \mathbf{A}(\hat{\delta}(\widehat{\mathbf{W}}))^{\prime}\right\}^{-1} \mathbf{a}(\hat{\delta}(\widehat{\mathbf{W}})) \underset{\mathrm{d}}{\rightarrow} \chi^{2}(\# \mathbf{a}) \tag{3.5.9}
\end{equation*}


For the Wald statistic for the nonlinear hypothesis, the same comment that we made for Proposition 2.3 applies here: the numerical value of the Wald statistic is not invariant to how the nonlinear restriction is represented by $\mathbf{a}(\cdot)$.

\section*{Estimation of S}
We have already studied the estimation of $\mathbf{S}\left(=\mathrm{E}\left(\mathbf{g}_{i} \mathbf{g}_{i}^{\prime}\right)=\mathrm{E}\left(\varepsilon_{i}^{2} \mathbf{x}_{i} \mathbf{x}_{i}^{\prime}\right)\right)$ in Section 2.5. The proposed estimator was, if adapted to the present estimation equation $y_{i}=\mathbf{z}_{i}^{\prime} \boldsymbol{\delta}+\varepsilon_{i}$,


\begin{equation*}
\widehat{\mathbf{S}} \equiv \frac{1}{n} \sum_{i=1}^{n} \hat{\varepsilon}_{i}^{2} \mathbf{x}_{i} \mathbf{x}_{i}^{\prime} \tag{3.5.10}
\end{equation*}


where $\hat{\varepsilon}_{i} \equiv y_{i}-\mathbf{z}_{i}^{\prime} \hat{\delta}$ and $\hat{\delta}$ is consistent for $\delta$. The fourth-moment assumption needed for this to be consistent for $\mathbf{S}$ is a generalization of Assumption 2.6.

Assumption 3.6 (finite fourth moments): $\mathrm{E}\left[\left(x_{i k} z_{i \ell}\right)^{2}\right]$ exists and is finite for all $k(=1, \ldots, K)$ and $\ell(=1, \ldots, L)$.

It is left as an analytical exercise to prove\\
Proposition 3.4 (consistent estimation of S): Suppose the coefficient estimate $\hat{\delta}$ used for calculating the residual $\hat{\varepsilon}_{i}$ for $\widehat{\mathbf{S}}$ in (3.5.10) is consistent, and suppose $\mathbf{S}=\mathrm{E}\left(\mathbf{g}_{i} \mathbf{g}_{i}^{\prime}\right)$ exists and is finite. Then, under Assumptions 3.1, 3.2, and 3.6, $\widehat{\mathbf{S}}$ given in (3.5.10) is consistent for S .

\section*{Efficient GMM Estimator}
Naturally, we wish to choose from the GMM estimators indexed by $\widehat{\mathbf{W}}$ one that has the least asymptotic variance. The next proposition provides a choice of $\mathbf{W}$ that minimizes the asymptotic variance.

Proposition 3.5 (optimal choice of the weighting matrix): A lower bound for the asymptotic variance of the GMM estimators (3.4.8) indexed by $\widehat{\mathbf{W}}$ is given by $\left(\boldsymbol{\Sigma}_{\mathbf{x z}}^{\prime} \mathbf{S}^{-1} \boldsymbol{\Sigma}_{\mathbf{x z}}\right)^{-1}$. The lower bound is achieved if $\widehat{\mathbf{W}}$ is such that $\mathbf{W}(\equiv \operatorname{plim} \widehat{\mathbf{W}})= \mathbf{S}^{-1} .^{11}$

Because the asymptotic variance for any given $\widehat{\mathbf{W}}$ is (3.5.1), this proposition is proved if we can show that

\footnotetext{${ }^{11}$ The condition that $\mathbf{W}=\mathbf{S}^{-1}$ is sufficient but not necessary for efficiency. A necessary and sufficient condition is that there exists a matrix $\mathbf{C}$ such that $\boldsymbol{\Sigma}_{\mathbf{x z}}^{\prime} \mathbf{W}=\mathbf{C} \boldsymbol{\Sigma}_{\mathbf{x z}}^{\prime} \mathbf{S}^{-1}$. See Newey and McFadden (1994, p. 2165).
}

\begin{equation*}
\left(\boldsymbol{\Sigma}_{\mathbf{x z}}^{\prime} \mathbf{W} \boldsymbol{\Sigma}_{\mathbf{x z}}\right)^{-1} \boldsymbol{\Sigma}_{\mathbf{x z}}^{\prime} \mathbf{W S W} \boldsymbol{\Sigma}_{\mathbf{x z}}\left(\boldsymbol{\Sigma}_{\mathbf{x z}}^{\prime} \mathbf{W} \boldsymbol{\Sigma}_{\mathbf{x z}}\right)^{-1} \geq\left(\boldsymbol{\Sigma}_{\mathbf{x z}}^{\prime} \mathbf{S}^{-1} \boldsymbol{\Sigma}_{\mathbf{x z}}\right)^{-1} \tag{3.5.11}
\end{equation*}

for any symmetric positive definite matrix $\mathbf{W}$. Proving this algebraic result is left as an analytical exercise.

A GMM estimator satisfying the efficiency condition that plim $\widehat{\mathbf{W}}=\mathbf{S}^{-1}$ will be called the efficient (or optimal) GMM estimator. Simply by replacing $\widehat{\mathbf{W}}$ by $\widehat{\mathbf{S}}^{-1}$ (which is consistent for $\mathbf{S}^{-1}$ ) in the formulas of Proposition 3.1, we obtain


\begin{align*}
\text { Efficient GMM estimator: } \hat{\delta}\left(\widehat{\mathbf{S}}^{-1}\right) & =\left(\mathbf{S}_{\mathbf{x z}}^{\prime} \widehat{\mathbf{S}}^{-1} \mathbf{S}_{\mathbf{x z}}\right)^{-1} \mathbf{S}_{\mathbf{x z}}^{\prime} \widehat{\mathbf{S}}^{-1} \mathbf{S}_{\mathbf{x y}}  \tag{3.5.12}\\
\operatorname{Avar}\left(\hat{\delta}\left(\widehat{\mathbf{S}}^{-1}\right)\right) & =\left(\boldsymbol{\Sigma}_{\mathbf{x z}}^{\prime} \mathbf{S}^{-1} \mathbf{\Sigma}_{\mathbf{x z}}\right)^{-1}  \tag{3.5.13}\\
\widehat{\operatorname{Avar}\left(\hat{\delta}\left(\widehat{\mathbf{S}}^{-1}\right)\right)} & =\left(\mathbf{S}_{\mathbf{x z}}^{\prime} \widehat{\mathbf{S}}^{-1} \mathbf{S}_{\mathbf{x z}}\right)^{-1} \tag{3.5.14}
\end{align*}


With $\widehat{\mathbf{W}}=\widehat{\mathbf{S}}^{-1}$, the formulas for the robust $t$ and the Wald statistics in Proposition 3.3 become


\begin{equation*}
t_{\ell}=\frac{\hat{\delta}_{\ell}\left(\widehat{\mathbf{S}}^{-1}\right)-\bar{\delta}_{\ell}}{S E_{\ell}^{*}}, \tag{3.5.15}
\end{equation*}


where $S E_{\ell}^{*}$ is the robust standard error given by

$$
S E_{\ell}^{*}=\sqrt{\frac{1}{n} \cdot\left(\left(\mathbf{S}_{\mathbf{x z}}^{\prime} \widehat{\mathbf{S}}^{-1} \mathbf{S}_{\mathbf{x z}}\right)^{-1}\right)_{\ell \ell}},
$$

and


\begin{equation*}
W=n \cdot \mathbf{a}\left(\hat{\delta}\left(\widehat{\mathbf{S}}^{-1}\right)\right)^{\prime}\left\{\mathbf{A}\left(\hat{\delta}\left(\widehat{\mathbf{S}}^{-1}\right)\right)\left(\mathbf{S}_{\mathbf{x z}}^{\prime} \widehat{\mathbf{S}}^{-1} \mathbf{S}_{\mathbf{x z}}\right)^{-1} \mathbf{A}\left(\hat{\delta}\left(\widehat{\mathbf{S}}^{-1}\right)\right)^{\prime}\right\}^{-1} \mathbf{a}\left(\hat{\delta}\left(\widehat{\mathbf{S}}^{-1}\right)\right) . \tag{3.5.16}
\end{equation*}


To calculate the efficient GMM estimator, we need the consistent estimator $\widehat{\mathbf{S}}$. But Proposition 3.4 assures us that the $\widehat{\mathbf{S}}$ based on any consistent estimator of $\boldsymbol{\delta}$ is consistent for $\mathbf{S}$. This leads us to the following two-step efficient GMM procedure:

Step 1: Choose a matrix $\widehat{\mathbf{W}}$ that converges in probability to a symmetric and positive definite matrix, and minimize $J(\tilde{\delta}, \widehat{\mathbf{W}})$ over $\tilde{\delta}$ to obtain $\hat{\delta}(\widehat{\mathbf{W}})$. There is no shortage of such matrices $\widehat{\mathbf{W}}$ (e.g., $\widehat{\mathbf{W}}=\mathbf{I}$ ), but usually we set $\widehat{\mathbf{W}}=\mathbf{S}_{\mathbf{x x}}^{-1}$. The resulting estimator $\hat{\boldsymbol{\delta}}\left(\mathbf{S}_{\mathbf{x x}}^{-\mathbf{1}}\right)$ is the celebrated two-stage least squares (as we will see in Section 3.8). Use this to calculate the residual $\hat{\varepsilon}_{i} \equiv y_{i}-\mathbf{z}_{i}^{\prime} \hat{\delta}(\widehat{\mathbf{W}})$ and obtain a consistent estimator $\widehat{\mathbf{S}}$ of $\mathbf{S}$ by (3.5.10).\\
Step 2: Minimize $J\left(\tilde{\delta}, \widehat{\mathbf{S}}^{-1}\right)$ over $\tilde{\delta}$. The minimizer is the efficient GMM estimator.

\section*{Asymptotic Power}
Like the coefficient estimator, the $t$ and Wald statistics depend on the choice of $\mathbf{W}$. It seems intuitively obvious that those statistics associated with the efficient GMM estimator should be preferred in large samples. This intuition can be formalized in terms of asymptotic power introduced in Section 2.4. Take, for example, the $t$-ratio for testing $\mathrm{H}_{0}: \delta_{\ell}=\bar{\delta}_{\ell}$. The $t$-ratio is written (reproducing (3.5.7)) as


\begin{equation*}
t_{\ell} \equiv \frac{\sqrt{n}\left(\hat{\delta}_{\ell}(\widehat{\mathbf{W}})-\bar{\delta}_{\ell}\right)}{\sqrt{(\operatorname{Avar}(\widehat{\boldsymbol{\delta}}(\widehat{\mathbf{W}})))_{\ell \ell}}} \tag{3.5.17}
\end{equation*}


The denominator converges to some finite number even when the null is false. In contrast, the numerator explodes (its plim is infinite) when the null is false. So the power under any fixed alternative approaches unity. That is, the test is consistent. This is true for any choice of $\mathbf{W}$, so test consistency cannot be a basis for choosing $\mathbf{W}$.

Next, consider a sequence of local alternatives subject to Pitman drift:


\begin{equation*}
\delta_{\ell}^{(n)}=\bar{\delta}_{\ell}+\frac{\gamma}{\sqrt{n}} \tag{3.5.18}
\end{equation*}


for some given $\gamma \neq 0$. Substituting (3.5.18) into (3.5.17), the $t$-ratio above can be rewritten as


\begin{equation*}
t_{\ell}=\frac{\sqrt{n}\left(\hat{\delta}_{\ell}(\widehat{\mathbf{W}})-\delta_{\ell}^{(n)}\right)}{\sqrt{(\operatorname{Avar}(\widehat{\boldsymbol{\delta}}(\widehat{\mathbf{W}})))_{\ell \ell}}}+\frac{\gamma}{\sqrt{(\operatorname{Avar}(\hat{\boldsymbol{\delta}}(\widehat{\mathbf{W}})))_{\ell \ell}}} \tag{3.5.19}
\end{equation*}


Applying the same argument for deriving the asymptotic distribution of the $t$-ratio in (2.4.4), we can show that $t_{\ell} \rightarrow_{\mathrm{d}} N(\mu, 1)$, where


\begin{equation*}
\mu \equiv \frac{\gamma}{\sqrt{(\operatorname{Avar}(\hat{\delta}(\widehat{\mathbf{W}})))_{\ell \ell}}} \tag{3.5.20}
\end{equation*}


If the significance level is $\alpha$, the asymptotic power is given by $\operatorname{Prob}\left(|x|>t_{\alpha / 2}\right)$, where $x \sim N(\mu, 1)$ and $t_{\alpha / 2}$ is the level $\alpha$ critical value. Evidently, the larger is $|\mu|$, the higher is the asymptotic power. And $|\mu|$ decreases with the asymptotic variance. Therefore, the asymptotic power against local alternatives is maximized by the efficient GMM estimator.

\section*{Small-Sample Properties}
Do these desirable asymptotic properties of the efficient GMM estimator and associated test statistics carry over to their finite-sample distributions? The efficient GMM estimator uses $\widehat{\mathbf{S}}^{-1}$, a function of estimated fourth moments, for $\widehat{\mathbf{W}}$. Generally, it takes a substantially larger sample size to estimate fourth moments reliably than to estimate first and second moments. So we would expect the efficient GMM estimator to have poorer small-sample properties than the GMM estimators that do not use fourth moments for $\widehat{\mathbf{W}}$. The July 1996 issue of the Journal of Business and Economic Statistics has a number of papers examining the small-sample distribution of GMM estimators and associated test statistics for various DGPs. Their overall conclusion is that the equally weighted GMM estimator with $\widehat{\mathbf{W}}=\mathbf{I}$ generally outperforms the efficient GMM in terms of the bias and variance in finite samples. They also find that the size of the efficient Wald statistic in small samples far exceeds the assumed significance level. That is, if $\alpha$ is the assumed significance level and $c_{\alpha}$ is the associated critical value so that $\operatorname{Prob}\left(\chi^{2}>c_{\alpha}\right)=\alpha$, the probability in finite samples that the Wald statistic is greater than $c_{\alpha}$ far exceeds $\alpha$; the test rejects the null too often. Unfortunately, however, like most other smallsample studies, those studies fail to produce clear quantitative guidance which the empirical researcher could follow.

\section*{QUESTIONS FOR REVIEW}
\begin{enumerate}
  \item Verify that all the results of Sections 2.3-2.5 are special cases of those of this section. In particular, verify that (3.5.1) reduces to (2.3.4). Hint: $\boldsymbol{\Sigma}_{\mathbf{x z}}$ is square if $\mathbf{z}_{i}=\mathbf{x}_{i}$.
  \item (Singular $\mathbf{W}$ ) Suppose $\mathbf{W}$ is singular but $\boldsymbol{\Sigma}_{\mathbf{x z}}^{\prime} \mathbf{W} \boldsymbol{\Sigma}_{\mathbf{x z}}$ is nonsingular. Verify that all the results of this section (except Proposition 3.5) remain valid.
  \item (Asymptotically equivalent choice of $\widehat{\mathbf{W}}$ ) Suppose $\widehat{\mathbf{W}}_{1}-\widehat{\mathbf{W}}_{2} \rightarrow_{\mathrm{p}} \mathbf{0}$. Show that
\end{enumerate}

$$
\sqrt{n} \hat{\boldsymbol{\delta}}\left(\widehat{\mathbf{W}}_{1}\right)-\sqrt{n} \hat{\boldsymbol{\delta}}\left(\widehat{\mathbf{W}}_{2}\right) \underset{\mathrm{p}}{\rightarrow} \mathbf{0}
$$

\section*{Hint:}
$$
\begin{aligned}
& \sqrt{n} \hat{\boldsymbol{\delta}}\left(\widehat{\mathbf{W}}_{1}\right)-\sqrt{n} \hat{\boldsymbol{\delta}}\left(\widehat{\mathbf{W}}_{2}\right) \\
& \quad=\left[\left(\mathbf{S}_{\mathbf{x z}}^{\prime} \widehat{\mathbf{W}}_{1} \mathbf{S}_{\mathbf{x z}}\right)^{-1} \mathbf{S}_{\mathbf{x z}}^{\prime} \widehat{\mathbf{W}}_{1}-\left(\mathbf{S}_{\mathbf{x z}}^{\prime} \widehat{\mathbf{W}}_{2} \mathbf{S}_{\mathbf{x z}}\right)^{-1} \mathbf{S}_{\mathbf{x z}}^{\prime} \widehat{\mathbf{W}}_{2}\right] \sqrt{n} \overline{\mathbf{g}}
\end{aligned}
$$

$\sqrt{n} \overline{\mathbf{g}}$ converges in distribution to a random variable. Use Lemma 2.4(b).\\
4. (Three-step GMM) Consider adding to the efficient two-step GMM procedure the following third step: Recompute $\widehat{\mathbf{S}}$ by (3.5.10), but this time using the residual from the second step. Calculate the GMM estimator with this recomputed $\widehat{\mathbf{S}}$. Is this estimator consistent? Asymptotically normal? Efficient? Hint: Verify by Proposition 3.4 that the recalculated $\widehat{\mathbf{S}}$ is consistent for $\mathbf{S}$. Does the asymptotic distribution of the GMM estimator depend on the choice of $\widehat{\mathbf{W}}$ if $\operatorname{plim}_{n \rightarrow \infty} \widehat{\mathbf{W}}$ is the same?\\
5. (When $\mathbf{z}_{i}$ is a strict subset of $\mathbf{x}_{i}$ ) Suppose $\mathbf{z}_{i}$ is a strict subset of $\mathbf{x}_{i}$. So $\mathbf{x}_{i}$ includes, in addition to the regressors (which are all predetermined), some other predetermined variables. Are the efficient two-step GMM estimator and the OLS estimator numerically the same? [Answer: No.]\\
6. (Data matrix representation of efficient GMM) Let $\mathbf{B}$ be the $n \times n$ diagonal matrix whose $i$-th element is $\hat{\varepsilon}_{i}^{2}$, where $\hat{\varepsilon}_{i}$ is the residual from the first-step consistent estimation. That is,

$$
\mathbf{B} \equiv\left[\begin{array}{ccc}
\hat{\varepsilon}_{1}^{2} & & \\
& \ddots & \\
& & \hat{\varepsilon}_{n}^{2}
\end{array}\right]
$$

Verify that

$$
\hat{\delta}\left(\widehat{\mathbf{S}}^{-1}\right)=\left[\mathbf{Z}^{\prime} \mathbf{X}\left(\mathbf{X}^{\prime} \mathbf{B X}\right)^{-1} \mathbf{X}^{\prime} \mathbf{Z}\right]^{-1} \mathbf{Z}^{\prime} \mathbf{X}\left(\mathbf{X}^{\prime} \mathbf{B X}\right)^{-1} \mathbf{X}^{\prime} \mathbf{y}
$$

where $\mathbf{X}, \mathbf{y}$, and $\mathbf{Z}$ are data matrices for the instruments, the dependent variable, and the regressors (they are defined in Section 3.8 below).\\
7. (GLS interpretation of efficient GMM) Let $\mathbf{X}, \mathbf{Z}$, and $\mathbf{y}$ be as in the previous question. Then the estimation equation can be written as $\mathbf{y}=\mathbf{Z} \boldsymbol{\delta}+\boldsymbol{\varepsilon}$. Premultiply both sides by $\mathbf{X}$ to obtain

$$
\mathbf{X}^{\prime} \mathbf{y}=\mathbf{X}^{\prime} \mathbf{Z} \boldsymbol{\delta}+\mathbf{X}^{\prime} \boldsymbol{\varepsilon}
$$

Taking $\mathbf{S}$ to be the variance matrix of $\mathbf{X}^{\prime} \boldsymbol{\varepsilon}$ and $\widehat{\mathbf{S}}$ to be its consistent estimate, apply FGLS. Verify that the FGLS estimator is the efficient GMM estimator (the FGLS estimator was defined in Section 1.6).\\
8. (Linear combination of orthogonality conditions) Derive the efficient GMM estimator that exploits a linear combination of orthogonality conditions,

$$
\mathbf{A} \mathrm{E}\left(\mathbf{g}_{i}\right)=\mathbf{0},
$$

where $\mathbf{A}$ is a $q \times K$ matrix of full row rank (so $q \leq K$ ). [Answer: Replace $\mathbf{S}_{\mathbf{x z}}$ by $\mathbf{A S}_{\mathbf{x y}}, \mathbf{s}_{\mathbf{x z}}$ by $\mathbf{A s}_{\mathbf{x y}}$, and $\mathbf{S}$ by $\mathbf{A S A}^{\prime}$. Formally, the estimator can be written as (3.4.8) with $\widehat{\mathbf{W}}=\mathbf{A}^{\prime}\left(\widehat{\mathbf{S}} \mathbf{A}^{\prime}\right)^{-1} \mathbf{A}$.] Verify: If $q=K$ (so that $\mathbf{A}$ is nonsingular), then the efficient GMM estimator is numerically equal to the efficient GMM estimator associated with the orthogonality conditions $\mathrm{E}\left(\mathbf{g}_{i}\right)=\mathbf{0}$.

\subsection*{3.6 Testing Overidentifying Restrictions}
If the equation is exactly identified, then it is possible to choose $\tilde{\delta}$ so that all the elements of the sample moments $\mathbf{g}_{n}(\tilde{\boldsymbol{\delta}})$ are zero and the distance

$$
J(\tilde{\delta}, \widehat{\mathbf{W}}) \equiv n \cdot \mathbf{g}_{n}(\tilde{\delta})^{\prime} \widehat{\mathbf{W}} \mathbf{g}_{n}(\tilde{\delta})
$$

is zero. (The $\tilde{\boldsymbol{\delta}}$ that does it is the IV estimator.) If the equation is overidentified, then the distance cannot be set to zero exactly, but we would expect the minimized distance to be close to zero. It turns out that, if the weighting matrix $\widehat{\mathbf{W}}$ is chosen optimally so that $\operatorname{plim} \widehat{\mathbf{W}}=\mathbf{S}^{-1}$, then the minimized distance is asymptotically chi-squared.

Let $\widehat{\mathbf{S}}$ be a consistent estimator of $\mathbf{S}$, and consider first the case where the distance is evaluated at the true parameter value $\boldsymbol{\delta}, J\left(\boldsymbol{\delta}, \widehat{\mathbf{S}}^{-1}\right)$. Since by definition $\mathbf{g}_{n}(\tilde{\boldsymbol{\delta}})=\overline{\mathbf{g}}\left(\equiv \frac{1}{n} \Sigma_{i} \mathbf{g}_{i}\right)$ for $\tilde{\boldsymbol{\delta}}=\boldsymbol{\delta}$, the distance equals


\begin{equation*}
J\left(\boldsymbol{\delta}, \widehat{\mathbf{S}}^{-1}\right)=n \cdot \overline{\mathbf{g}}^{\prime} \widehat{\mathbf{S}}^{-1} \overline{\mathbf{g}}=(\sqrt{n} \overline{\mathbf{g}})^{\prime} \widehat{\mathbf{S}}^{-1}(\sqrt{n} \overline{\mathbf{g}}) . \tag{3.6.1}
\end{equation*}


Since $\sqrt{n} \overline{\mathbf{g}} \rightarrow_{\mathrm{d}} N(\mathbf{0}, \mathbf{S})$ and $\widehat{\mathbf{S}} \rightarrow_{\mathrm{p}} \mathbf{S}$, its asymptotic distribution is $\chi^{2}(K)$ by Lemma 2.4(d). Now if $\boldsymbol{\delta}$ is replaced by $\hat{\boldsymbol{\delta}}\left(\widehat{\mathbf{S}}^{-1}\right)$, then the degrees of freedom change from $K$ to $K-L$. The intuitive reason is that we have to estimate $L$ parameters $\delta$ before forming the sample average of $\mathbf{g}_{i}$. (We encountered a similar situation in Chapter 1 in the context of the unbiased estimation of $\sigma^{2}$.) We summarize this as

Proposition 3.6 (Hansen's test of overidentifying restrictions (Hansen, 1982)): Suppose there is available a consistent estimator, $\widehat{\mathbf{S}}$, of $\mathbf{S}\left(=\mathrm{E}\left(\mathbf{g}_{i} \mathbf{g}_{i}^{\prime}\right)\right)$. Under Assumptions 3.1-3.5,

$$
J\left(\hat{\delta}\left(\widehat{\mathbf{S}}^{-1}\right), \widehat{\mathbf{S}}^{-1}\right)\left(=n \cdot \mathbf{g}_{n}\left(\hat{\delta}\left(\widehat{\mathbf{S}}^{-1}\right)\right)^{\prime} \widehat{\mathbf{S}}^{-1} \mathbf{g}_{n}\left(\hat{\delta}\left(\widehat{\mathbf{S}}^{-1}\right)\right)\right) \underset{\mathrm{d}}{\longrightarrow} \chi^{2}(K-L) .
$$

A formal proof is left as an optional exercise. Because the $\widehat{\mathbf{S}}$ given in (3.5.10) is consistent (under the additional condition of Assumption 3.6), the minimized distance calculated in the second step of the efficient two-step GMM is asymptotically $\chi^{2}(K-L)$.

Three points on the use and interpretation of the $J$ test are in order.

\begin{itemize}
  \item This is a specification test, testing whether all the restrictions of the model (which are the assumptions maintained in Proposition 3.6) are satisfied. If the $J$ statistic of Proposition 3.6 is surprisingly large, it means that either the orthogonality conditions (Assumption 3.3) or the other assumptions (or both) are likely to be false. Only when we are confident about those other assumptions can we interpret the large $J$ statistic as evidence for the endogeneity of some of the $K$ instruments included in $\mathbf{x}_{i}$.
  \item Unlike the tests we have encountered so far, the test is not consistent against some failures of the orthogonality conditions. The essential reason is the loss of degrees of freedom from $K$ to $K-L$. It is easy to show that $\overline{\mathbf{g}}$ is related to its sample counterpart, $\mathbf{g}_{n}\left(\hat{\boldsymbol{\delta}}\left(\widehat{\mathbf{S}}^{-1}\right)\right)$, as
\end{itemize}


\begin{equation*}
\sqrt{n} \mathbf{g}_{n}\left(\hat{\boldsymbol{\delta}}\left(\widehat{\mathbf{S}}^{-1}\right)\right)=\widehat{\mathbf{B}} \sqrt{n} \overline{\mathbf{g}}, \quad \widehat{\mathbf{B}} \equiv \mathbf{I}_{K}-\mathbf{S}_{\mathbf{x z}}\left(\mathbf{S}_{\mathbf{x z}}^{\prime} \widehat{\mathbf{S}}^{-1} \mathbf{S}_{\mathbf{x z}}\right)^{-1} \mathbf{S}_{\mathbf{x z}}^{\prime} \widehat{\mathbf{S}}^{-1} \tag{3.6.2}
\end{equation*}


The problem is that, since $\widehat{\mathbf{B}} \mathbf{S}_{\mathbf{x z}}=\mathbf{0}$, this matrix $\widehat{\mathbf{B}}$ is not of full column rank. If the orthogonality conditions fail and $\mathrm{E}\left(\mathbf{g}_{i}\right) \neq \mathbf{0}$, then the elements of $\sqrt{n} \overline{\mathbf{g}}$ will diverge to $+\infty$ or $-\infty$. But, since $\widehat{\mathbf{B}}$ is not of full column rank, $\widehat{\mathbf{B}} \sqrt{n} \overline{\mathbf{g}}$ and hence $J\left(\hat{\delta}\left(\widehat{\mathbf{S}}^{-1}\right), \widehat{\mathbf{S}}^{-1}\right)$ may remain finite for some pattern of nonzero orthogonalities. ${ }^{12}$

\begin{itemize}
  \item It is only recently that the small-sample properties of the test became a concern. Several papers in the July 1996 issue of the Journal of Business and Economic Statistics report that the finite-sample or actual size of the $J$ test in small samples far exceeds the nominal size (i.e., the test rejects too often).
\end{itemize}

\section*{Testing Subsets of Orthogonality Conditions}
Suppose we can divide the $K$ instruments into two groups: the vector $\mathbf{x}_{i 1}$ of $K_{1}$ variables that are known to satisfy the orthogonality conditions, and the vector $\mathbf{x}_{i 2}$ of remaining $K-K_{1}$ variables that are suspect. Since the ordering of instruments does not change the numerical values of the estimator and test statistics, we can assume without loss of generality that the last $K-K_{1}$ elements of $\mathbf{x}_{i}$ are the suspect

\footnotetext{${ }^{12}$ See Newey (1985, Section 3) for a thorough treatment of this issue.
}
instruments:
\[
\left.\mathbf{x}_{i}=\left[\begin{array}{l}
\mathbf{x}_{i 1}  \tag{3.6.3}\\
\mathbf{x}_{i 2}
\end{array}\right]\right\} \begin{aligned}
& K_{1} \text { rows } \\
& K-K_{1} \text { rows }
\end{aligned}
\]

The part of the model we wish to test is


\begin{equation*}
\mathrm{E}\left(\mathbf{x}_{i 2} \cdot \varepsilon_{i}\right)=\mathbf{0} \tag{3.6.4}
\end{equation*}


This restriction is testable if there are at least as many nonsuspect instruments as there are coefficients so that $K_{1} \geq L$. The basic idea is to compare two $J$ statistics from two separate GMM estimators of the same coefficient vector $\delta$, one using only the instruments included in $\mathbf{x}_{i 1}$, and the other using also the suspect instruments $\mathbf{x}_{i 2}$ in addition to $\mathbf{x}_{i 1}$. If the inclusion of the suspect instruments significantly increases the $J$ statistic, that is a good reason for doubting the predeterminedness of $\mathbf{x}_{i 2}$.

In accordance with the partition of $\mathbf{x}_{i}$, the sample orthogonality conditions $\mathbf{g}_{n}(\tilde{\boldsymbol{\delta}})$ and $\mathbf{S}$ can be written as

\[
\underset{(K \times 1)}{\mathbf{g}_{n}(\tilde{\delta})} \equiv\left[\begin{array}{c}
\mathbf{g}_{1 n}(\tilde{\delta})  \tag{3.6.5}\\
\left(K_{1} \times 1\right) \\
\mathbf{g}_{2 n}(\tilde{\delta}) \\
\left(\left(K-K_{1}\right) \times 1\right)
\end{array}\right], \underset{(K \times K)}{\mathbf{S}} \equiv\left[\begin{array}{ll}
\mathbf{S}_{11} & \mathbf{S}_{12} \\
\mathbf{S}_{21} & \mathbf{S}_{22}
\end{array}\right],
\]

where

$$
\mathbf{S}_{11}=\mathrm{E}\left(\varepsilon_{i}^{2} \mathbf{x}_{i 1} \mathbf{x}_{i 1}^{\prime}\right), \mathbf{S}_{12}=\mathrm{E}\left(\varepsilon_{i}^{2} \mathbf{x}_{i 1} \mathbf{x}_{i 2}^{\prime}\right), \mathbf{S}_{21}=\mathrm{E}\left(\varepsilon_{i}^{2} \mathbf{x}_{i 2} \mathbf{x}_{i 1}^{\prime}\right), \mathbf{S}_{22}=\mathrm{E}\left(\varepsilon_{i}^{2} \mathbf{x}_{i 2} \mathbf{x}_{i 2}^{\prime}\right) .
$$

In particular, $\mathbf{g}_{1 n}(\tilde{\delta})$ can be written as

$$
\mathbf{g}_{1 n}(\tilde{\boldsymbol{\delta}})=\frac{1}{n} \sum_{i=1}^{n} \mathbf{x}_{i 1} \cdot\left(y_{i}-\mathbf{z}_{i}^{\prime} \tilde{\boldsymbol{\delta}}\right) \equiv \mathbf{s}_{\mathbf{x}_{1} \mathbf{y}}-\mathbf{S}_{\mathbf{x}_{1} \mathbf{z}} \tilde{\boldsymbol{\delta}}
$$

where


\begin{equation*}
\mathbf{s}_{\mathbf{x}_{1} \mathbf{y}} \equiv \frac{1}{n} \sum_{i=1}^{n} \mathbf{x}_{i 1} \cdot y_{i}, \quad \mathbf{S}_{\mathbf{x}_{1} \mathbf{z}} \equiv \frac{1}{n} \sum_{i=1}^{n} \mathbf{x}_{i 1} \mathbf{z}_{i}^{\prime} \tag{3.6.6}
\end{equation*}


For a consistent estimate $\widehat{\mathbf{S}}$ of $\mathbf{S}$, the efficient GMM estimator using all the $K$ instruments and its associated $J$ statistic have already been derived in this and the\\
previous section. Reproducing them,


\begin{align*}
& \hat{\boldsymbol{\delta}}=\left(\mathbf{S}_{\mathbf{x z}}^{\prime} \widehat{\mathbf{S}}^{-1} \mathbf{S}_{\mathbf{x z}}\right)^{-1} \mathbf{S}_{\mathbf{x z}}^{\prime} \widehat{\mathbf{S}}^{-1} \mathbf{S}_{\mathbf{x y}}  \tag{3.6.7}\\
& J=n \cdot \mathbf{g}_{n}(\hat{\boldsymbol{\delta}})^{\prime} \widehat{\mathbf{S}}^{-1} \mathbf{g}_{n}(\hat{\boldsymbol{\delta}}) \tag{3.6.8}
\end{align*}


The efficient GMM estimator of the same coefficient vector $\boldsymbol{\delta}$ using only the first $K_{1}$ instruments and its associated $J$ statistic are obtained by replacing $\mathbf{x}_{i}$ by $\mathbf{x}_{i 1}$ in these expressions. So


\begin{align*}
\overline{\boldsymbol{\delta}} & =\left[\mathbf{S}_{\mathbf{x}_{1} \mathbf{z}}^{\prime}\left(\widehat{\mathbf{S}}_{11}\right)^{-1} \mathbf{S}_{\mathbf{x}_{1} \mathbf{z}}\right]^{-1} \mathbf{S}_{\mathbf{x}_{1} \mathbf{z}}^{\prime}\left(\widehat{\mathbf{S}}_{11}\right)^{-1} \mathbf{S}_{\mathbf{x}_{1} \mathbf{y}}  \tag{3.6.9}\\
J_{1} & =n \cdot \mathbf{g}_{1 n}(\overline{\boldsymbol{\delta}})^{\prime}\left(\widehat{\mathbf{S}}_{11}\right)^{-1} \mathbf{g}_{1 n}(\overline{\boldsymbol{\delta}}) \tag{3.6.10}
\end{align*}


where $\widehat{\mathbf{S}}_{11}$ is a consistent estimate of $\mathbf{S}_{11}$.\\
The test is based on the following proposition specifying the asymptotic distribution of $J-J_{1}$ (the proof is left as an optional exercise).

Proposition 3.7 (testing a subset of orthogonality conditions ${ }^{13}$ ): Suppose Assumptions 3.1-3.5 hold. Let $\mathbf{x}_{i 1}$ be a subvector of $\mathbf{x}_{i}$, and strengthen Assumption 3.4 by requiring that the rank condition for identification is satisfied for $\mathbf{x}_{i 1}$ (so $\mathrm{E}\left(\mathbf{x}_{i 1} \mathbf{z}_{i}^{\prime}\right)$ is of full column rank). Then, for any consistent estimators $\widehat{\mathbf{S}}$ of $\mathbf{S}$ and $\widehat{\mathbf{S}}_{11}$ of $\mathbf{S}_{11}$,

$$
C \equiv J-J_{1} \underset{\mathrm{~d}}{\rightarrow} \chi^{2}\left(K-K_{1}\right)
$$

where $K=\# \mathbf{x}_{i}$ (dimension of $\mathbf{x}_{i}$ ), $K_{1}=\# \mathbf{x}_{i 1}$ (dimension of $\mathbf{x}_{i 1}$ ), and $J$ and $J_{1}$ are defined in (3.6.8) and (3.6.10).

Clearly, the choice of $\widehat{\mathbf{S}}$ and $\widehat{\mathbf{S}}_{11}$ does not matter asymptotically as long as they are consistent. But in finite samples, the test statistic $C$ can be negative. This problem can be avoided and $C$ can be made nonnegative in finite samples if the same $\widehat{\mathbf{S}}$ is used throughout, that is, if $\widehat{\mathbf{S}}_{11}$ in (3.6.9) and (3.6.10) is the submatrix of $\widehat{\mathbf{S}}$ in (3.6.7) and (3.6.8). This is accomplished by taking the following steps:\\
(1) do the efficient two-step GMM with full instruments $\mathbf{x}_{i}$ to obtain $\widehat{\mathbf{S}}$ from the first step, $\hat{\delta}$ and $J$ from the second step;\\
(2) extract the submatrix $\widehat{\mathbf{S}}_{11}$ from $\widehat{\mathbf{S}}$ obtained from (1), calculate $\overline{\boldsymbol{\delta}}$ by (3.6.9) using this $\widehat{\mathbf{S}}_{11}$ and $J_{1}$ by (3.6.10). Then take the difference in $J$.

\footnotetext{${ }^{13}$ The test was developed in Newey (1985, Section 4) and Eichenbaum, Hansen, and Singleton (1985, Appendix C).
}It is left as an optional analytical exercise to prove that $C$ calculated as just described is nonnegative in finite samples.

We can use Proposition 3.7 to test for the endogeneity of a subset of regressors, as the next example illustrates.

Example 3.3 (testing whether schooling is predetermined in the wage equation): In the wage equation of Example 3.2, suppose schooling $S_{i}$ is suspected to be endogenous. To test for the endogeneity of $S_{i}$, partition $\mathbf{x}_{i}$ as

$$
\mathbf{x}_{i 1}=\left[\begin{array}{c}
1 \\
E X P R_{i} \\
A G E_{i} \\
M E D_{i}
\end{array}\right], \quad \mathbf{x}_{i 2}=S_{i} .
$$

The vector of regressors, $\mathbf{z}_{i}$, is the same as in Example 3.2. The first step is to do the efficient two-step GMM estimation of $\boldsymbol{\delta}$ with $\mathbf{x}_{i} \equiv\left(1, E X P R_{i}, A G E_{i}\right.$, $\left.M E D_{i}, S_{i}\right)^{\prime}$ as the instruments. This produces $J$ and the $5 \times 5$ matrix $\widehat{\mathbf{S}}$. Second, extract the leading $4 \times 4$ submatrix $\widehat{\mathbf{S}}_{11}$ corresponding to $\mathbf{x}_{i 1}$ from $\widehat{\mathbf{S}}$ and estimate the same coefficients $\boldsymbol{\delta}$ by GMM, this time with the fewer instruments $\mathbf{x}_{i 1}$ and using this $\widehat{\mathbf{S}}_{11}$. The difference in the $J$ statistic from the two different GMM estimators of $\delta$ should be asymptotically $\chi^{2}(1)$.

\section*{QUESTIONS FOR REVIEW}
\begin{enumerate}
  \item Does $J\left(\hat{\boldsymbol{\delta}}\left(\widehat{\mathbf{S}}^{-1}\right), \overline{\mathbf{S}}^{-1}\right) \rightarrow_{\mathrm{d}} \chi^{2}(K-L)$ where $\widehat{\mathbf{S}}$ and $\overline{\mathbf{S}}$ are two different consistent estimators of $\mathbf{S}$ ?
  \item Show: $J\left(\hat{\boldsymbol{\delta}}\left(\widehat{\mathbf{S}}^{-1}\right), \widehat{\mathbf{S}}^{-1}\right)=n \cdot \mathbf{S}_{\mathbf{x y}}^{\prime} \widehat{\mathbf{S}}^{-1}\left(\mathbf{S}_{\mathbf{x y}}-\mathbf{S}_{\mathbf{x z}} \hat{\boldsymbol{\delta}}\left(\widehat{\mathbf{S}}^{-1}\right)\right)$.
  \item Can the degrees of freedom of the $C$ statistic be greater than $K-L$ ? [Answer: No.]
  \item Suppose $K_{1}=L$. Does the numerical value of $C$ depend on the partition of $\mathbf{x}_{i}$ between $\mathbf{x}_{i 1}$ and $\mathbf{x}_{i 2}$ ?
\end{enumerate}

\subsection*{3.7 Hypothesis Testing by the Likelihood-Ratio Principle}
We have derived in Section 3.5 the chi-squared test statistics for the null hypothesis $\mathrm{H}_{0}: \mathbf{a}(\boldsymbol{\delta})=\mathbf{0}$ by the Wald principle. This section does the same by the likelihoodratio (LR) principle, which is to examine the difference in the objective function with and without the imposition of the null. Derivation of test statistics by the Lagrange Multiplier principle for GMM and extension to nonlinear equations are given in Section 7.4.

In the efficient GMM estimation, the objective function is $J\left(\tilde{\delta}, \widehat{\mathbf{S}}^{-1}\right)$ for a given consistent estimate $\widehat{\mathbf{S}}$ of $\mathbf{S}$. The restricted efficient GMM estimator is defined as


\begin{equation*}
\text { restricted efficient GMM: } \bar{\delta}\left(\widehat{\mathbf{S}}^{-1}\right) \equiv \underset{\tilde{\delta}}{\operatorname{argmin}} J\left(\tilde{\delta}, \widehat{\mathbf{S}}^{-1}\right) \text { subject to } \mathrm{H}_{0} \tag{3.7.1}
\end{equation*}


The LR principle suggests that


\begin{equation*}
L R \equiv J\left(\bar{\delta}\left(\widehat{\mathbf{S}}^{-1}\right), \widehat{\mathbf{S}}^{-1}\right)-J\left(\hat{\delta}\left(\widehat{\mathbf{S}}^{-1}\right), \widehat{\mathbf{S}}^{-1}\right) \tag{3.7.2}
\end{equation*}


should be asymptotically chi-squared. Indeed it is.

Proposition 3.8 (test statistic by the LR principle): Suppose Assumptions 3.13.5 hold and suppose there is available a consistent estimator, $\widehat{\mathbf{S}}$, of $\mathbf{S}\left(=\mathrm{E}\left(\mathbf{g}_{i} \mathbf{g}_{i}^{\prime}\right)\right)$. Consider the null hypothesis of \#a restrictions $\mathrm{H}_{0}: \mathbf{a}(\boldsymbol{\delta})=\mathbf{0}$ such that $\mathbf{A}(\boldsymbol{\delta})$, the \#a × $L$ matrix of first derivatives, is continuous and of full row rank. Define two statistics, $W$ and $L R$, by (3.5.16) and (3.7.2), respectively. Then, under the null, the following holds:\\
(a) The two statistics are asymptotically equivalent in that their asymptotic distributions are the same (namely, $\left.\chi^{2}(\# \mathbf{a})\right)$.\\
(b) The two statistics are asymptotically equivalent in the stronger sense that their numerical difference converges in probability to zero: $L R-W \rightarrow_{\mathrm{p}} 0$. (By Lemma 2.4(a), this result is stronger than (a).)\\
(c) Furthermore, if the hypothesis is linear so that the restrictions can be written as $\mathbf{R} \boldsymbol{\delta}=\mathbf{r}$, then the two statistics are numerically equal.

Proving this for the linear case, which is an algebraic result, is left as an analytical exercise. For a full proof, see Section 7.4.

Several comments about Proposition 3.8:

\begin{itemize}
  \item The advantage of $L R$ over $W$ is invariance: the numerical value of $L R$ does not depend on how the nonlinear restrictions are represented by $\mathbf{a}(\cdot)$. On the other hand, you have to write a nonlinear optimization computer program to find the restricted efficient GMM when the hypothesis is nonlinear.
  \item Proposition 3.8 requires that the distance matrix $\widehat{\mathbf{W}}$ satisfy the efficiency condition $\operatorname{plim} \widehat{\mathbf{W}}=\mathbf{S}^{-1}$. Otherwise $L R$ is not asymptotically chi-squared. In contrast, the Wald statistic is asymptotically chi-squared without $\widehat{\mathbf{W}}$ satisfying the efficiency condition.
  \item The same estimate of $\mathbf{S}$ should be used throughout in the calculation of $L R$. Let $\widetilde{\mathbf{S}}$ and $\overline{\mathbf{S}}$ be two different consistent estimators of $\mathbf{S}$, and consider the statistic
\end{itemize}

$$
J\left(\bar{\delta}\left(\overline{\mathbf{S}}^{-1}\right), \overline{\mathbf{S}}^{-1}\right)-J\left(\hat{\delta}\left(\widetilde{\mathbf{S}}^{-1}\right), \widetilde{\mathbf{S}}^{-1}\right) .
$$

This statistic is what you end up with if you perform two separate two-step efficient GMMs with and without the constraint of the null; $\overline{\mathbf{S}}$ is the estimate of $\mathbf{S}$ from the first step with the constraint, while $\widetilde{\mathbf{S}}$ is from the first step without the constraint. The statistic is asymptotically equivalent to $L R$ (in that the difference between the two converges in probability to zero), but in finite samples it may be negative. Having the same estimate of $\mathbf{S}$ throughout ensures the nonnegativity of the statistic in finite samples. Researchers usually use the estimate from unconstrained estimation ( $\widetilde{\mathbf{S}}$ ) here, but the estimate from constrained estimation is also valid because it is consistent for $\mathbf{S}$ under the null.

\begin{itemize}
  \item Part (b) of the proposition (the asymptotic equivalence in a stronger sense) means that if the sample size is large enough and the hypothesis is true, then the outcome (not just the probability of rejection or acceptance) of the test based on $L R$ will be the same as that based on $W$ because the probability that the two statistics differ numerically by an even tiny amount is zero in sufficiently large samples.
  \item For the numerical equivalence $L R=W$ for the linear case to hold, the same $\widehat{\mathbf{S}}$ must be used throughout to calculate not only $L R$ but also the Wald statistic. Otherwise $L R$ and $W$ are only asymptotically equivalent.
\end{itemize}

\section*{The $L R$ Statistic for the Regression Model}
Because the regression model of Chapter 2 is a special case of the GMM model of this chapter, it may be useful to know how $L R$ would look for that special case.

Because $\mathbf{x}_{i}=\mathbf{z}_{i}$ in the regression model, the (unrestricted) efficient GMM estimator is OLS and $J\left(\hat{\boldsymbol{\delta}}\left(\widehat{\mathbf{S}}^{-1}\right), \widehat{\mathbf{S}}^{-1}\right)=0$. Therefore,


\begin{equation*}
L R=J\left(\bar{\delta}\left(\widehat{\mathbf{S}}^{-1}\right), \widehat{\mathbf{S}}^{-1}\right) \tag{3.7.3}
\end{equation*}


where $\bar{\delta}\left(\widehat{\mathbf{S}}^{-1}\right)$ is the restricted efficient GMM estimator. By Proposition 3.8, this statistic is asymptotically chi-squared and is numerically equal to the Wald statistic if the null is linear. As will be shown below, under conditional homoskedasticity, this statistic can be written as the difference in the sum of squared residuals normalized to the error variance.

\section*{Variable Addition Test (optional)}
In the previous section, we considered a specification test based on the $C$ statistic for the endogeneity of a subset, $\mathbf{x}_{i 2}$ of instruments $\mathbf{x}_{i}$ while assuming that the other instruments $\mathbf{x}_{i 1}$ are predetermined. Occasionally, we encounter a special case where $\mathbf{z}_{i}=\mathbf{x}_{i 1}$ :


\begin{equation*}
y_{i}=\mathbf{x}_{i 1}^{\prime} \boldsymbol{\delta}+\varepsilon_{i} . \tag{3.7.4}
\end{equation*}


A popular method to test whether the suspect instruments $\mathbf{x}_{i 2}$ are predetermined is to estimate the augmented equation

\[
y_{i}=\mathbf{x}_{i 1}^{\prime} \boldsymbol{\delta}+\mathbf{x}_{i 2}^{\prime} \boldsymbol{\alpha}+\varepsilon_{i}=\mathbf{x}_{i}^{\prime} \boldsymbol{\gamma}+\varepsilon_{i} \quad \text { with } \boldsymbol{\gamma}=\left[\begin{array}{l}
\boldsymbol{\delta}  \tag{3.7.5}\\
\boldsymbol{\alpha}
\end{array}\right]
\]

and test the null $\mathrm{H}_{0}: \boldsymbol{\alpha}=\mathbf{0}$. Testing is either by the Wald statistic $W$ or by the $L R$ statistic, which is numerically equal to $W$. This test is sometimes called the variable addition test. How is the test related to the $C$ test of Proposition 3.7?

To calculate $L R$, we have to calculate two efficient GMM estimators of $\boldsymbol{\gamma}$ with the same instrument set $\mathbf{x}_{i}$ : one with the constraint $\boldsymbol{\alpha}=\mathbf{0}$ and one without. The unrestricted efficient GMM estimator is the OLS estimator on the unconstrained equation (3.7.5). The associated $J$ statistic is zero. Let


\begin{equation*}
\widehat{\mathbf{S}} \equiv \frac{1}{n} \sum_{i=1}^{n} e_{i}^{2} \mathbf{x}_{i} \mathbf{x}_{i}^{\prime} \tag{3.7.6}
\end{equation*}


where $e_{i}$ is the OLS residual from the unrestricted regression (3.7.5). If we use this estimate of $\mathbf{S}$, the restricted efficient GMM estimator of $\boldsymbol{\gamma}$ minimizes the $J$ function

$$
J\left(\tilde{\boldsymbol{\gamma}}, \widehat{\mathbf{S}}^{-1}\right)=n \cdot\left(\mathbf{s}_{\mathbf{x y}}-\mathbf{S}_{\mathbf{x x}} \tilde{\boldsymbol{\gamma}}\right)^{\prime} \widehat{\mathbf{S}}^{-1}\left(\mathbf{s}_{\mathbf{x y}}-\mathbf{S}_{\mathbf{x x}} \tilde{\boldsymbol{\gamma}}\right)
$$

subject to $\boldsymbol{\alpha}=\mathbf{0}$. Clearly, the restricted estimator can be written as

$$
\overline{\boldsymbol{\gamma}}=\left[\begin{array}{c}
\hat{\boldsymbol{\delta}}\left(\widehat{\mathbf{S}}^{-1}\right) \\
0
\end{array}\right]
$$

where $\hat{\delta}\left(\widehat{\mathbf{S}}^{-1}\right)$ is the efficient GMM estimator with $\mathbf{x}_{i}$ as instruments on the restricted equation (3.7.4). So

$$
\begin{aligned}
L R & =n \cdot\left(\mathbf{s}_{\mathrm{xy}}-\mathbf{S}_{\mathrm{xx}} \overline{\boldsymbol{\gamma}}\right)^{\prime} \widehat{\mathbf{S}}^{-1}\left(\mathbf{s}_{\mathrm{xy}}-\mathbf{S}_{\mathrm{xx}} \overline{\boldsymbol{\gamma}}\right) \quad(\text { since } J=0 \text { for unrestricted GMM }) \\
& =n \cdot\left(\mathbf{s}_{\mathrm{xy}}-\mathbf{S}_{\mathrm{xx} 1} \hat{\delta}\left(\widehat{\mathbf{S}}^{-1}\right)\right)^{\prime} \widehat{\mathbf{S}}^{-1}\left(\mathbf{s}_{\mathrm{xy}}-\mathbf{S}_{\mathrm{xx}_{1}} \hat{\delta}\left(\widehat{\mathbf{S}}^{-1}\right)\right) \quad\left(\text { since } \mathbf{S}_{\mathrm{xx}} \overline{\boldsymbol{\gamma}}=\mathbf{S}_{\mathrm{xx}_{1}} \hat{\delta}\left(\widehat{\mathbf{S}}^{-1}\right)\right)
\end{aligned}
$$

This is none other than Hansen's $J$ statistic for the restricted equation (3.7.4) when $\mathbf{x}_{i}$ is the instrument vector. This statistic, in turn, equals the $C$ statistic because the $J_{1}$ in Proposition 3.7 is zero in the present case. Therefore, all three statistics, $W$ from the unrestricted regression, $L R$, and $C$, are numerically equal to Hansen's $J$, provided that the same $\widehat{\mathbf{S}}$ is used throughout. That is, the variable addition test is numerically equivalent to Hansen's test for overidentifying restrictions.

\section*{QUESTIONS FOR REVIEW}
\begin{enumerate}
  \item (LR for the regression model) Verify that, for the regression model where $\mathbf{z}_{i}=\mathbf{x}_{i}$,
\end{enumerate}

$$
L R=\mathbf{y}^{\prime} \mathbf{X}(n \widehat{\mathbf{S}})^{-1} \mathbf{X}^{\prime} \mathbf{y}-2 \mathbf{y}^{\prime} \mathbf{X}(n \widehat{\mathbf{S}})^{-1}\left(\mathbf{X}^{\prime} \mathbf{X}\right) \bar{\delta}+\bar{\delta}^{\prime}\left(\mathbf{X}^{\prime} \mathbf{X}\right)(n \widehat{\mathbf{S}})^{-1}\left(\mathbf{X}^{\prime} \mathbf{X}\right) \bar{\delta}
$$

\begin{enumerate}
  \setcounter{enumi}{1}
  \item (Choice of $\widehat{\mathbf{S}}$ in variable addition test) Suppose you form $\widehat{\mathbf{S}}$ from the residual from the restricted regression (3.7.4) and use it to form $W, L R$, and $C$. Are they numerically equal? Are they asymptotically chi-squared? Hint: Is this $\widehat{\mathbf{S}}$ consistent under the null?
\end{enumerate}

\subsection*{3.8 Implications of Conditional Homoskedasticity}
So far, we have not assumed conditional homoskedasticity in developing the asymptotics of the GMM estimator. This section considers the implication of imposing

\section*{Assumption 3.7 (conditional homoskedasticity):}
$$
\mathrm{E}\left(\varepsilon_{i}^{2} \mid \mathbf{x}_{i}\right)=\sigma^{2}
$$

Under conditional homoskedasticity, the matrix of fourth moments $\mathbf{S}$ (= $\left.\mathrm{E}\left(\mathbf{g}_{i} \mathbf{g}_{i}^{\prime}\right)=\mathrm{E}\left(\varepsilon_{i}^{2} \mathbf{x}_{i} \mathbf{x}_{i}^{\prime}\right)\right)$ can be written as a product of second moments:


\begin{equation*}
\mathbf{S}=\sigma^{2} \mathbf{\Sigma}_{\mathbf{x x}}, \tag{3.8.1}
\end{equation*}


where $\mathbf{\Sigma}_{\mathbf{x x}}=\mathrm{E}\left(\mathbf{x}_{i} \mathbf{x}_{i}^{\prime}\right)$. As in Chapter 2, this decomposition of $\mathbf{S}$ has several implications.

\begin{itemize}
  \item Since $\mathbf{S}$ is nonsingular by Assumption 3.5, this decomposition implies that $\sigma^{2}>$ 0 and $\boldsymbol{\Sigma}_{\mathbf{x x}}$ is nonsingular.
  \item The estimator exploiting this structure of $\mathbf{S}$ is
\end{itemize}


\begin{equation*}
\widehat{\mathbf{S}}=\hat{\sigma}^{2} \frac{1}{n} \sum_{i=1}^{n} \mathbf{x}_{i} \mathbf{x}_{i}^{\prime}=\hat{\sigma}^{2} \mathbf{S}_{\mathbf{x x}}, \tag{3.8.2}
\end{equation*}


where $\hat{\sigma}^{2}$ is some consistent estimator to be specified below. By ergodic stationarity, $\mathbf{S}_{\mathbf{x x}} \rightarrow_{\text {a.s. }} \boldsymbol{\Sigma}_{\mathbf{x x}}$. Thus, provided that $\hat{\sigma}^{2}$ is consistent, we do not need the fourth-moment assumption (Assumption 3.6) for $\widehat{\mathbf{S}}$ to be consistent.

Needless to say, all the results presented so far are valid under the extra condition of conditional homoskedasticity. But many of the results and formulas can be simplified substantially under this extra condition, by just replacing $\mathbf{S}$ by $\sigma^{2} \boldsymbol{\Sigma}_{\mathbf{x x}}$ and the expression (3.5.10) for $\widehat{\mathbf{S}}$ by (3.8.2). This section collects those simplifications.

\section*{Efficient GMM Becomes 2SLS}
In the efficient GMM estimation, the weighting matrix is $\widehat{\mathbf{S}}^{-1}$. If we set $\widehat{\mathbf{S}}$ to (3.8.2), the GMM estimator becomes


\begin{align*}
\hat{\delta}\left(\widehat{\mathbf{S}}^{-1}\right) & =\left[\mathbf{S}_{\mathbf{x z}}^{\prime}\left(\hat{\sigma}^{2} \mathbf{S}_{\mathbf{x x}}\right)^{-1} \mathbf{S}_{\mathbf{x z}}\right]^{-1} \mathbf{S}_{\mathbf{x z}}^{\prime}\left(\hat{\sigma}^{2} \mathbf{S}_{\mathbf{x x}}\right)^{-1} \mathbf{S}_{\mathbf{x y}} \\
& =\left(\mathbf{S}_{\mathbf{x z}}^{\prime} \mathbf{S}_{\mathbf{x x}}^{-1} \mathbf{S}_{\mathbf{x z}}\right)^{-1} \mathbf{S}_{\mathbf{x z}}^{\prime} \mathbf{S}_{\mathbf{x x}}^{-1} \mathbf{S}_{\mathbf{x y}} \\
& =\hat{\delta}\left(\mathbf{S}_{\mathbf{x x}}^{-1}\right) \equiv \hat{\boldsymbol{\delta}}_{2 \text { SLS }}, \tag{3.8.3}
\end{align*}


which does not depend on $\hat{\sigma}^{2}$. In the general case, the whole point of the first step in the efficient two-step GMM was to obtain a consistent estimator of $\mathbf{S}$. Under conditional homoskedasticity, there is no need to do the first step because the second-\\
step estimator collapses to the GMM estimator with $\mathbf{S}_{\mathbf{x x}}$ used for $\widehat{\mathbf{S}}, \hat{\boldsymbol{\delta}}\left(\mathbf{S}_{\mathbf{x x}}^{-1}\right)$. This estimator, $\hat{\boldsymbol{\delta}}_{\text {2SLS }}$, is called the Two-Stage Least Squares (2SLS or TSLS). ${ }^{14}$ The same equation can be estimated by maximum likelihood. Section 8.6 will cover the ML counterpart of 2SLS, called the "limited-information maximum likelihood estimator."

The expression for $\operatorname{Avar}\left(\hat{\delta}_{2 \text { SLS }}\right)$ can be obtained by substituting (3.8.1) into (3.5.13):


\begin{equation*}
\operatorname{Avar}\left(\hat{\boldsymbol{\delta}}_{2 \mathrm{SLS}}\right)=\sigma^{2} \cdot\left(\boldsymbol{\Sigma}_{\mathbf{x z}}^{\prime} \boldsymbol{\Sigma}_{\mathbf{x x}}^{-1} \boldsymbol{\Sigma}_{\mathbf{x z}}\right)^{-1} \tag{3.8.4}
\end{equation*}


A natural estimator of this is


\begin{equation*}
\widehat{\operatorname{Avar}\left(\hat{\delta}_{2 S L S}\right)}=\hat{\sigma}^{2} \cdot\left(\mathbf{S}_{\mathbf{x z}}^{\prime} \mathbf{S}_{\mathbf{x x}}^{-1} \mathbf{S}_{\mathbf{x z}}\right)^{-1} . \tag{3.8.5}
\end{equation*}


For $\hat{\sigma}^{2}$, consider the sample variance of the 2SLS residuals:


\begin{equation*}
\hat{\sigma}^{2} \equiv \frac{1}{n} \sum_{i=1}^{n}\left(y_{i}-\mathbf{z}_{i}^{\prime} \hat{\boldsymbol{\delta}}_{2 \mathrm{SLS}}\right)^{2} \tag{3.8.6}
\end{equation*}


(Some authors divide the sum of squares by $n-L$, not by $n$, to calculate $\hat{\sigma}^{2}$.) By Proposition 3.2, (3.8.6) $\rightarrow_{\mathrm{p}} \sigma^{2}$ if $\mathrm{E}\left(\mathbf{z}_{i} \mathbf{z}_{i}^{\prime}\right)$ exists and is finite. Thus $\widehat{\mathbf{S}}$ defined in (3.8.2) with this $\hat{\sigma}^{2}$ is consistent for $\mathbf{S}$.

Substituting (3.8.2) into (3.5.15) and (3.5.16), the $t$-ratio and the Wald statistic become


\begin{align*}
t_{\ell} & =\frac{\hat{\delta}_{2 S L S, \ell}-\bar{\delta}_{\ell}}{S E_{\ell}} \quad \text { with } \quad S E_{\ell}=\sqrt{\frac{\hat{\sigma}^{2}}{n} \cdot\left(\left(\mathbf{S}_{\mathbf{x z}}^{\prime} \mathbf{S}_{\mathbf{x x}}^{-1} \mathbf{S}_{\mathbf{x z}}\right)^{-1}\right)_{\ell \ell}}  \tag{3.8.7}\\
W & =n \cdot \frac{\mathbf{a}\left(\hat{\delta}_{2 S L S}\right)^{\prime}\left[\mathbf{A}\left(\hat{\delta}_{2 S L S}\right)\left(\mathbf{S}_{\mathbf{x z}}^{\prime} \mathbf{S}_{\mathbf{x x}}^{-1} \mathbf{S}_{\mathbf{x z}}\right)^{-1} \mathbf{A}\left(\hat{\delta}_{2 S L S}\right)^{\prime}\right]^{-1} \mathbf{a}\left(\hat{\delta}_{2 S L S}\right)}{\hat{\sigma}^{2}} . \tag{3.8.8}
\end{align*}


\section*{$\boldsymbol{J}$ Becomes Sargan's Statistic}
When $\widehat{\mathbf{W}}$ is set to $\left(\hat{\sigma}^{2} \mathbf{S}_{\mathbf{x x}}\right)^{-1}$, the distance defined in (3.4.6) becomes


\begin{equation*}
J\left(\tilde{\delta},\left(\hat{\sigma}^{2} \cdot \mathbf{S}_{\mathbf{x x}}\right)^{-1}\right)=n \cdot \frac{\left(\mathbf{S}_{\mathbf{x y}}-\mathbf{S}_{\mathbf{x z}} \tilde{\boldsymbol{\delta}}\right)^{\prime} \mathbf{S}_{\mathbf{x x}}^{-1}\left(\mathbf{S}_{\mathbf{x y}}-\mathbf{S}_{\mathbf{x z}} \tilde{\boldsymbol{\delta}}\right)}{\hat{\sigma}^{2}} \tag{3.8.9}
\end{equation*}


\footnotetext{${ }^{14}$ This estimator was first proposed by Theil (1953).
}Proposition 3.6 then implies that the distance evaluated at the efficient GMM estimator under conditional homoskedasticity, $\hat{\boldsymbol{\delta}}_{\text {2SLS }}$, is asymptotically chi-squared. This distance is called Sargan's statistic (Sargan, 1958):


\begin{equation*}
\text { Sargan's statistic }=n \cdot \frac{\left(\mathbf{s}_{\mathbf{x y}}-\mathbf{S}_{\mathbf{x z}} \hat{\boldsymbol{\delta}}_{2 S L S}\right)^{\prime} \mathbf{S}_{\mathbf{x x}}^{-1}\left(\mathbf{s}_{\mathbf{x y}}-\mathbf{S}_{\mathbf{x z}} \hat{\boldsymbol{\delta}}_{2 S L S}\right)}{\hat{\sigma}^{2}} \text {. } \tag{3.8.10}
\end{equation*}


We summarize our results so far as

\section*{Proposition 3.9 (asymptotic properties of 2SLS):}
(a) Under Assumptions 3.1-3.4, the 2SLS estimator (3.8.3) is consistent. If Assumption 3.5 is added, the estimator is asymptotically normal with the asymptotic variance given by (3.5.1) with $\mathbf{W}=\left(\sigma^{2} \mathbf{\Sigma}_{\mathbf{x x}}\right)^{-1}$. If Assumption 3.7 (conditional homoskedasticity) is added to Assumptions 3.1-3.5, then the estimator is the efficient GMM estimator.\\
Furthermore, if $\mathrm{E}\left(\mathbf{z}_{i} \mathbf{z}_{i}^{\prime}\right)$ exists and is finite, ${ }^{15}$ then\\
(b) the asymptotic variance is consistently estimated by (3.8.5),\\
(c) $t_{\ell}$ in (3.8.7) $\rightarrow_{\mathrm{d}} N(0,1), W$ in (3.8.8) $\rightarrow_{\mathrm{d}} \chi^{2}(\# \mathrm{r})$, and\\
(d) the Sargan statistic in $(3.8 .10) \rightarrow_{\mathrm{d}} \chi^{2}(K-L)$.

Proposition 3.8 states that the $L R$ statistic, which is the difference in $J$ with and without the imposition of the null hypothesis, is asymptotically chi-squared. Since $J$ can be written as (3.8.9), we have


\begin{equation*}
L R=n \cdot \frac{\left(\mathbf{S}_{\mathbf{x y}}-\mathbf{S}_{\mathbf{x z}} \overline{\boldsymbol{\delta}}\right)^{\prime} \mathbf{S}_{\mathbf{x x}}^{-1}\left(\mathbf{S}_{\mathbf{x y}}-\mathbf{S}_{\mathbf{x z}} \overline{\boldsymbol{\delta}}\right)-\left(\mathbf{S}_{\mathbf{x y}}-\mathbf{S}_{\mathbf{x z}} \hat{\boldsymbol{\delta}}_{2 \mathrm{SLS}}\right)^{\prime} \mathbf{S}_{\mathbf{x x}}^{-1}\left(\mathbf{S}_{\mathbf{x y}}-\mathbf{S}_{\mathbf{x z}} \hat{\boldsymbol{\delta}}_{2 \mathrm{SLS}}\right)}{\hat{\sigma}^{2}}, \tag{3.8.11}
\end{equation*}


where $\overline{\boldsymbol{\delta}}$ is the restricted 2SLS estimator which minimizes (3.8.9) under the null hypothesis. ${ }^{16}$ In Proposition 3.8, the use of the same $\widehat{\mathbf{S}}$ guaranteed the statistic to be nonnegative in finite samples. Here, deflation by the same $\hat{\sigma}^{2}$ ensures the statistic to be nonnegative. If the hypothesis is linear, then this $L R$ is numerically equal to the Wald statistic $W$.

\footnotetext{${ }^{15}$ This additional assumption is needed for the consistency of $\hat{\sigma}^{2}$; see Proposition 3.2.\\
${ }^{16}$ This statistic was first derived by Gallant and Jorgenson (1979).
}\section*{Small-Sample Properties of 2SLS}
There is a fairly large literature on the finite-sample distribution of the 2SLS estimator (see, e.g., Judge et al. (1985, Section 15.4) and Staiger and Stock (1997, Section 1)). Some studies manage to derive analytically the exact finite-sample distribution of the estimator, while others are Monte Carlo studies for various DGPs. The analytical results, however, are not useful for empirical researchers, because they are derived under the restrictive assumptions of fixed instruments and normal errors and the analytical expressions for distributions are computationally intractable.

For the case of a single regressor and a single (stochastic) instrument with normal errors, Nelson and Startz (1990) derive the exact finite-sample distribution of the 2SLS estimator that is fairly simple and easy to calculate. They also show that, when the instrument is "weak" in the sense of low explanatory power in the first-stage regression of the regressor on the instrument, a mass of the distribution of the sampling error $\hat{\delta}_{2 \text { SLS }}-\delta$ remains apart from zero until the sample size gets really large.

Their work illustrates the need for reporting the $R^{2}$ for the first-stage regressions; if the $R^{2}$ is low, we should suspect the large-sample approximation to the distribution of the 2SLS estimator to be poor (you will see a dramatic example in part (g) of the empirical exercise). Recently, Staiger and Stock (1997) proposed an alternative asymptotic approximation to the finite-sample distribution of the 2SLS and other estimators and associated test statistics, for the case of "weak" instruments. Their device is to look at a sequence of models along which the coefficients of the instruments in the first-stage regressions converge to zero. (Analytical Exercise 10 works out this type of asymptotics for the simple case of one regressor and one instrument.)

\section*{Alternative Derivations of 2SLS}
If we define data matrices as

$$
\underset{(n \times K)}{\mathbf{X}}=\left[\begin{array}{c}
\mathbf{x}_{1}^{\prime} \\
\mathbf{x}_{2}^{\prime} \\
\vdots \\
\mathbf{x}_{n}^{\prime}
\end{array}\right], \quad \underset{(n \times L)}{\mathbf{Z}}=\left[\begin{array}{c}
\mathbf{z}_{1}^{\prime} \\
\mathbf{z}_{2}^{\prime} \\
\vdots \\
\mathbf{z}_{n}^{\prime}
\end{array}\right], \quad \underset{(n \times 1)}{\mathbf{y}}=\left[\begin{array}{c}
y_{1} \\
y_{2} \\
\vdots \\
y_{n}
\end{array}\right] \text {, }
$$

then it is easy to see that the 2SLS estimator and associated statistics can be written as


\begin{align*}
\hat{\boldsymbol{\delta}}_{2 \text { SLS }} & =\left[\mathbf{Z}^{\prime} \mathbf{X}\left(\mathbf{X}^{\prime} \mathbf{X}\right)^{-1} \mathbf{X}^{\prime} \mathbf{Z}\right]^{-1} \mathbf{Z}^{\prime} \mathbf{X}\left(\mathbf{X}^{\prime} \mathbf{X}\right)^{-1} \mathbf{X}^{\prime} \mathbf{y} \\
& =\left(\mathbf{Z}^{\prime} \mathbf{P} \mathbf{Z}\right)^{-1} \mathbf{Z}^{\prime} \mathbf{P} \mathbf{y}, \tag{$\prime$}
\end{align*}


where $\mathbf{P} \equiv \mathbf{X}\left(\mathbf{X}^{\prime} \mathbf{X}\right)^{-1} \mathbf{X}^{\prime}$ is the projection matrix,


\begin{gather*}
\widehat{\operatorname{Avar}\left(\hat{\boldsymbol{\delta}}_{2 \text { SLS }}\right)}=n \cdot \hat{\sigma}^{2} \cdot\left[\mathbf{Z}^{\prime} \mathbf{X}\left(\mathbf{X}^{\prime} \mathbf{X}\right)^{-1} \mathbf{X}^{\prime} \mathbf{Z}\right]^{-1}=n \cdot \hat{\sigma}^{2} \cdot\left(\mathbf{Z}^{\prime} \mathbf{P} \mathbf{Z}\right)^{-1} .  \tag{3.8.5'}\\
\hat{\sigma}^{2} \equiv \frac{\hat{\boldsymbol{\varepsilon}}^{\prime} \hat{\boldsymbol{\varepsilon}}}{n} \quad \text { where } \hat{\boldsymbol{\varepsilon}} \equiv \mathbf{y}-\mathbf{Z} \hat{\boldsymbol{\delta}}_{2 \text { SLS }} .  \tag{$\prime$}\\
t_{\ell}=\frac{\hat{\delta}_{2 \text { SLS }, \ell}-\bar{\delta}_{\ell}}{\sqrt{\hat{\sigma}^{2} \cdot\left(\left[\mathbf{Z}^{\prime} \mathbf{X}\left(\mathbf{X}^{\prime} \mathbf{X}\right)^{-1} \mathbf{X}^{\prime} \mathbf{Z}\right]^{-1}\right)_{\ell \ell}}} .  \tag{$\prime$}\\
W=\frac{\mathbf{a}\left(\hat{\boldsymbol{\delta}}_{2 \text { SLS }}\right)^{\prime}\left[\mathbf{A}\left(\hat{\boldsymbol{\delta}}_{2 \text { SLS }}\right)\left(\mathbf{Z}^{\prime} \mathbf{X}\left(\mathbf{X}^{\prime} \mathbf{X}\right)^{-1} \mathbf{X}^{\prime} \mathbf{Z}\right)^{-1} \mathbf{A}\left(\hat{\delta}_{2 \text { SLS }}\right)^{\prime}\right]^{-1} \mathbf{a}\left(\hat{\boldsymbol{\delta}}_{2 \text { SLS }}\right)}{\hat{\sigma}^{2}} .  \tag{$\prime$}\\
J\left(\tilde{\boldsymbol{\delta}},\left(\hat{\sigma}^{2} \cdot \mathbf{S}_{\mathbf{x} \mathbf{x}}\right)^{-1}\right)=\frac{(\mathbf{y}-\mathbf{Z} \tilde{\boldsymbol{\delta}})^{\prime} \mathbf{P}(\mathbf{y}-\mathbf{Z} \tilde{\boldsymbol{\delta}})}{\hat{\sigma}^{2}} .  \tag{$\prime$}\\
\text { Sargan's statistic }=\frac{\hat{\boldsymbol{\varepsilon}}^{\prime} \mathbf{P} \hat{\boldsymbol{\varepsilon}}}{\hat{\sigma}^{2}} . \tag{3.8.10'}
\end{gather*}


Using this formula, we can provide two other derivations of the 2SLS estimator.

\section*{2SLS as an IV Estimator}
Let $\hat{\mathbf{z}}_{i}(L \times 1)$ be the vector of $L$ instruments (which will be generated from $\mathbf{x}_{i}$ as described below) for the $L$ regressors, and let $\widehat{\mathbf{Z}}$ be the $n \times L$ data matrix of those $L$ instruments. So the $i$-th row of $\widehat{\mathbf{Z}}$ is $\hat{\mathbf{z}}_{i}^{\prime}$. The IV estimator of $\boldsymbol{\delta}$ with $\hat{\mathbf{z}}_{i}$ serving as instruments is, by (3.4.4),


\begin{equation*}
\hat{\boldsymbol{\delta}}_{\mathrm{IV}}=\left(\frac{1}{n} \sum_{i=1}^{n} \hat{\mathbf{z}}_{i} \mathbf{z}_{i}^{\prime}\right)^{-1} \frac{1}{n} \sum_{i=1}^{n} \hat{\mathbf{z}}_{i} \cdot y_{i}=\left(\widehat{\mathbf{Z}}^{\prime} \mathbf{Z}\right)^{-1} \widehat{\mathbf{Z}}^{\prime} \mathbf{y} \tag{3.8.12}
\end{equation*}


Now we generate those $L$ instruments from $\mathbf{x}_{i}$ as follows. The $\ell$-th instrument is the fitted value from regressing $z_{i \ell}$ (the $\ell$-th regressor) on $\mathbf{x}_{i}$. The $n$-vector of fitted value is $\mathbf{X}\left(\mathbf{X}^{\prime} \mathbf{X}\right)^{-1} \mathbf{X}^{\prime} \mathbf{z}_{\ell}$, where $\mathbf{z}_{\ell}$ is the $n$-vector of the $\ell$-th regressor (i.e., the $\ell$-th column of $\mathbf{Z}$ ). Therefore, the $n \times L$ data matrix of instruments is


\begin{equation*}
\widehat{\mathbf{Z}}=\left(\mathbf{X}\left(\mathbf{X}^{\prime} \mathbf{X}\right)^{-1} \mathbf{X}^{\prime} \mathbf{z}_{1}, \ldots, \mathbf{X}\left(\mathbf{X}^{\prime} \mathbf{X}\right)^{-1} \mathbf{X}^{\prime} \mathbf{z}_{L}\right)=\mathbf{X}\left(\mathbf{X}^{\prime} \mathbf{X}\right)^{-1} \mathbf{X}^{\prime} \mathbf{Z}=\mathbf{P} \mathbf{Z} \tag{3.8.13}
\end{equation*}


where $\mathbf{P}$ is the projection matrix. Substituting this into the IV format (3.8.12) yields the 2SLS estimator (see (3.8.3')).

\section*{2SLS as Two Regressions}
Instead of substituting the generated instruments $\hat{\mathbf{z}}_{i}$ into the IV format to estimate the equation $y_{i}=\mathbf{z}_{i}^{\prime} \boldsymbol{\delta}+\varepsilon_{i}$, consider regressing $y_{i}$ on $\hat{\mathbf{z}}_{i}$. The coefficient estimate is

$$
\begin{aligned}
& \left(\widehat{\mathbf{Z}}^{\prime} \mathbf{Z}\right)^{-1} \widehat{\mathbf{Z}}^{\prime} \mathbf{y}=\left(\mathbf{Z}^{\prime} \mathbf{P}^{\prime} \mathbf{P} \mathbf{Z}\right)^{-1} \mathbf{Z}^{\prime} \mathbf{P} \mathbf{y} \\
& =\left(\mathbf{Z}^{\prime} \mathbf{P} \mathbf{Z}\right)^{-1} \mathbf{Z}^{\prime} \mathbf{P} \mathbf{y} \quad \text { (since } \mathbf{P} \text { is symmetric and idempotent). }
\end{aligned}
$$

This is the 2SLS estimator. So the 2SLS coefficient estimate can be obtained in two stages: the first stage is to regress the $L$ regressors on $\mathbf{x}_{i}$ and obtain fitted values $\hat{\mathbf{z}}_{i}$, and the second stage is to regress $y_{i}$ on those fitted values.

For those regressors that are predetermined and hence included in $\mathbf{x}_{i}$, there is no need to carry out the first-stage regression because the fitted value is the regressor itself. To see this, if $z_{i \ell}$ is predetermined and included in $\mathbf{x}_{i}$ as the $k$-th instrument, the $n$-vector of fitted values for the $\ell$-th regressor is $\mathbf{P} \mathbf{z}_{\ell}$, where $\mathbf{z}_{\ell}$ is the $n$-vector whose $i$-th element is $z_{i \ell}$. But since $\mathbf{z}_{\ell}$ is also the $k$-th column of $\mathbf{X}, \mathbf{P z}_{\ell}=\mathbf{P x}_{k}$. Since $\mathbf{P}$ is the projection matrix, $\mathbf{P x}_{k}=\mathbf{x}_{k}$.

This derivation of the 2SLS is useful as it justifies the naming of the estimator, but there is a pitfall. In the second-stage regression where $y_{i}$ is regressed on $\hat{\mathbf{z}}_{i}$, the standard errors routinely calculated by the OLS package are based on the residual vector $\mathbf{y}-\widehat{\mathbf{Z}} \hat{\boldsymbol{\delta}}_{2 \text { SLS }}$. This does not equal the 2SLS residual $\mathbf{y}-\mathbf{Z} \hat{\boldsymbol{\delta}}_{\text {2SLS }}$. Therefore, the OLS standard errors and estimated asymptotic variance from the second stage cannot be used for statistical inference.

\section*{When Regressors Are Predetermined}
When all the regressors are predetermined and the errors are conditionally homoskedastic, there is a close connection between the distance function $J$ for efficient GMM and the sum of squared residuals ( $S S R$ ). From (3.8.9') on page 230,


\begin{align*}
J\left(\tilde{\boldsymbol{\delta}},\left(\hat{\sigma}^{2} \cdot \mathbf{S}_{\mathbf{x} \mathbf{x}}\right)^{-1}\right) & =\frac{(\mathbf{y}-\mathbf{Z} \tilde{\boldsymbol{\delta}})^{\prime} \mathbf{P}(\mathbf{y}-\mathbf{Z} \tilde{\boldsymbol{\delta}})}{\hat{\sigma}^{2}} \\
& =\frac{\mathbf{y}^{\prime} \mathbf{P} \mathbf{y}-2 \mathbf{y}^{\prime} \mathbf{P} \mathbf{Z} \tilde{\boldsymbol{\delta}}+\tilde{\boldsymbol{\delta}}^{\prime} \mathbf{Z}^{\prime} \mathbf{P} \mathbf{Z} \tilde{\boldsymbol{\delta}}}{\hat{\sigma}^{2}} \\
& =\frac{\mathbf{y}^{\prime} \mathbf{P} \mathbf{y}-2 \mathbf{y}^{\prime} \mathbf{Z} \tilde{\boldsymbol{\delta}}+\tilde{\boldsymbol{\delta}}^{\prime} \mathbf{Z}^{\prime} \mathbf{Z} \tilde{\boldsymbol{\delta}}}{\hat{\sigma}^{2}}\left(\text { since } \mathbf{P} \mathbf{Z}=\mathbf{Z} \text { when } \mathbf{z}_{i} \subset \mathbf{x}_{i}\right) \\
& =\frac{(\mathbf{y}-\mathbf{Z} \tilde{\boldsymbol{\delta}})^{\prime}(\mathbf{y}-\mathbf{Z} \tilde{\boldsymbol{\delta}})}{\hat{\sigma}^{2}}-\frac{\mathbf{y}^{\prime} \mathbf{y}-\mathbf{y}^{\prime} \mathbf{P} \mathbf{y}}{\hat{\sigma}^{2}} \\
& =\frac{(\mathbf{y}-\mathbf{Z} \tilde{\boldsymbol{\delta}})^{\prime}(\mathbf{y}-\mathbf{Z} \tilde{\boldsymbol{\delta}})}{\hat{\sigma}^{2}}-\frac{(\mathbf{y}-\hat{\mathbf{y}})(\mathbf{y}-\hat{\mathbf{y}})^{\prime}}{\hat{\sigma}^{2}} \tag{3.8.14}
\end{align*}


where $\hat{\mathbf{y}} \equiv \mathbf{P y}$ is the vector of fitted values from unrestricted OLS. Since the last term does not depend on $\tilde{\delta}$, minimizing $J$ amounts to minimizing the sum of squared residuals $(\mathbf{y}-\mathbf{Z} \tilde{\boldsymbol{\delta}})^{\prime}(\mathbf{y}-\mathbf{Z} \tilde{\boldsymbol{\delta}})$. It then follows that (1) the efficient GMM estimator is OLS (which actually is true without conditional homoskedasticity as long as $\mathbf{z}_{i}=\mathbf{x}_{i}$ ), (2) the restricted efficient GMM estimator subject to the constraints of the null hypothesis is the restricted OLS (whose objective function is not $J$ but $S S R$ ), and (3) the Wald statistic, which is numerically equal to the $L R$ statistic, can be calculated as the difference in $S S R$ with and without the imposition of the null, normalized to $\hat{\sigma}^{2}$. This last result confirms the derivation in Section 2.6 of the Wald statistic by the Likelihood Ratio principle.

\section*{Testing a Subset of Orthogonality Conditions}
In Section 3.6 we introduced the statistic $C$ for testing a subset of orthogonality conditions. It utilizes two efficient GMM estimators of the same equation, one using the full set $\mathbf{x}_{i}$ of instruments and the other using only a subset, $\mathbf{x}_{i 1}$, of $\mathbf{x}_{i}$. To examine what expression it takes under conditional homoskedasticity, let $\mathbf{X}_{1}$ ( $n \times K_{1}$ ) be the data matrix whose $i$-th row is $\mathbf{x}_{i 1}^{\prime}$. Because the two-step efficient GMM estimator is the 2SLS estimator under conditional homoskedasticity, the two GMM estimators are given by


\begin{align*}
& \hat{\delta}=\left(\mathbf{Z}^{\prime} \mathbf{P} \mathbf{Z}\right)^{-1} \mathbf{Z}^{\prime} \mathbf{P} \mathbf{y} \quad \text { with } \quad \mathbf{P}=\mathbf{X}\left(\mathbf{X}^{\prime} \mathbf{X}\right)^{-1} \mathbf{X}^{\prime}  \tag{3.8.15}\\
& \bar{\delta}=\left(\mathbf{Z}^{\prime} \mathbf{P}_{1} \mathbf{Z}\right)^{-1} \mathbf{Z}^{\prime} \mathbf{P}_{1} \mathbf{y} \quad \text { with } \quad \mathbf{P}_{1}=\mathbf{X}_{1}\left(\mathbf{X}_{1}^{\prime} \mathbf{X}_{1}\right)^{-1} \mathbf{X}_{1}^{\prime} . \tag{3.8.16}
\end{align*}


And the $C$ statistic becomes the difference in two Sargan statistics:


\begin{equation*}
C=\frac{\hat{\varepsilon}^{\prime} \mathbf{P} \hat{\varepsilon}-\bar{\varepsilon}^{\prime} \mathbf{P}_{1} \bar{\varepsilon}}{\hat{\sigma}^{2}}, \tag{3.8.17}
\end{equation*}


where

$$
\hat{\varepsilon} \equiv \mathbf{y}-\mathbf{Z} \hat{\delta}, \quad \bar{\varepsilon} \equiv \mathbf{y}-\mathbf{Z} \bar{\delta}, \quad \hat{\sigma}^{2} \equiv \frac{\hat{\varepsilon}^{\prime} \hat{\varepsilon}}{n}
$$

As seen for the case without conditional homoskedasticity, $C$ is guaranteed to be nonnegative in finite samples if the same matrix $\widehat{\mathbf{S}}$ is used throughout, which under conditional homoskedasticity amounts to using the same estimate of the error variance, $\hat{\sigma}^{2}$, to deflate both $\hat{\boldsymbol{\varepsilon}}^{\prime} \mathbf{P} \hat{\boldsymbol{\varepsilon}}$ and $\overline{\boldsymbol{\varepsilon}}^{\prime} \mathbf{P}_{1} \overline{\boldsymbol{\varepsilon}}$, as in (3.8.17). By Proposition 3.7, this $C$ is asymptotically $\chi^{2}\left(K-K_{1}\right)$.

Perhaps more popular is the test statistic proposed by Hausman and Taylor (1980) and further examined by Newey (1985). $\hat{\delta}$ is asymptotically more efficient than $\bar{\delta}$ because $\hat{\delta}$ exploits more orthogonality conditions. Therefore, $\operatorname{Avar}(\bar{\delta}) \geq$

Avar( $\hat{\boldsymbol{\delta}}$ ). Furthermore, as you will be asked to show (Analytical Exercise 9), under the same hypothesis guaranteeing $C$ to be asymptotically chi-squared,


\begin{equation*}
\operatorname{Avar}(\overline{\boldsymbol{\delta}}-\hat{\boldsymbol{\delta}})=\operatorname{Avar}(\overline{\boldsymbol{\delta}})-\operatorname{Avar}(\hat{\boldsymbol{\delta}}) \tag{3.8.18}
\end{equation*}


(If you have been exposed to the Hausman test in the ML context, you can recognize this as the GMM version of the Hausman principle.) By Proposition 3.9, $\operatorname{Avar}(\hat{\delta})$ is consistently estimated by


\begin{equation*}
\widehat{\operatorname{Avar}(\hat{\boldsymbol{\delta}})} \equiv n \cdot \hat{\sigma}^{2} \cdot\left(\mathbf{Z}^{\prime} \mathbf{P} \mathbf{Z}\right)^{-1} . \tag{3.8.19}
\end{equation*}


Similarly, a consistent estimator of $\operatorname{Avar}(\overline{\boldsymbol{\delta}})$ is


\begin{equation*}
\widehat{\operatorname{Avar}(\overline{\boldsymbol{\delta}})} \equiv n \cdot \hat{\sigma}^{2} \cdot\left(\mathbf{Z}^{\prime} \mathbf{P}_{1} \mathbf{Z}\right)^{-1} . \tag{3.8.20}
\end{equation*}


Here, as in the calculation of the $C$ statistic, the same estimate, $\hat{\sigma}^{2}$, is used throughout, in order to guarantee the test statistic below to be nonnegative. The resulting estimator of $\operatorname{Avar}(\overline{\boldsymbol{\delta}}-\hat{\boldsymbol{\delta}})$ is


\begin{equation*}
\widehat{\operatorname{Avar}(\bar{\delta}-\hat{\delta})}=\hat{\sigma}^{2} \cdot n \cdot\left[\left(\mathbf{Z}^{\prime} \mathbf{P}_{1} \mathbf{Z}\right)^{-1}-\left(\mathbf{Z}^{\prime} \mathbf{P} \mathbf{Z}\right)^{-1}\right] . \tag{3.8.21}
\end{equation*}


Hausman and Taylor (1980) have shown that (1) this matrix in finite samples is positive semidefinite (nonnegative definite) but not necessarily nonsingular, but (2) for any generalized inverse ${ }^{17}$ of this matrix, the Hausman statistic


\begin{align*}
H & \equiv \sqrt{n}(\overline{\boldsymbol{\delta}}-\hat{\boldsymbol{\delta}})^{\prime}\left\{\hat{\sigma}^{2} \cdot n \cdot\left[\left(\mathbf{Z}^{\prime} \mathbf{P}_{1} \mathbf{Z}\right)^{-1}-\left(\mathbf{Z}^{\prime} \mathbf{P} \mathbf{Z}\right)^{-1}\right]\right\}^{-} \sqrt{n}(\overline{\boldsymbol{\delta}}-\hat{\boldsymbol{\delta}}) \\
& =\frac{(\overline{\boldsymbol{\delta}}-\hat{\boldsymbol{\delta}})^{\prime}\left[\left(\mathbf{Z}^{\prime} \mathbf{P} 1 \mathbf{Z}\right)^{-1}-\left(\mathbf{Z}^{\prime} \mathbf{P} \mathbf{Z}\right)^{-1}\right]^{-}(\overline{\boldsymbol{\delta}}-\hat{\boldsymbol{\delta}})}{\hat{\sigma}^{2}} \tag{3.8.22}
\end{align*}


is invariant to the choice of the generalized inverse and is asymptotically chisquared with $\min \left(K-K_{1}, L-s\right)$ degrees of freedom, where

$$
s=\# \mathbf{z}_{i} \cap \mathbf{x}_{i 1}=\text { number of regressors which are retained as instruments in } \mathbf{x}_{i 1} \text {. }
$$

What is the relationship between $C$ and $H$ under conditional homoskedasticity? It can be shown (see Newey, 1985) that:

\footnotetext{${ }^{17}$ A generalized inverse, $\mathbf{A}^{-}$, of a matrix $\mathbf{A}$ is any matrix satisfying $\mathbf{A A}^{-} \mathbf{A}=\mathbf{A}$. If $\mathbf{A}$ is square and nonsingular, then $\mathbf{A}^{-}$is unique and equal to $\mathbf{A}^{-1}$.
}If $K-K_{1} \leq L-s$ so that both $H$ and $C$ have the same degrees of freedom, then $H=C$ (numerically equal). Otherwise, the two statistics are numerically different and have different degrees of freedom.

One frequent case where $K-K_{1} \leq L-s$ holds is when $\mathbf{x}_{i 2}$ is a subset of $\mathbf{z}_{i}$, that is, when the suspect instruments are a subset of regressors. In this case, the formulas (3.8.17) and (3.8.22) are two alternative ways of calculating numerically the same statistic. In the case where $K-K_{1}>L-s$, the degrees of freedom for $H$ are less than those for the $C$ statistic ( $K-K_{1}$ ). For this reason, the Hausman test, unlike the $C$ test, is not consistent against some local alternatives. ${ }^{18}$

\section*{Testing Conditional Homoskedasticity}
For the OLS case with predetermined regressors, as shown in Section 2.7, there is a convenient $n R^{2}$ test of conditional homoskedasticity. In the present case, the $n R^{2}$ statistic obtained by regressing the squared 2SLS residuals on a constant and second-order cross products of the instrumental variables turns out not to have the desired asymptotic distribution. A test statistic that is asymptotically chi-squared is available but is extremely cumbersome. See White (1982, note 2).

\section*{Testing for Serial Correlation}
For the OLS case, we developed in Section 2.10 tests for serial correlation in the error term. More specifically, under (2.10.15) (which is stronger than Assumption 2.3 or 3.3 ) and (2.10.16) (which is stronger than Assumption 2.7 or 3.7 of conditional homoskedasticity), the modified Box-Pierce $Q$ given in (2.10.20) can be used for testing the null of no serial correlation in the error term, and this statistic is asymptotically equivalent (in the stronger sense of the plim of the difference being zero) to the $n R^{2}$ statistic from regression (2.10.21). Can this test be extended to the case where the regressors $\mathbf{z}_{i}$ are endogenous? If the instruments $\mathbf{x}_{i}$ satisfy (2.10.15) and (2.10.16), then the argument in Appendix 2.B of Chapter 2 can be generalized to produce a modified $Q$ statistic that is asymptotically chi-squared under the null of no serial correlation. However, the expression for the statistic is more complicated than (2.10.20) and so is not presented here. This modified $Q$ statistic for the case of endogenous regressors does not seem to be asymptotically equivalent to the $n R^{2}$ statistic from regression (2.10.21).

\footnotetext{${ }^{18}$ The Hausman statistic can be generalized to the case of conditional heteroskedasticity, but it is not practical because the degrees of freedom depend on the unknown values of the matrices $\boldsymbol{\Sigma}_{\mathbf{X Z}}$ and $\mathbf{S}$. See Newey (1985).
}\section*{QUESTIONS FOR REVIEW}
\begin{enumerate}
  \item (GMM with conditional homoskedasticity) In efficient two-step GMM estimation, $\widehat{\mathbf{S}}$ in the second step is calculated from the first step by (3.5.10) using the first-step residuals. Under conditional homoskedasticity, derive the asymptotic variance of the two-step estimator. Is it the same as (3.8.4)? Hint: Under conditional homoskedasticity, (3.5.10) $\rightarrow_{\mathrm{p}} \sigma^{2} \boldsymbol{\Sigma}_{\mathbf{x x}}$.
  \item (2SLS without conditional homoskedasticity) Is the 2SLS consistent when conditional homoskedasticity does not hold? Derive the plim of (3.8.2) and $\operatorname{Avar}\left(\hat{\boldsymbol{\delta}}_{\text {2SLS }}\right)$ without assuming conditional homoskedasticity. Is the 2SLS as efficient as the two-step GMM (i.e., is its asymptotic variance as small as (3.5.13))? Hint: 2SLS is a GMM estimator with a choice of $\widehat{\mathbf{W}}$ that is not necessarily efficient without conditional homoskedasticity.
  \item (GLS interpretation of 2SLS) Verify that the 2SLS estimator can be written as a GLS estimator if $\mathbf{S}_{\mathbf{x z}}$ and $\mathbf{s}_{\mathbf{x y}}$ are interpreted as the data matrix of regressors and the data vector of the dependent variable and $\mathbf{S}_{\mathbf{x x}}$ as the variance matrix of the error term.
  \item Provide two cases in which $\hat{\boldsymbol{\delta}}_{\text {2SLS }}$ and $\hat{\boldsymbol{\delta}}\left(\widehat{\mathbf{S}}^{-1}\right)$ are asymptotically equivalent in the sense that
\end{enumerate}

$$
\sqrt{n}\left[\hat{\boldsymbol{\delta}}_{2 \mathrm{SLS}}-\hat{\boldsymbol{\delta}}\left(\widehat{\mathbf{S}}^{-1}\right)\right] \underset{\mathrm{p}}{\rightarrow} \mathbf{0} .
$$

Hint: Keywords are "conditional homoskedasticity" and "exactly identified."\\
5. Suppose the equation is just identified. Show that 2SLS (3.8.3) reduces to IV (3.4.4).\\
6. (Sargan as $n R^{2}$ ) Prove that Sargan's statistic (3.8.10') equals $n R_{u c}^{2}$, where $R_{u c}^{2}$ is the uncentered $R$-squared from a regression of $\hat{\varepsilon}$ on $\mathbf{X}$. Hint: Review Question 8 of Section 1.2.\\
7. (When $\mathbf{z}_{i}$ is a strict subset of $\mathbf{x}_{i}$ ) Suppose $\mathbf{z}_{i}$ is a strict subset of $\mathbf{x}_{i}$. So $\mathbf{x}_{i}$ includes, in addition to the regressors (which are all predetermined), some other predetermined variables. We have shown that the efficient GMM estimator is OLS under conditional homoskedasticity. Does the result remain true without conditional homoskedasticity? Hint: Review Question 5 of Section 3.5.

\subsection*{3.9 Application: Returns from Schooling}
Since Mincer's (1958) pioneering study, the relationship between the wage rate and schooling has been the subject of a large number of empirical and theoretical investigations. You might find the amount of attention puzzling because the explanation of the positive relationship seems to be obvious: education enhances the individual's productivity. There are, however, other explanations. For example, in the Spencian job market signaling model, more educated individuals receive higher wages only because education is used as a signal of higher ability; although education does not increase the individual's earning capacity, there is a correlation between the wage rate and schooling because both variables are influenced by the third variable, ability. This section shows how to use the technique of this chapter to isolate the effect of education on the wage rate from that of ability. One of the earliest studies to address this issue is Griliches (1976).

\section*{The NLS-Y Data}
The data used by Griliches are the Young Men's Cohort of the National Longitudinal Survey (NLS-Y). This cohort was first surveyed in 1966 at ages 14-24, with 5,225 respondents, and was resurveyed at one- or two-year intervals thereafter. By 1969, about a quarter of the original sample was lost, but there are 2,026 individuals who reported earnings in 1969 and whose records are complete enough to allow derivation of all the variables for analysis. A very attractive feature of the NLS-Y is its inclusion of two measures of ability. One of them is the score on the Knowledge of the World of Work (KWW) test administered by the NLS interviewers in 1966. The other measure is the IQ score. All youths in the survey who had completed ninth grade by 1966 were asked to sign waivers letting their school supply the survey administrator their scores on various tests and other background materials. The resulting School Survey, conducted in 1968, yielded data on different mental ability scores, which were combined into IQ equivalents. Of the 2,026 individuals with information about the 1969 wage rate and other variables, the IQ score is available for 1,362 individuals, reflecting the fact that the School Survey was able to cover only two-thirds of the original sample. Our discussion here concerns Griliches's results based on this smaller sample. Table 3.1 reports the means and standard deviation of key variables.

Means and Standard Deviations (in parentheses)

\begin{table}[h]
\begin{center}
\captionsetup{labelformat=empty}
\caption{Table 3.1: Characteristics of Young Men from the National Longitudinal Survey}
\begin{tabular}{|l|l|}
\hline
\multicolumn{2}{|l|}{Variable} \\
\hline
Sample size & 1,362 \\
\hline
Age in 1969 & 22.3 (3.2) \\
\hline
Schooling in years in 1969 (S) & 12.5 (1.9) \\
\hline
Logarithm of hourly wages (in cents) in 1969 (LW) & 5.68 (0.40) \\
\hline
Score on the Knowledge of the World of Work test ( $K W W$ ) & 35.1 (7.9) \\
\hline
IQ score (IQ) & 97.7 (15.3) \\
\hline
Experience in years in 1969 & 3.7 (2.8) \\
\hline
\end{tabular}
\end{center}
\end{table}

Source: Griliches (1976, Table 1).

\section*{The Semi-Log Wage Equation}
The typical wage equation estimated in the literature is the semi-log form:


\begin{equation*}
L W=\alpha+\beta S+\gamma A+\delta^{\prime} \mathbf{h}+\varepsilon, \tag{3.9.1}
\end{equation*}


where $L W$ is the log wage rate for the individual, $S$ is schooling in years, $A$ is a measure of ability, $\mathbf{h}$ is the vector of observable characteristics of the individual (such as experience and location dummies), $\boldsymbol{\delta}$ is the associated vector of coefficients, and $\varepsilon$ is the unobservable error term with zero mean (in this section, the individual subscript $i$ will be dropped for notational simplicity). The semi-log specification for schooling $S$ is often justified by appealing to the well-established stylized fact from large cross-section data (such as the Current Population Survey) that the relationship between $\log$ wages and schooling is linear. ${ }^{19}$

The schooling coefficient $\beta$ measures the percentage increase in the wage rate the individual would receive if she had one more year of education. It therefore

\footnotetext{${ }^{19}$ For the functional form issue, see Card (1995).
}
represents the marginal return from investing in human capital, which should be in the same order of magnitude as the rate of return from financial assets. We assume that the nonconstant regressors, ( $S, A, \mathbf{h}$ ), are uncorrelated with the error term $\varepsilon$ so that the OLS is an appropriate estimation technique if ability $A$ is included in the regression along with $S$ and $\mathbf{h}$. In the rest of this section, we examine what Griliches called the "ability bias"-biases on the OLS estimate of $\beta$ that would arise when ability $A$ is not included in the regression and when its imperfect measure is included in its place.

\section*{Omitted Variable Bias}
Sometimes the data set you work with has no measures of $A$ (this is true for the Current Population Survey, for example). What is the consequence of ignoring ability by omitting $A$ from the wage equation? We know from Section 2.9 that the regression of $L W$ on a constant, $S$, and $\mathbf{h}$ provides a consistent estimator of the corresponding least squares projection, which can be written as


\begin{align*}
\widehat{\mathrm{E}}^{*}(L W \mid 1, S, \mathbf{h}) & =\widehat{\mathrm{E}}^{*}\left(\alpha+\beta S+\gamma A+\delta^{\prime} \mathbf{h}+\varepsilon \mid 1, S, \mathbf{h}\right) \quad(\text { by }(3.9 .1)) \\
& =\alpha+\beta S+\delta^{\prime} \mathbf{h}+\gamma \widehat{\mathrm{E}}^{*}(A \mid 1, S, \mathbf{h})+\widehat{\mathrm{E}}^{*}(\varepsilon \mid 1, S, \mathbf{h}) \tag{3.9.2}
\end{align*}


Since the regressors in (3.9.1) are all predetermined, we have $\mathrm{E}(\varepsilon)=0, \mathrm{E}(S \varepsilon)=$ 0 , and $\mathrm{E}(\mathbf{h} \cdot \varepsilon)=\mathbf{0}$. So $\widehat{\mathrm{E}}^{*}(\varepsilon \mid 1, S, \mathbf{h})=0$. Writing $\widehat{\mathrm{E}}^{*}(A \mid 1, S, \mathbf{h})=\theta_{1}+\theta_{S} S+ \boldsymbol{\theta}_{\mathbf{h}}^{\prime} \mathbf{h}$, (3.9.2) becomes

$$
\widehat{\mathrm{E}}^{*}(L W \mid 1, S, \mathbf{h})=\left(\alpha+\gamma \theta_{1}\right)+\left(\beta+\gamma \theta_{S}\right) S+\left(\boldsymbol{\delta}+\gamma \cdot \boldsymbol{\theta}_{\mathbf{h}}\right)^{\prime} \mathbf{h} .
$$

Therefore, the OLS coefficient estimate can be asymptotically biased for all the included regressors. This phenomenon is called the omitted variable bias. In particular,

$$
\operatorname{plim} \widehat{\beta}_{\mathrm{OLS}}=\beta+\gamma \theta_{S} .
$$

That is, the OLS estimate of the schooling coefficient $\beta$ includes the indirect effect of ability on log wages through schooling ( $\gamma \theta_{S}$ ) as well as the direct effect of schooling ( $\beta$ ). If $\theta_{S}$ is positive, then $\widehat{\beta}_{\text {OLS }}$ is asymptotically biased upward.

Using the sample of 1,362 individuals from the NLS-Y described above, Griliches estimated the wage equation for 1969. The list of variables included in $\mathbf{h}$ will not be given here; suffice to say that it includes experience in years and some region and city-size dummies. Griliches's estimate of the schooling coefficient when ability is ignored in the wage equation is reproduced on line 1 of Table 3.2.

On the face of it, the estimate looks good: the point estimate resembles the rate of return one might get from financial assets, and it is sharply estimated as evidenced by the high $t$-value.

\section*{IQ as the Measure of Ability}
As already mentioned, the NLS-Y has two measures of ability, $K W W$ (the score on the Knowledge of the World of Work test) collected in 1966 and IQ (the IQ score). Since the NLS-Y respondents were at least fourteen years old in 1966, $K W W$ would reflect the effect of schooling already undertaken, and so it cannot be a measure of raw ability. The IQ score does not have this problem. If IQ were a perfect measure of ability so that $A$ can be equated with it, then the wage equation (3.9.1) could be estimated consistently with $I Q$ substituting for $A$. Griliches's OLS estimates of $\beta$ (schooling coefficient) and $\gamma$ (ability coefficient) when $I Q$ is included in the equation are reported in line 2 of Table 3.2. Now the estimated schooling coefficient is lower, confirming our prediction that the estimated schooling coefficient includes the effect of ability on the wage rate when ability is omitted from the regression.

\section*{Errors-in-Variables}
Of course the IQ score may not be an error-free measure of ability. If $\eta$ is the measurement error, $I Q$ is related to $A$ as


\begin{equation*}
I Q=\phi+A+\eta, \tag{3.9.3}
\end{equation*}


with $\mathrm{E}(\eta)=0$. The interpretation of $\eta$ as measurement error makes it reasonable to assume that $\eta$ is uncorrelated with $A, S, \mathbf{h}$, and the wage equation error term $\varepsilon$. Substituting (3.9.3) into (3.9.1), we obtain


\begin{equation*}
L W=(\alpha-\gamma \phi)+\beta S+\gamma I Q+\delta^{\prime} \mathbf{h}+(\varepsilon-\gamma \eta) . \tag{3.9.4}
\end{equation*}


We now illustrate for this example the general result that, if at least one of the regressors is measured with error, the OLS estimates of all the regression coefficients can be asymptotically biased.

To examine the consequence of using the error-ridden measure $I Q$ for $A$, consider the corresponding least squares projection:


\begin{align*}
\widehat{\mathrm{E}}^{*}(L W \mid 1, S, I Q, \mathbf{h})= & (\alpha-\gamma \phi)+\beta S+\gamma I Q+\delta^{\prime} \mathbf{h} \\
& +\widehat{\mathrm{E}}^{*}(\varepsilon \mid 1, S, I Q, \mathbf{h})-\gamma \widehat{\mathrm{E}}^{*}(\eta \mid 1, S, I Q, \mathbf{h}) . \tag{3.9.5}
\end{align*}


Table 3.2: Parameter Estimates

Estimation equation: $L W=\alpha+\beta S+\gamma I Q+$ other controls

\begin{center}
\begin{tabular}{|l|l|l|l|l|l|l|l|}
\hline
\multirow{2}{*}{Line no.} & \multirow[b]{2}{*}{Estimation technique} & \multicolumn{2}{|c|}{Coefficient of} & \multirow{2}{*}{SER} & \multirow{2}{*}{$R^{2}$} & \multirow[b]{2}{*}{Which regressor is endogenous?} & \multirow{2}{*}{Excluded predetermined variables} \\
\hline
 &  & $S$ & $I Q$ &  &  &  &  \\
\hline
1 & OLS & 0.065 (13.2) & - & 0.332 & 0.309 & none & - \\
\hline
2 & OLS & 0.059 (10.7) & 0.0019 (2.8) & 0.331 & 0.313 & none & - \\
\hline
3 & 2SLS & 0.052 (7.0) & 0.0038 (2.4) & 0.332 & 一 & IQ & $M E D, K W W$, age, age squared, background variables \\
\hline
\end{tabular}
\end{center}

\footnotetext{Source: Line 1: equation (B1) in Griliches (1976. Table 2). Line 2: equation (B3) in Griliches's Table 2.\\
Line 3: line 3 in Griliches's Table 5. Figures in parentheses are $t$-values rather than standard errors.
}Now consider $\widehat{\mathrm{E}}^{*}(\varepsilon \mid 1, S, I Q, \mathbf{h})$ in (3.9.5). Since $\mathrm{E}(\varepsilon)=0$ and $(S, \mathbf{h})$ are uncorrelated with $\varepsilon,(1, S, \mathbf{h})$ are orthogonal to $\varepsilon . I Q$ is also orthogonal to $\varepsilon$ because

$$
\begin{aligned}
\mathrm{E}(I Q \varepsilon) & =\mathrm{E}[(\phi+A+\eta) \varepsilon] \quad(\text { by }(3.9 .3)) \\
& =\mathrm{E}(\eta \varepsilon) \quad(\text { since } \mathrm{E}(\varepsilon)=0, \operatorname{Cov}(A, \varepsilon)=0) \\
& =0
\end{aligned}
$$

Thus, $\widehat{\mathrm{E}}^{*}(\varepsilon \mid 1, S, I Q, \mathbf{h})=0$ in (3.9.5).\\
So the biases, if any, equal $-\gamma \widehat{\mathrm{E}}^{*}(\eta \mid 1, S, I Q, \mathbf{h})$ in (3.9.5). Let $\left(\theta_{S}, \theta_{I Q}, \theta_{\mathbf{h}}\right)$ be the projection coefficients of $(S, I Q, \mathbf{h})$ in $\widehat{\mathrm{E}}^{*}(\eta \mid 1, S, I Q, \mathbf{h})$. By the formula (2.9.7) from Chapter 2,

\[
\left[\begin{array}{c}
\theta_{S}  \tag{3.9.6}\\
\theta_{I Q} \\
\theta_{\mathbf{h}}
\end{array}\right]=\left[\begin{array}{ccc}
\operatorname{Var}(S) & \operatorname{Cov}(S, I Q) & \operatorname{Cov}\left(S, \mathbf{h}^{\prime}\right) \\
\operatorname{Cov}(I Q, S) & \operatorname{Var}(I Q) & \operatorname{Cov}\left(I Q, \mathbf{h}^{\prime}\right) \\
\operatorname{Cov}(\mathbf{h}, S) & \operatorname{Cov}(\mathbf{h}, I Q) & \operatorname{Var}(\mathbf{h})
\end{array}\right]^{-1}\left[\begin{array}{c}
\operatorname{Cov}(S, \eta) \\
\operatorname{Cov}(I Q, \eta) \\
\operatorname{Cov}(\mathbf{h}, \eta)
\end{array}\right] .
\]

In this expression, $\operatorname{Cov}(S, \eta)$ and $\operatorname{Cov}(\mathbf{h}, \eta)$ are zero by assumption. $\operatorname{Cov}(I Q, \eta)$, however, is not zero because

$$
\begin{aligned}
\operatorname{Cov}(I Q, \eta) & =\mathrm{E}(I Q \eta) \quad(\text { since } \mathrm{E}(\eta)=0) \\
& =\mathrm{E}[(\phi+A+\eta) \eta] \quad(\text { by }(3.9 .3)) \\
& =\phi \mathrm{E}(\eta)+\mathrm{E}(A \eta)+\mathrm{E}\left(\eta^{2}\right) \\
& =\operatorname{Var}(\eta) \quad(\text { since } \mathrm{E}(\eta)=0 \text { and } \operatorname{Cov}(A, \eta)=0) .
\end{aligned}
$$

That is, the measurement error, if uncorrelated with the true value, is necessarily correlated with the measured value. Using the fact that

$$
\operatorname{Cov}(S, \eta)=0, \quad \operatorname{Cov}(I Q, \eta)=\operatorname{Var}(\eta), \quad \text { and } \quad \operatorname{Cov}(\mathbf{h}, \eta)=\mathbf{0}
$$

the projection coefficients can be rewritten as

$$
\left[\begin{array}{c}
\theta_{S} \\
\theta_{I Q} \\
\theta_{\mathbf{h}}
\end{array}\right]=\operatorname{Var}(\eta) \cdot \mathbf{a},
$$

where

$$
\mathbf{a} \equiv \text { second column of }\left[\begin{array}{ccc}
\operatorname{Var}(S) & \operatorname{Cov}(S, I Q) & \operatorname{Cov}\left(S, \mathbf{h}^{\prime}\right) \\
\operatorname{Cov}(I Q, S) & \operatorname{Var}(I Q) & \operatorname{Cov}\left(I Q, \mathbf{h}^{\prime}\right) \\
\operatorname{Cov}(\mathbf{h}, S) & \operatorname{Cov}(\mathbf{h}, I Q) & \operatorname{Var}(\mathbf{h})
\end{array}\right]^{-1} .
$$

Therefore, in the regression of $L W$ on a constant, $S, I Q$, and $\mathbf{h}$,


\begin{align*}
& \operatorname{plim} \widehat{\beta}_{\mathrm{OLS}}=\beta-\gamma \cdot \operatorname{Var}(\eta) \cdot(1 \text { st element of } \mathbf{a})  \tag{3.9.7a}\\
& \operatorname{plim} \hat{\gamma}_{\mathrm{OLS}}=\gamma-\gamma \cdot \operatorname{Var}(\eta) \cdot(2 \text { nd element of } \mathbf{a}) \tag{3.9.7b}
\end{align*}


Since the second element of $\mathbf{a}$ is positive (it is a diagonal element of the inverse of a variance-covariance matrix), the OLS estimate of the ability coefficient is biased downward. The direction of the asymptotic bias for the schooling coefficient, however, depends on the sign of the first element of $\mathbf{a}$. Typically, $\operatorname{Cov}(S, I Q)$ would be positive, so, barring unusually strong correlation of ( $S, I Q$ ) with $\mathbf{h}$, the first element of a would be negative. Thus, we conjecture that $\widehat{\beta}_{\text {OLS }}$ is biased upward for schooling.

\section*{2SLS to Correct for the Bias}
To control for the bias, Griliches applies the 2SLS to the wage equation. To do so, the set of instruments needs to be specified. The predetermined regressors, $(1, S, \mathbf{h})$, can be included in the set. The additional variables included in the set to instrument $I Q$ are age, age squared, $K W W$, mother's education, and some other background variables of the individual (such as father's occupation). Those variables are thus assumed to be predetermined. Griliches's 2SLS estimate for this specification is reproduced in line 3 of Table 3.2. ${ }^{20}$ In accordance with our prediction that the OLS estimate of the ability coefficient is biased downward, the OLS estimate of the ability coefficient of 0.0019 in line 2 is lower than the 2SLS estimate of 0.0038 in line 3 . Our conjecture that $\widehat{\beta}_{\text {OLS }}$ of the schooling coefficient when $I Q$ is included in the regression is biased upward, too, is borne out by data because the estimate of $\beta$ of 0.059 in line 2 is higher than the estimate of 0.052 in line 3.

To summarize, the "ability bias" studied by Griliches is that the schooling coefficient is biased upward if ability is ignored and is biased in an unknown direction (but perhaps upward) if an imperfect measure of ability ( $I Q$ in the present case) is included. The 2SLS provides a solution, but it is predicated on the assumption that the instruments used (such as $M E D$ and $K W W$ ) are uncorrelated with unobserved wage determinants $\varepsilon$ and $\eta$. Those determinants would include personal characteristics like diligence. It seems not unreasonable to suppose that $K W W$ depends on such characteristics. If so, $K W W$ cannot serve as an instrument. One could go further and argue that $M E D$ is not predetermined either, because some

\footnotetext{${ }^{20}$ The $R^{2}$ is not reported for the 2SLS estimate because, unlike in the OLS case, the sum of squared 2SLS residuals cannot be divided between the "explained variation" and the "unexplained variation."
}
of the mother's personal characteristics that influenced her education and that are valued in the marketplace would be inherited by the child. Whether those variables included in the set of instruments are predetermined can be tested by Sargan's statistic. But the statistic is not reported in Griliches's paper (which was written long before doing so became a standard practice).

\section*{Subsequent Developments}
Perhaps because it is so difficult to come up with valid instruments that are uncorrelated with unobserved characteristics but correlated with $I Q$, controlling for the "ability bias" continued to attract much attention in the literature. There is a fairly large literature starting with Behrman et al. (1980), which compares identical twins with different levels of education to control for genetic characteristics and family background. One can also compare the same individual at two points in time. In either case, the appropriate estimation technique is what is called the fixed-effect estimator, to be covered in Chapter 5.

The endogeneity of schooling is another major issue. If one takes the view that the error term includes a host of unobservable individual characteristics that might affect the individual's choice of schooling, then schooling needs to be treated as an endogenous variable. But again, finding valid instruments uncorrelated with unobservable characteristics but correlated with schooling is extremely difficult. The recent literature can be viewed as a search for the determinants of schooling having little to do with the individual's unobserved characteristics. The variable used in Angrist and Krueger (1991), for example, is whether the individual is subject to compulsory schooling laws.

Finally, the literature on the relationship between school quality and earnings is attracting renewed attention. See Card and Krueger (1996) for a survey.

\section*{QUESTIONS FOR REVIEW}
\begin{enumerate}
  \item List all the assumptions made in this section about the means and covariances of ( $S, A, I Q, K W W, \mathbf{h}, \varepsilon, \eta$ ). Study how Griliches states those assumptions. Which ones are crucial for the 2SLS to be consistent?
  \item Suppose the relationship between $I Q$ and $A$ is given not by (3.9.3) but by
\end{enumerate}

$$
I Q=\phi+\pi \cdot A+\eta .
$$

Is the 2SLS estimator of $\beta$ (the schooling coefficient) still consistent? [Answer: Yes.]\\
3. In the Spencian model of job market signaling, individuals are sorted into two groups, with low-ability individuals choosing a low level of education and high-ability ones choosing a high level of education. Education does not increase the earning capacity of individuals, so the individual's wage rate is determined by his or her ability. If we had data on the wage rate, ability, and schooling for a sample of individuals that includes both high- and low-ability individuals, can we test the hypothesis that education by itself does not contribute to higher wages? Hint: The data would have a multicollinearity problem.

\section*{PROBLEM SET FOR CHAPTER 3}
$\_\_\_\_$

\section*{ANALYTICAL EXERCISES}
\begin{enumerate}
  \item Prove: A symmetric and idempotent matrix is positive semidefinite.
  \item We wish to prove Proposition 3.4. To avoid inessential complications, consider the case where $K=1$ and $L=1$ so that $\mathbf{x}_{i}$ and $\mathbf{z}_{i}$ are scalars. So (3.5.5) becomes
\end{enumerate}

$$
\hat{\varepsilon}_{i}^{2}=\varepsilon_{i}^{2}-2(\hat{\delta}-\delta) z_{i} \varepsilon_{i}+(\hat{\delta}-\delta)^{2} z_{i}^{2}
$$

and the $\widehat{\mathbf{S}}$ in (3.5.10) becomes


\begin{equation*}
\frac{1}{n} \sum_{i=1}^{n} x_{i}^{2} \varepsilon_{i}^{2}-2(\hat{\delta}-\delta) \frac{1}{n} \sum_{i=1}^{n} z_{i} x_{i}^{2} \varepsilon_{i}+(\hat{\delta}-\delta)^{2} \frac{1}{n} \sum_{i=1}^{n} z_{i}^{2} x_{i}^{2} . \tag{*}
\end{equation*}


(a) Show: The first term on the RHS of (*) converges in probability to $S \left(\equiv \mathrm{E}\left(x_{i}^{2} \varepsilon_{i}^{2}\right)\right)$.\\
(b) Use the Cauchy-Schwartz inequality

$$
\mathrm{E}(|X \cdot Y|) \leq \sqrt{\mathrm{E}\left(X^{2}\right) \cdot \mathrm{E}\left(Y^{2}\right)}
$$

to show that $\mathrm{E}\left(z_{i} x_{i}^{2} \varepsilon_{i}\right)$ exists and is finite. Hint: $z_{i} x_{i}^{2} \varepsilon_{i}$ is the product of $x_{i} \varepsilon_{i}$ and $x_{i} z_{i}$. Fact: $\mathrm{E}(x)$ exists and is finite if and only if $\mathrm{E}(|x|)<\infty$. Where do we use Assumption 3.6?\\
(c) Show that the second term on the RHS of (\textit{) converges to zero in probability.\\
(d) In a similar fashion, show that the third term on the RHS of (}) converges to zero in probability.\\
3. We wish to prove the algebraic result (3.5.11). What needs to be proved is that $\mathbf{A}-\mathbf{B}$ is positive semidefinite where

$$
\mathbf{A} \equiv\left(\boldsymbol{\Sigma}_{\mathbf{x z}}^{\prime} \mathbf{W} \boldsymbol{\Sigma}_{\mathbf{x z}}\right)^{-1} \boldsymbol{\Sigma}_{\mathbf{x z}}^{\prime} \mathbf{W S W} \boldsymbol{\Sigma}_{\mathbf{x z}}\left(\boldsymbol{\Sigma}_{\mathbf{x z}}^{\prime} \mathbf{W} \boldsymbol{\Sigma}_{\mathbf{x z}}\right)^{-1}
$$

and

$$
\mathbf{B} \equiv\left(\boldsymbol{\Sigma}_{\mathbf{x z}}^{\prime} \mathbf{S}^{-1} \boldsymbol{\Sigma}_{\mathbf{x z}}\right)^{-1} .
$$

A result from matrix algebra states:\\
Let $\mathbf{A}$ and $\mathbf{B}$ be two positive definite and symmetric matrices. $\mathbf{A}-\mathbf{B}$ is positive semidefinite if and only if $\mathbf{B}^{-1}-\mathbf{A}^{-1}$ is positive semidefinite.

Therefore, what needs to be proved is that

$$
\mathbf{Q} \equiv \boldsymbol{\Sigma}_{\mathbf{x z}}^{\prime} \mathbf{S}^{-1} \boldsymbol{\Sigma}_{\mathbf{x z}}-\boldsymbol{\Sigma}_{\mathbf{x z}}^{\prime} \mathbf{W} \boldsymbol{\Sigma}_{\mathbf{x z}}\left(\boldsymbol{\Sigma}_{\mathbf{x z}}^{\prime} \mathbf{W S W} \boldsymbol{\Sigma}_{\mathbf{x z}}\right)^{-1} \boldsymbol{\Sigma}_{\mathbf{x z}}^{\prime} \mathbf{W} \boldsymbol{\Sigma}_{\mathbf{x z}}
$$

is positive semidefinite.\\
(a) Since $\mathbf{S}(K \times K)$ is positive definite, there exists a nonsingular $K \times K$ matrix $\mathbf{C}$ such that $\mathbf{C}^{\prime} \mathbf{C}=\mathbf{S}^{-1}$. Verify that $\mathbf{Q}$ can be written as

$$
\mathbf{Q}=\mathbf{H}^{\prime} \mathbf{M}_{\mathbf{G}} \mathbf{H},
$$

where

$$
\mathbf{H}=\mathbf{C} \boldsymbol{\Sigma}_{\mathbf{x z}}, \quad \mathbf{M}_{\mathbf{G}}=\mathbf{I}_{K}-\mathbf{G}\left(\mathbf{G}^{\prime} \mathbf{G}\right)^{-1} \mathbf{G}^{\prime}, \quad \mathbf{G}=\mathbf{C}^{\prime-1} \mathbf{W} \boldsymbol{\Sigma}_{\mathbf{x z}} .
$$

Hint: $\mathbf{C}^{-1} \mathbf{C}^{-1}=\mathbf{S}$.\\
(b) Show that $\mathbf{Q}$ is positive semidefinite. Hint: First show that $\mathbf{M}_{\mathbf{G}}$ is symmetric and idempotent. As you showed in Exercise 1, a symmetric and idempotent matrix is positive semidefinite.\\
4. Suppose $\mathbf{z}_{i}$ is a strict subset of $\mathbf{x}_{i}$. Which one is smaller in the matrix sense, the asymptotic variance of the OLS estimator or the asymptotic variance of the efficient two-step GMM estimator? Hint: Without loss of generality, we can assume that the first $L$ elements of the $K$-dimensional vector $\mathbf{x}_{i}$ are $\mathbf{z}_{i}$. If $K \times K$ matrix $\mathbf{W}$ is given by

$$
\mathbf{W}=\left[\begin{array}{cc}
\mathbf{A}^{-1} & \mathbf{0} \\
\mathbf{0} & \mathbf{0}
\end{array}\right]
$$

where

$$
\underset{(L \times L)}{\mathbf{A}} \equiv \mathrm{E}\left(\varepsilon_{i}^{2} \mathbf{z}_{i} \mathbf{z}_{i}^{\prime}\right)=\text { leading }(L \times L) \text { submatrix of } \mathbf{S},
$$

then the asymptotic variance of the OLS estimator can be written as (3.5.1). If you did Exercise 3, you can verify that the algebraic relation (3.5.11) does not require $\mathbf{W}$ to be positive definite, as long as $\boldsymbol{\Sigma}_{\mathbf{x z}}^{\prime} \mathbf{W S W} \boldsymbol{\Sigma}_{\mathbf{x z}}$ is nonsingular.\\
5. (optional) Prove Proposition 3.6 by following the steps below.\\
(a) Show: $\mathbf{g}_{n}\left(\hat{\delta}\left(\widehat{\mathbf{S}}^{-1}\right)\right)=\widehat{\mathbf{B}} \overline{\mathbf{g}}$ where $\widehat{\mathbf{B}} \equiv \mathbf{I}_{K}-\mathbf{S}_{\mathbf{x z}}\left(\mathbf{S}_{\mathbf{x z}}^{\prime} \widehat{\mathbf{S}}^{-1} \mathbf{S}_{\mathbf{x z}}\right)^{-1} \mathbf{S}_{\mathbf{x z}}^{\prime} \widehat{\mathbf{S}}^{-1}$.\\
(b) Since $\widehat{\mathbf{S}}$ is positive definite, there exists a nonsingular $K \times K$ matrix $\mathbf{C}$ such that $\mathbf{C}^{\prime} \mathbf{C}=\widehat{\mathbf{S}}^{-1}$. Define $\mathbf{A} \equiv \mathbf{C S}_{\mathbf{x z}}$. Show that $\widehat{\mathbf{B}}^{\prime} \widehat{\mathbf{S}}^{-1} \widehat{\mathbf{B}}=\mathbf{C}^{\prime} \mathbf{M C}$, where $\mathbf{M} \equiv \mathbf{I}_{K}-\mathbf{A}\left(\mathbf{A}^{\prime} \mathbf{A}\right)^{-1} \mathbf{A}^{\prime}$. What is the rank of $\mathbf{M}$ ? Hint: The rank of an idempotent matrix equals its trace.\\
(c) Let $\mathbf{v} \equiv \sqrt{n} \mathbf{C} \overline{\mathbf{g}}$. Show that $\mathbf{v} \rightarrow_{\mathrm{d}} N\left(\mathbf{0}, \mathbf{I}_{K}\right)$.\\
(d) Show that $J\left(\hat{\delta}\left(\widehat{\mathbf{S}}^{-1}\right), \widehat{\mathbf{S}}^{-1}\right)=\mathbf{v}^{\prime} \mathbf{M v}$. What is its asymptotic distribution?\\
6. Show that the $J$ statistic in Proposition 3.6 is numerically equal to

$$
n \cdot \mathbf{s}_{\mathbf{x y}}^{\prime} \widehat{\mathbf{B}} \widehat{\mathbf{S}}^{-1} \widehat{\mathbf{B}} \mathbf{s}_{\mathbf{x y}}
$$

where $\widehat{\mathbf{B}}$ is defined in Exercise 5(a). Hint: From Exercise 5, $J=n \cdot \overline{\mathbf{g}}^{\prime} \widehat{\mathbf{B}}^{\prime} \widehat{\mathbf{S}}^{-1} \widehat{\mathbf{B}} \overline{\mathbf{g}}$. Show that $\widehat{\mathbf{B}} \overline{\mathbf{g}}=\widehat{\mathbf{B}} \mathbf{s}_{\mathbf{x y}}$.\\
7. (optional) Prove Proposition 3.7 by taking the following steps. The notation not introduced in this exercise is the same as in Exercise 5.\\
(a) Prove:

$$
\mathbf{g}_{1 n}(\overline{\boldsymbol{\delta}})=\widehat{\mathbf{B}}_{1} \overline{\mathbf{g}}_{1} \quad \text { and } \quad \widehat{\mathbf{B}}_{1}^{\prime}\left(\widehat{\mathbf{S}}_{11}\right)^{-1} \widehat{\mathbf{B}}_{1}=\mathbf{C}_{1}^{\prime} \mathbf{M}_{1} \mathbf{C}_{1},
$$

where

$$
\begin{aligned}
\overline{\mathbf{g}}_{1} & \equiv \frac{1}{n} \sum_{i=1}^{n} \mathbf{x}_{i 1} \cdot \varepsilon_{i}, \\
\widehat{\mathbf{B}}_{1} & \equiv \mathbf{I}_{K_{1}}-\mathbf{S}_{\mathbf{x}_{1} \mathbf{z}}\left[\mathbf{S}_{\mathbf{x}_{1} \mathbf{z}}^{\prime}\left(\widehat{\mathbf{S}}_{11}\right)^{-1} \mathbf{S}_{\mathbf{x}_{1} \mathbf{z}}\right]^{-1} \mathbf{S}_{\mathbf{x}_{1} \mathbf{z}}^{\prime}\left(\widehat{\mathbf{S}}_{11}\right)^{-1}, \\
\mathbf{C}_{1}^{\prime} \mathbf{C}_{1} & \equiv\left(\widehat{\mathbf{S}}_{11}\right)^{-1},
\end{aligned}
$$

$$
\begin{aligned}
\mathbf{M}_{1} & \equiv \mathbf{I}_{K_{1}}-\mathbf{A}_{1}\left(\mathbf{A}_{1}^{\prime} \mathbf{A}_{1}\right)^{-1} \mathbf{A}_{1}^{\prime} \\
\mathbf{A}_{1} & \equiv \mathbf{C}_{1} \mathbf{S}_{\mathbf{x}_{1} \mathbf{z}}
\end{aligned}
$$

(b) Prove: $\operatorname{rank}(\mathbf{M})=K-L, \operatorname{rank}\left(\mathbf{M}_{1}\right)=K_{1}-L$.\\
(c) Prove: You have proved in the previous exercise that $J=\mathbf{v}^{\prime} \mathbf{M v}$, where $\mathbf{v} \equiv \sqrt{n} \mathbf{C} \overline{\mathbf{g}}$. Here, prove: $J_{1}=\mathbf{v}_{1}^{\prime} \mathbf{M}_{1} \mathbf{v}_{1}$, where $\mathbf{v}_{1} \equiv \sqrt{n} \mathbf{C}_{1} \overline{\mathbf{g}}_{1}$.

To proceed further, define the $K \times K_{1}$ matrix $\mathbf{F}$ as

$$
\left.\underset{\left(K \times K_{1}\right)}{\mathbf{F}}=\left[\begin{array}{c}
\mathbf{I}_{K_{1}} \\
\mathbf{0}
\end{array}\right]\right\} \begin{aligned}
& K_{1} \text { rows } \\
& K-K_{1} \text { rows. }
\end{aligned}
$$

Then: $\mathbf{x}_{i 1}=\mathbf{F}^{\prime} \mathbf{x}_{i}, \mathbf{S}_{\mathbf{x}_{1} \mathbf{z}}=\mathbf{F}^{\prime} \mathbf{S}_{\mathbf{x z}}, \overline{\mathbf{g}}_{1}=\mathbf{F}^{\prime} \overline{\mathbf{g}}$.\\
(d) Prove $J-J_{1}=\mathbf{v}^{\prime}(\mathbf{M}-\mathbf{D}) \mathbf{v}$ where $\mathbf{D}=\mathbf{C}^{\prime-1} \mathbf{F C}_{1}^{\prime} \mathbf{M}_{1} \mathbf{C}_{1} \mathbf{F}^{\prime} \mathbf{C}^{-1}$.\\
(e) Prove that D is symmetric and idempotent, with rank $K_{1}-L$. Hint:

$$
\begin{gathered}
\mathbf{F}^{\prime} \mathbf{C}^{-1} \mathbf{C}^{\prime-1} \mathbf{F}=\mathbf{F}^{\prime}\left(\mathbf{C}^{\prime} \mathbf{C}\right)^{-1} \mathbf{F}=\mathbf{F}^{\prime} \widehat{\mathbf{S}} \mathbf{F}=\widehat{\mathbf{S}}_{11} \\
\mathbf{C}_{1} \widehat{\mathbf{S}}_{11} \mathbf{C}_{1}^{\prime}=\mathbf{C}_{1}\left(\mathbf{C}_{1}^{\prime} \mathbf{C}_{1}\right)^{-1} \mathbf{C}_{1}^{\prime}=\mathbf{I}_{K_{1}}
\end{gathered}
$$

(f) Prove: A'D = 0. Hint:

$$
\begin{gathered}
\mathbf{A}^{\prime} \mathbf{D}=\left(\mathbf{S}_{\mathbf{x z}}^{\prime} \mathbf{C}^{\prime}\right)\left(\mathbf{C}^{\prime-1} \mathbf{F C}_{1}^{\prime} \mathbf{M}_{1} \mathbf{C}_{1} \mathbf{F}^{\prime} \mathbf{C}^{-1}\right)=\left(\mathbf{S}_{\mathbf{x z}}^{\prime} \mathbf{F C}_{1}^{\prime} \mathbf{M}_{1}\right)\left(\mathbf{C}_{1} \mathbf{F}^{\prime} \mathbf{C}^{-1}\right) \\
\mathbf{S}_{\mathbf{x z}}^{\prime} \mathbf{F C}_{1}^{\prime}=\mathbf{A}_{1}^{\prime}
\end{gathered}
$$

so

$$
\mathbf{S}_{\mathbf{x z}}^{\prime} \mathbf{F C}_{1}^{\prime} \mathbf{M}_{1}=\mathbf{A}_{1}^{\prime} \mathbf{M}_{1}
$$

(g) Prove: $\mathbf{M}-\mathbf{D}$ is symmetric and idempotent, with rank $K-K_{1}$.\\
(h) Prove the desired result: $J-J_{1} \rightarrow_{\mathrm{d}} \chi^{2}\left(K-K_{1}\right)$.\\
(i) Step (d) has established that $C=n \cdot \overline{\mathbf{g}}^{\prime} \mathbf{C}^{\prime}(\mathbf{M}-\mathbf{D}) \mathbf{C} \overline{\mathbf{g}}$. Show that $C$ can be written as

$$
n \cdot \mathbf{s}_{\mathbf{x y}}^{\prime} \mathbf{C}^{\prime}(\mathbf{M}-\mathbf{D}) \mathbf{C s}_{\mathbf{x y}}
$$

Hint: $\mathbf{g}_{n}(\hat{\boldsymbol{\delta}})=\widehat{\mathbf{B}} \mathbf{s}_{\mathbf{x y}}$. So the numerical value of $C$ would have been the same if we replaced $\overline{\mathbf{g}}$ by $\mathbf{s}_{\mathbf{x y}}$ from the beginning in (a).\\
(j) Show that $\mathbf{M}-\mathbf{D}$ is positive semidefinite. Thus, the quadratic form $C$ is nonnegative for any sample.\\
8. (Optional, proof of numerical equivalence in Proposition 3.8) In the restricted efficient GMM, the $J$ function is minimized subject to the restriction $\mathbf{R} \tilde{\boldsymbol{\delta}}=\mathbf{r}$. For the time being, do not restrict $\widehat{\mathbf{W}}$ to be equal to $\widehat{\mathbf{S}}^{-1}$ and form the Lagrangian as

$$
\mathcal{L}=n \cdot\left(\mathbf{S}_{\mathbf{x y}}-\mathbf{S}_{\mathbf{x z}} \tilde{\boldsymbol{\delta}}\right)^{\prime} \widehat{\mathbf{W}}\left(\mathbf{s}_{\mathbf{x y}}-\mathbf{S}_{\mathbf{x z}} \tilde{\boldsymbol{\delta}}\right)+\lambda^{\prime}(\mathbf{R} \tilde{\boldsymbol{\delta}}-\mathbf{r}),
$$

where $\lambda$ is the \#r-dimensional Lagrange multipliers (recall: $\mathbf{r}$ is \# $\mathbf{r} \times 1, \mathbf{R}$ is $\# \mathbf{r} \times L$, and $\tilde{\boldsymbol{\delta}}$ is $L \times 1$ ). Let $\overline{\boldsymbol{\delta}}$ be the solution to the constrained minimization problem. (If $\widehat{\mathbf{W}}=\widehat{\mathbf{S}}^{-1}$, it is the restricted efficient GMM estimator of $\boldsymbol{\delta}$.)\\
(a) Let $\hat{\delta}(\widehat{\mathbf{W}})$ be the unrestricted GMM estimator that minimizes $J(\tilde{\delta}, \widehat{\mathbf{W}})$.

Show:

$$
\begin{gathered}
\bar{\delta}=\hat{\delta}(\widehat{\mathbf{W}})-\left(\mathbf{S}_{\mathbf{x z}}^{\prime} \widehat{\mathbf{W}} \mathbf{S}_{\mathbf{x z}}\right)^{-1} \mathbf{R}^{\prime}\left[\mathbf{R}\left(\mathbf{S}_{\mathbf{x z}}^{\prime} \widehat{\mathbf{W}} \mathbf{S}_{\mathbf{x z}}\right)^{-1} \mathbf{R}^{\prime}\right]^{-1}[\mathbf{R} \hat{\delta}(\widehat{\mathbf{W}})-\mathbf{r}], \\
\lambda=2 n \cdot\left[\mathbf{R}\left(\mathbf{S}_{\mathbf{x z}}^{\prime} \widehat{\mathbf{W}} \mathbf{S}_{\mathbf{x z}}\right)^{-1} \mathbf{R}^{\prime}\right]^{-1}[\mathbf{R} \hat{\delta}(\widehat{\mathbf{W}})-\mathbf{r}] .
\end{gathered}
$$

Hint: The first-order conditions are

$$
-2 n \mathbf{S}_{\mathbf{x z}}^{\prime} \widehat{\mathbf{W}} \mathbf{s}_{\mathbf{x y}}+2 n\left(\mathbf{S}_{\mathbf{x z}}^{\prime} \widehat{\mathbf{W}} \mathbf{S}_{\mathbf{x z}}\right) \overline{\boldsymbol{\delta}}+\mathbf{R}^{\prime} \boldsymbol{\lambda}=\mathbf{0}
$$

Combine this with the constraint $\mathbf{R} \overline{\boldsymbol{\delta}}=\mathbf{r}$ to solve for $\lambda$ and $\overline{\boldsymbol{\delta}}$.\\
(b) Show:

$$
J(\overline{\boldsymbol{\delta}}, \widehat{\mathbf{W}})=J(\hat{\boldsymbol{\delta}}(\widehat{\mathbf{W}}), \widehat{\mathbf{W}})+n \cdot(\hat{\boldsymbol{\delta}}(\widehat{\mathbf{W}})-\overline{\boldsymbol{\delta}})^{\prime}\left(\mathbf{S}_{\mathbf{x z}}^{\prime} \widehat{\mathbf{W}} \mathbf{S}_{\mathbf{x z}}\right)(\hat{\boldsymbol{\delta}}(\widehat{\mathbf{W}})-\overline{\boldsymbol{\delta}}) .
$$

Hint:

$$
\begin{aligned}
& J(\overline{\boldsymbol{\delta}}, \widehat{\mathbf{W}})=n \cdot\left(\mathbf{S}_{\mathbf{x y}}-\mathbf{S}_{\mathbf{x z}} \overline{\boldsymbol{\delta}}\right)^{\prime} \widehat{\mathbf{W}}\left(\mathbf{S}_{\mathbf{x y}}-\mathbf{S}_{\mathbf{x z}} \overline{\boldsymbol{\delta}}\right) \\
&=n \cdot\left[\left(\mathbf{s}_{\mathbf{x y}}-\mathbf{S}_{\mathbf{x z}} \hat{\boldsymbol{\delta}}(\widehat{\mathbf{W}})\right)+\mathbf{S}_{\mathbf{x z}}(\hat{\boldsymbol{\delta}}(\widehat{\mathbf{W}})-\overline{\boldsymbol{\delta}})\right]^{\prime} \widehat{\mathbf{W}} \\
& \cdot\left[\left(\mathbf{s}_{\mathbf{x y}}-\mathbf{S}_{\mathbf{x z}} \hat{\boldsymbol{\delta}}(\widehat{\mathbf{W}})\right)+\mathbf{S}_{\mathbf{x z}}(\hat{\boldsymbol{\delta}}(\widehat{\mathbf{W}})-\overline{\boldsymbol{\delta}})\right]
\end{aligned}
$$

Use the first-order conditions for the unrestricted GMM:

$$
\mathbf{S}_{\mathbf{x z}}^{\prime} \widehat{\mathbf{W}}\left(\mathbf{s}_{\mathbf{x y}}-\mathbf{S}_{\mathbf{x z}} \hat{\delta}(\widehat{\mathbf{W}})\right)=\mathbf{0} .
$$

(c) Now set $\widehat{\mathbf{W}}=\widehat{\mathbf{S}}^{-1}$ and show that the Wald statistic $W$ defined in (3.5.8) is numerically equal to $L R$ in (3.7.2).\\
9. (GMM Hausman principle) Let $\hat{\boldsymbol{\delta}}_{1} \equiv \hat{\boldsymbol{\delta}}\left(\widehat{\mathbf{W}}_{1}\right)$ and $\hat{\boldsymbol{\delta}}_{2} \equiv \hat{\boldsymbol{\delta}}\left(\widehat{\mathbf{W}}_{2}\right)$ be two GMM estimators with two different choices of the weighting matrix, $\widehat{\mathbf{W}}_{1}$ and $\widehat{\mathbf{W}}_{2}$.\\
(a) Show:

$$
\left[\begin{array}{c}
\sqrt{n}\left(\hat{\boldsymbol{\delta}}_{1}-\boldsymbol{\delta}\right) \\
\sqrt{n}\left(\hat{\boldsymbol{\delta}}_{2}-\boldsymbol{\delta}\right)
\end{array}\right] \rightarrow N\left(\mathbf{0},\left[\begin{array}{ll}
\mathbf{A}_{11} & \mathbf{A}_{12} \\
\mathbf{A}_{21} & \mathbf{A}_{22}
\end{array}\right]\right)
$$

where

$$
\begin{aligned}
& \mathbf{A}_{11}=\mathbf{Q}_{1}^{-1} \boldsymbol{\Sigma}_{\mathbf{x z}}^{\prime} \mathbf{W}_{1} \mathbf{S} \mathbf{W}_{1} \boldsymbol{\Sigma}_{\mathbf{x z}} \mathbf{Q}_{1}^{-1}, \quad \mathbf{A}_{12}=\mathbf{Q}_{1}^{-1} \boldsymbol{\Sigma}_{\mathbf{x z}}^{\prime} \mathbf{W}_{1} \mathbf{S} \mathbf{W}_{2} \boldsymbol{\Sigma}_{\mathbf{x z}} \mathbf{Q}_{2}^{-1} \\
& \mathbf{A}_{21}=\mathbf{Q}_{2}^{-1} \boldsymbol{\Sigma}_{\mathbf{x z}}^{\prime} \mathbf{W}_{2} \mathbf{S} \mathbf{W}_{1} \boldsymbol{\Sigma}_{\mathbf{x z}} \mathbf{Q}_{1}^{-1}, \quad \mathbf{A}_{22}=\mathbf{Q}_{2}^{-1} \boldsymbol{\Sigma}_{\mathbf{x z}}^{\prime} \mathbf{W}_{2} \mathbf{S} \mathbf{W}_{2} \boldsymbol{\Sigma}_{\mathbf{x z}} \mathbf{Q}_{2}^{-1} \\
& \mathbf{Q}_{\mathbf{1}}=\mathbf{\Sigma}_{\mathbf{x z}}^{\prime} \mathbf{W}_{1} \boldsymbol{\Sigma}_{\mathbf{x z}}, \quad \mathbf{Q}_{2}=\boldsymbol{\Sigma}_{\mathbf{x z}}^{\prime} \mathbf{W}_{2} \boldsymbol{\Sigma}_{\mathbf{x z}}, \quad \operatorname{plim}_{n \rightarrow \infty} \widehat{\mathbf{W}}_{j}=\mathbf{W}_{j} \quad(j=1,2)
\end{aligned}
$$

(b) Let $\mathbf{q} \equiv \hat{\boldsymbol{\delta}}_{1}-\hat{\boldsymbol{\delta}}_{2}$. Show that $\sqrt{n} \mathbf{q} \rightarrow_{\mathrm{d}} N(\mathbf{0}, \operatorname{Avar}(\mathbf{q}))$, where

$$
\operatorname{Avar}(\mathbf{q})=\mathbf{A}_{11}+\mathbf{A}_{22}-\mathbf{A}_{12}-\mathbf{A}_{21} .
$$

(c) Set $\widehat{\mathbf{W}}_{2}=\widehat{\mathbf{S}}^{-1}$ so that $\hat{\boldsymbol{\delta}}_{2}$ is efficient GMM. Show that

$$
\operatorname{Avar}(\mathbf{q})=\mathbf{A}_{11}-\left(\boldsymbol{\Sigma}_{\mathbf{x z}}^{\prime} \mathbf{S}^{-1} \boldsymbol{\Sigma}_{\mathbf{x z}}\right)^{-1}=\operatorname{Avar}\left(\hat{\delta}\left(\widehat{\mathbf{W}}_{1}\right)\right)-\operatorname{Avar}\left(\hat{\delta}\left(\widehat{\mathbf{S}}^{-1}\right)\right) .
$$

\begin{enumerate}
  \setcounter{enumi}{9}
  \item (Weak instruments. This question is due to M. Watson.) Consider the following model:
\end{enumerate}

$$
y_{i}=z_{i} \delta+\varepsilon_{i} \quad(i=1,2, \ldots, n),
$$

where $y_{i}, z_{i}$, and $\varepsilon_{i}$ are scalar random variables. Let $x_{i}$ denote a scalar instrumental variable that is related to $z_{i}$ via the equation:

$$
z_{i}=x_{i} \beta+v_{i}
$$

Let

$$
\boldsymbol{\eta}_{i} \equiv\left[\begin{array}{c}
\varepsilon_{i} \\
v_{i}
\end{array}\right], \quad \mathbf{g}_{i} \equiv \boldsymbol{\eta}_{i} \cdot x_{i} \equiv\left[\begin{array}{c}
g_{1 i} \\
g_{2 i}
\end{array}\right]=\left[\begin{array}{c}
x_{i} \varepsilon_{i} \\
x_{i} v_{i}
\end{array}\right] .
$$

Assume that\\
(1) $\left\{x_{i}, \eta_{i}\right\}$ follows a stationary and ergodic process.\\
(2) $\mathbf{g}_{i}$ is a martingale difference sequence with $\mathrm{E}\left(\mathbf{g}_{i} \mathbf{g}_{i}^{\prime}\right)=\mathbf{S}$, which is a positive definite matrix.\\
(3) $\mathrm{E}\left(x_{i}^{2}\right)=\sigma_{x}^{2}>0$.\\
(4) $\beta \neq 0$.

Let $\hat{\delta} \equiv s_{x z}^{-1} s_{x y}$ where

$$
s_{x z} \equiv \frac{1}{n} \sum_{i=1}^{n} x_{i} z_{i}, \quad s_{x y} \equiv \frac{1}{n} \sum_{i=1}^{n} x_{i} y_{i}
$$

(a) Prove that $\sigma_{x z} \equiv \mathrm{E}\left(x_{i} z_{i}\right) \neq 0$ (so that the rank condition for identification is satisfied).\\
(b) Prove that $\hat{\delta} \rightarrow_{\mathrm{p}} \delta$.

We now want to consider a case in which the rank condition is just barely satisfied, that is, $\sigma_{x z}$ is very close to zero. One way to do this is to replace Assumption (4) with\\
${ }^{(4 \prime)} \beta=1 / \sqrt{n}$.\\
For the remainder of the problem, assume that (1)-(3) and (4') hold. (This assumption and the following questions are based on Staiger and Stock (1997).) Note that $\hat{\delta}-\delta=s_{x z}^{-1} \bar{g}_{1}$, where $\bar{g}_{1}$ is the sample mean of $g_{1 i}\left(\equiv x_{i} \varepsilon_{i}\right)$.\\
(c) Show that $s_{x z} \rightarrow_{\mathrm{p}} 0$.\\
(d) Show that $\sqrt{n} s_{x z} \rightarrow_{\mathrm{d}} \sigma_{x}^{2}+a$, where $a$ is distributed $N\left(0, s_{22}\right)$ and $s_{22}$ is the $(2,2)$ element of $\mathbf{S}$.\\
(e) Show that $\hat{\delta}-\delta \rightarrow_{\mathrm{d}}\left(\sigma_{x}^{2}+a\right)^{-1} b$, where ( $a, b$ ) are jointly normally distributed with mean zero and covariance matrix $\mathbf{S}$.\\
(f) Is $\hat{\delta}$ consistent?

\section*{EMPIRICAL EXERCISES}
Read Griliches (1976) before answering. We will estimate the type of the wage equation estimated by Griliches using an extract from the NLS-Y used by Blackburn and Neumark (1992). ${ }^{21}$ The NLS-Y is panel data, with the same set of young men surveyed at several points in time (we will not exploit the panel feature of the data set in this exercise, though). The extract contains information about those individuals at two points in time: first, the earliest year in which wages and other variables are available, and second, in 1980. In a data file GRILIC.ASC, data are provided on RNS, RNS80, MRT, MRT80, SMSA, SMSA80, MED, IQ, KWW, YEAR, AGE, AGE80, S, S80, EXPR, EXPR80, TENURE, TENURE80, LW, and LW80 (in this order, with the columns corresponding to the variables, as usual). The variable

\footnotetext{${ }^{21}$ See Blackburn and Neumark (1992, Section III) for a more detailed description of the sample.
}YEAR is the year of the first point in time. Variables without " 80 " are for the first point, and those with " 80 " are for 1980. The definition of the variables for the first point is:

$$
\begin{aligned}
R N S & =\text { dummy for residency in the southern states } \\
M R T & =\text { dummy for marital status (1 if married) } \\
S M S A & =\text { dummy for residency in metropolitan areas } \\
M E D & =\text { mother's education in years } \\
K W W & =\text { score on the "Knowledge of the World of Work" test } \\
I Q & =\text { IQ score } \\
A G E & =\text { age of the individual } \\
S & =\text { completed years of schooling } \\
E X P R & =\text { experience in years } \\
T E N U R E & =\text { tenure in years } \\
L W & =\log \text { wage. }
\end{aligned}
$$

The Blackburn-Neumark sample has 815 observations after deleting black individuals. Also deleted are cases with missing information on mother's education, reducing the sample size to 758 .\\
(a) Calculate means and standard deviations of all the provided variables (including those for 1980) and prepare a table similar to Griliches's Table 1. Also, calculate the correlation between $I Q$ and $S$ (use the MSD (CORR) command for TSP, CMOM (PRINT, CORR) for RATS).

Since the year the wage rate is observed differs across individuals, the wage rate will have the year effect. Generate eight year dummies for $Y E A R=66, \ldots, 73$. (Note: There is no observation for 1972.) The year dummies will be included in the log wage equation to control for the year effect.\\
(b) Consider the wage equation (dropping the individual subscript $i$ )

$$
L W=\beta S+\gamma I Q+\delta^{\prime} \mathbf{h}+\varepsilon,
$$

where $L W$ is log wages, $S$ is schooling, and

$$
\mathbf{h} \equiv(E X P R, T E N U R E, R N S, S M S A, \text { year dummies })^{\prime} .
$$

(There is no need to include a constant because it is a linear combination of\\
the year dummies.) Prepare a table similar to Table 3.2. In your 2SLS estimation of the wage equation, the set of instruments should consist of predetermined regressors ( $S$ and $\mathbf{h}$ ) and excluded predetermined variables ( $M E D$, $K W W, M R T, A G E)$. Is the relative magnitude of the three different estimates of the schooling coefficient consistent with the prediction based on the omitted variable and errors-in-variables biases?

TSP Tip: To do 2SLS, use the INST or 2SLS command.\\
RATS Tip: Use the LINREG command with the INST option. To define the instrument set, use the INSTRUMENTS command.\\
(c) For your 2SLS estimation of the wage equation in part (b), calculate Sargan's statistic (it should be 87.655 ). What should the degrees of freedom be? Calculate the $p$-value.

TSP Tip: Sargan's statistic can be calculated as @phi / (@ssr/@nob), where the variables with "@" are outputs from TSP commands. @nob is sample size, @phi is $\hat{\boldsymbol{\varepsilon}}^{\prime} \mathbf{X}\left(\mathbf{X}^{\prime} \mathbf{X}\right)^{-1} \mathbf{X}^{\prime} \hat{\boldsymbol{\varepsilon}}$ where $\hat{\boldsymbol{\varepsilon}}$ is the 2SLS residual vector, and @ssr is $\hat{\boldsymbol{\varepsilon}}^{\prime} \hat{\boldsymbol{\varepsilon}}$.

RATS Tip: Sargan's statistic is \%uzwzu / (\%rss / \%nobs), where variables with "\%" are outputs from RATS commands. \%nobs is sample size, \%uzwzu is $\hat{\boldsymbol{\varepsilon}}^{\prime} \mathbf{X}\left(\mathbf{X}^{\prime} \mathbf{X}\right)^{-1} \mathbf{X}^{\prime} \hat{\boldsymbol{\varepsilon}}$, and \%rss is $\hat{\boldsymbol{\varepsilon}}^{\prime} \hat{\boldsymbol{\varepsilon}}$.\\
(d) Obtain the 2SLS estimate by actually running two regressions. Verify that the standard errors given by the second stage regression are different from those you obtained in (b).\\
(e) Griliches mentions that schooling, too, may be endogenous. What is his argument? Estimate by 2SLS the wage equation, treating both $I Q$ and $S$ as endogenous. What happens to the schooling coefficient? How would you explain the difference between your 2SLS estimate of the schooling coefficient here and your 2SLS estimate in (b)? Calculate Sargan's statistic (it should be 13.268) and its $p$-value.\\
(f) Estimate the wage equation by GMM, treating schooling as predetermined as in the 2SLS estimation in part (b). Are 2SLS standard errors smaller than GMM standard errors? Test whether schooling is predetermined by the $C$ statistic. To calculate $C$, you need to do two GMMs with and without schooling as an instrument. Use the same $\widehat{\mathbf{S}}$ throughout. To calculate $\widehat{\mathbf{S}}$, use the 2SLS residual from part (b). The $C$ statistic should be 58.168.

TSP Tip: It is a bit awkward to calculate $\widehat{\mathbf{S}}$ in TSP. The following codes, though not elegant, get the job done. Here, res0 is the residuals from the 2SLS estimation in (b).

\begin{verbatim}
smpl 1 758;
x0=res0*s;
x1=res0*expr;
x2=res0* tenure;
x3=res0*rns;
x4=res0* smsa;
x5=res0*y66;
    \vdots
x11=res0*y73;
x12 =res0*med;
x13=res0*kww;
x14=res0*mrt;
x15=res0*age;
msd(noprint,moment) x0-x15;
\end{verbatim}

@mom, an output from the TSP command MSD, is $\widehat{\mathbf{S}}$. To do GMM, use TSP's GMM command. The GMM option COVOC=matrix name forces the GMM command to use the assigned matrix for $\widehat{\mathbf{S}}$. The $J$ statistic is given by @nob*@phi.

RATS Tip: The LINREG with the INST and ROBUSTERRORS option does GMM for linear equations. The WMATRIX=matrix name option of LINREG accepts $\widehat{\mathbf{S}}^{-1}$, not $\widehat{\mathbf{S}}$. To calculate $\widehat{\mathbf{S}}^{-1}$ in RATS, do the following. Here, res0 is the residuals from the 2SLS estimation in part (b).

\begin{verbatim}
mcov / res0
# s expr tenure rns smsa y66 ... y73
    med kww mrt age
compute sinv=inv(%cmom)
sinv is }\mp@subsup{\widehat{\mathbf{S}}}{}{-1}\mathrm{ . The J statistic is given by %uzwzu.
\end{verbatim}

(g) Go back to the wage equation you estimated in (e) by 2SLS endogenous schooling. The large Sargan's statistic (and the large $J$ statistic in part (f)) is a concern. So consider dropping both $M E D$ and $K W W$ from the list of instruments. Is the order condition for identification still satisfied? What happens to the

2SLS estimates? (The schooling coefficient will be $-529.3 \%$ !) What do you think is the problem? Hint: There are two endogenous regressors, $S$ and $I Q$. $M R T$ and $A G E$ are the only excluded predetermined variables. If $M R T$ and $A G E$ do not have explanatory power in the first-stage regression of $S$ and $I Q$ on the instruments, then the fitted values of $S$ and $I Q$ will be very close to linear combinations of the included predetermined variables. This will lead to a nearmulticollinearity problem in the second-stage regression. Make sure you examine the first-stage regressions to check the explanatory power of $M R T$ and $A G E$.

\section*{ANSWERS TO SELECTED QUESTIONS}
\section*{ANALYTICAL EXERCISES}
\begin{enumerate}
  \setcounter{enumi}{3}
  \item The asymptotic variance of the OLS estimator is (by setting $\mathbf{x}=\mathbf{z}$ in (3.5.1))
\end{enumerate}

$$
\boldsymbol{\Sigma}_{\mathbf{z z}}^{-1} \mathbf{A} \boldsymbol{\Sigma}_{\mathbf{z z}}^{-1}=\left(\boldsymbol{\Sigma}_{\mathbf{z z}} \mathbf{A}^{-1} \boldsymbol{\Sigma}_{\mathbf{z z}}\right)^{-1}
$$

where $\mathbf{A} \equiv \mathrm{E}\left(\varepsilon_{i}^{2} \mathbf{z}_{i} \mathbf{z}_{i}^{\prime}\right)$. The asymptotic variance of the two-step GMM is, by (3.5.13),

$$
\left(\boldsymbol{\Sigma}_{\mathbf{x z}}^{\prime} \mathbf{S}^{-1} \boldsymbol{\Sigma}_{\mathbf{x z}}\right)^{-1}
$$

where $\mathbf{S} \equiv \mathrm{E}\left(\varepsilon_{i}^{2} \mathbf{x}_{i} \mathbf{x}_{i}^{\prime}\right)$. If $\mathbf{W}$ is as defined in the hint, then

$$
\mathbf{W S W}=\mathbf{W} \quad \text { and } \quad \boldsymbol{\Sigma}_{\mathbf{x z}}^{\prime} \mathbf{W} \boldsymbol{\Sigma}_{\mathbf{x z}}=\boldsymbol{\Sigma}_{\mathbf{z z}} \mathbf{A}^{-1} \boldsymbol{\Sigma}_{\mathbf{z z}}
$$

So (3.5.1) reduces to the asymptotic variance of the OLS estimator. By (3.5.11), it is no smaller than $\left(\boldsymbol{\Sigma}_{\mathbf{x z}}^{\prime} \mathbf{S}^{-1} \boldsymbol{\Sigma}_{\mathbf{x z}}\right)^{-1}$, which is the asymptotic variance of the two-step GMM estimator.

\section*{EMPIRICAL EXERCISES}
(b) See Table 3.3, lines 1-5.\\
(e) He gives three reasons for the endogeneity of schooling. First, the measurement error $\eta$ may be related to success in schooling. The second reason is the argument we made in the text: the choice of $S$ is influenced by unobservable characteristics. Third, schooling may be measured with error. The general

\begin{table}[h]
\begin{center}
\captionsetup{labelformat=empty}
\caption{Table 3.3: Parameter Estimates, Blackburn-Neumark Sample\\
Estimation equation: dependent variable $=L W$, regressors: $S, I Q, E X P R, T E N U R E$, and other controls}
\begin{tabular}{|l|l|l|l|l|l|l|l|l|l|l|}
\hline
\multirow{2}{*}{Line no.} & \multirow{2}{*}{Estimation technique} & \multicolumn{4}{|c|}{Coefficient of} & \multirow{2}{*}{SEE} & \multirow{2}{*}{$R^{2}$} & \multirow{2}{*}{Test of overidentifying restrictions} & \multirow{2}{*}{Which regressor is endogenous?} & \multirow{2}{*}{Excluded predetermined variables besides regional dummies} \\
\hline
 &  & S & IQ & EXPR & TENURE &  &  &  &  &  \\
\hline
1 & OLS & 0.070 (0.0067) & - & 0.030 (0.0065) & 0.043 (0.0075) & 0.328 & 0.425 & - & none & - \\
\hline
2 & OLS & 0.062 (0.0073) & 0.0027 (0.0010) & 0.031 (0.0065) & 0.042 (0.0075) & 0.326 & 0.430 & - & none & - \\
\hline
3 & 2SLS & 0.069 (0.013) & 0.0002 (0.0039) & 0.030 (0.0066) & 0.043 (0.0076) & 0.328 & - & 87.6 ( $p=0.0000$ ) & IQ & MED, KWW, AGE, MRT \\
\hline
4 & 2SLS & 0.172 (0.021) & -0.009 (0.0047) & 0.049 (0.0082) & 0.042 (0.0088) & 0.380 & - & 13.3 ( $p=0.00131$ ) & S, IQ & MED, KWW, AGE, MRT \\
\hline
5 & GMM & 0.176 (0.021) & -0.009 (0.0049) & 0.050 (0.0080) & 0.043 (0.0095) & 0.379 & - & 11.6 ( $p=0.00303$ ) & S, IQ & MED, KWW, AGE, MRT \\
\hline
6 & 2SLS & 0.117 (0.027) & 0.002 (0.0050) & 0.033 (0.0052) & 0.0051 (0.0029) & 0.380 & - & 14.9 ( $p=0.00058$ ) & S80, IQ & MED, KWW, AGE80, MRT80 \\
\hline
\end{tabular}
\end{center}
\end{table}

Note: Standard errors in parentheses. Line 6 is for 1980.\\
formula for the asymptotic bias for the 2SLS estimator is given by

$$
\operatorname{plim} \hat{\boldsymbol{\delta}}_{2 \mathrm{SLS}}-\boldsymbol{\delta}=\left(\boldsymbol{\Sigma}_{\mathbf{x z}}^{\prime} \boldsymbol{\Sigma}_{\mathbf{x x}}^{-1} \boldsymbol{\Sigma}_{\mathbf{x z}}\right)^{-1} \boldsymbol{\Sigma}_{\mathbf{x z}}^{\prime} \boldsymbol{\Sigma}_{\mathbf{x x}}^{-1} \mathrm{E}\left(\mathbf{x}_{i} \cdot \varepsilon_{i}\right) .
$$

If $S$ is endogenous, the instrument set $\mathbf{x}_{i}$ for the 2SLS estimation in (b) erroneously includes a variable (which is $S$ ) that is not predetermined. So one of the elements of $\mathrm{E}\left(\mathbf{x}_{i} \cdot \varepsilon_{i}\right)$ is not zero. Then it is clear from the formula that asymptotic bias can exist at least for some coefficients.\\
(g) Without $M E D$ and $K W W$ as instruments, the 2SLS estimator becomes imprecise. The order condition is satisfied with equality. The culprit is the low explanatory power of $M R T$ (look at its $t$-value in the first-stage regressions for $S$ and $I Q$ ).

\section*{References}
Angrist, J., and A. Krueger, 1991, "Does Compulsory School Attendance Affect Schooling and Earnings?" Quarterly Journal of Economics, 106, 979-1014.\\
Behrman, J., Z. Hrubec, P. Taubman, and T. Wales, 1980, Socioeconomic Success: A Study of the Effects of Genetic Endowments, Family Environment, and Schooling, Amsterdam: North-Holland.\\
Blackburn, M., and D. Neumark, 1992, "Unobserved Ability, Efficiency Wages, and Interindustry Wage Differentials," Quarterly Journal of Economics, 107, 1421-1436.\\
Card, D., 1995, "Earnings, Schooling, and Ability Revisited," Research in Labor Economics, 14, 23-48.\\
Card, D., and A. Krueger, 1996, "Labor Market Effects of School Quality: Theory and Evidence," Working Paper 5450, NBER.\\
Eichenbaum, M., L. Hansen, and K. Singleton, 1985, "A Time Series Analysis of Representative Agent Models of Consumption and Leisure Choice under Uncertainty,"Quarterly Journal of Economics, 103, 51-78.\\
Friedman, M., 1957, A Theory of the Consumption Function, Princeton: Princeton University Press.\\
Gallant, R., and D. Jorgenson, 1979, "Statistical Inference for a System of Simultaneous, Non-linear, Implicit Equations in the Context of Instrumental Variable Estimation,"Journal of Econometrics, 11, 275-302.\\
Griliches, Z., 1976, "Wages of Very Young Men," Journal of Political Economy, 84, S69-S85.\\
Haavelmo, T., 1943, "The Statistical Implications of a System of Simultaneous Equations," Econometrica, 11, 1-12.\\
Hansen, L., 1982, "Large Sample Properties of Generalized Method of Moments Estimators," Econometrica, 50, 1029-1054.\\
Hausman, J., and W. Taylor, 1980, "Comparing Specification Tests and Classical Tests," MIT Working Paper No. 266.\\
Judge, G., W. Griffiths, R. Hill, H. Lütkepohl, and T. Lee, 1985, The Theory and Practice of Econometrics (2d ed.), New York: Wiley.\\
Mincer, J., 1958, "Investment in Human Capital and Personal Income Distribution," Journal of Political Economy, 66, 281-302.

Nelson, C., and R. Startz, 1990, "Some Further Results on the Exact Small Sample Properties of the Instrumental Variable Estimator," Econometrica, 58, 967-976.\\
Newey, W., 1985, "Generalized Method of Moments Specification Testing," Journal of Econometrics, 29, 229-256.\\
Newey, W., and D. McFadden, 1994, "Large Sample Estimation and Hypothesis Testing," Chapter 36 in R. Engle and D. McFadden (eds.), Handbook of Econometrics, Volume IV, New York: North-Holland.\\
Sargan, D., 1958, "The Estimation of Economic Relationships Using Instrumental Variables," Econometrica, 26, 393-415.\\
Searle, S., 1982, Matrix Algebra Useful for Statistics, New York: Wiley.\\
Staiger, D., and J. Stock, 1997, "Instrumental Variables Regression with Weak Instruments," Econometrica, 65, 557-586.\\
Theil, H., 1953, "Repeated Least-Squares Applied to Complete Equation Systems," MS., The Hague: Central Planning Bureau.\\
White, H., 1982, "Instrumental Variables Regression with Independent Observations," Econometrica, 50, 483-499.\\
\_\_, 1984, Asymptotic Theory for Econometricians, Academic.\\
Working, E. J., 1927, "What Do 'Statistical Demand' Curves Show?" Quarterly Journal of Economics, 41, 212-235.


\end{document}
